%!TEX root = ./../../springer-mono-format/book.tex
%!TEX root = ./../../thesis-trier.tex

%Further automata-theoretical notions are introduced in \cref{prelim:sec:automata},
%in particular various constructions and automata as language acceptors.
%%most of this section is devoted to introducing various constructions
%%as sub-, (Rees) factor, product, power, and pair automata. We also introduces automata as language acceptors 
%%and regular languages.

%\Cref{prelim:remarks_trans_semigrp_semiautomata,prelim:sec:remarks-reduction,prelim:sec:remarks_func_comp} 
%discusses conventions and notational choices.

Recall the definition of a (semi-)automaton and the set of synchronizing words
from \cref{intro:sec:basic,intro:sec:defsyncaut}.
Here, we give additional background on automata. %introduce additional notions. 
\cref{prelim:subsec:morenotions} 
collects more basic notions. The reachability quasiorder and related
notions are introduced in~\cref{prelim:subsec:reachorder}.
\Cref{prelim:subsec:transsemigrp} introduces
the transformation semigroup of an automaton
and a representation by transition matrices.
\Cref{prelim:subsubsec:endo_aut} investigates morphisms
on automata. We introduce endo-, iso- and automorphisms, 
and collect several properties, facts, and formulas; in particular
how morphisms act w.r.t.~the reachability order, strongly connected
components, and upper sets.
% The following sections are devoted
Then a sequence of sections follows
devoted to various constructions. 
Specifically, \SEC~\ref{prelim:sec:aut:sub}, \ref{prelim:sec:aut:factor}, \ref{prelim:sec:rees}, \ref{prelim:sec:productaut},
\ref{prelim:sec:aut:power}, and \ref{sec:pair_aut}
introduce the sub- and restricted, quotient, Rees quotient, product, power, and pair automaton, respectively.
\Cref{prelim:sec:aut:sequppersets} discusses
the set of maximal states and shows ways to decompose an automaton
into a sequence of upper sets. In~\cref{prelim:sec:autacceptors} 
we connect automata to formal languages.
We introduce regular languages
as languages accepted by automata when adjoined with a start state
and a set of final states (which we then call automata
as language acceptors).
\Cref{prelim:subsubsec:min_aut} investigates the minimal automaton
accepting a given regular language.

\section{More Basic Notions}
\label{prelim:subsec:morenotions}

%!TEX root = ./../../springer-mono-format/book.tex
%!TEX root = ./../../thesis-trier.tex

\label{def:inverse_transition}
Let $\semiautomat{A}{Q}{\Sigma}{\delta}$ be an automaton. For $u,v \in \Sigma^*$, $S \subseteq Q$ and $q \in Q$, we set 
$\applytrans{\delta^{-1}}{S}{u} %= f_u^{-1}(S) 
= \{ p \in Q \mid \applytrans{\delta}{p}{u} \in S \}$
and $\applytrans{\delta^{-1}}{q}{u} = \applytrans{\delta^{-1}}{\{q\}}{u}$. 
We have the following formula (related to \cref{prelim:eqn:invcomp}):
\begin{equation}\label{prelim:eqn:delta_inverse}
	\applytrans{\delta^{-1}}{S}{uv}
	= \applytrans{\delta^{-1}}{\applytrans{\delta^{-1}}{S}{v}}{u}.
\end{equation}

\IN{automaton graph}
The \emph{automaton (di-)graph} corresponding to $\mathcal A$ 
is the digraph $G(\mathcal A) = (Q, E)$ with $E = \{ (s,t) \mid \exists x \in \Sigma : \applytrans{\delta}{s}{x} = t \}$. \label{def:aut_graph}
%Via the automaton graph, graph theoretical concept transfer
%to automata.

\IN{self-loop}
\IN{trap state}
\index{sink state|see{trap state}}
We say $\mathcal A$ has a \emph{self-loop} at $q \in Q$ w.r.t.~$x \in \Sigma$ (or $x$ induces a self-loop at $q$)
if $\applytrans{\delta}{q}{x} = q$.
The state $q \in Q$ is called a \emph{trap} (or \emph{sink}) \emph{state} if
$\applytrans{\delta}{q}{x} = q$ for all $x \in \Sigma$, i.e., every letter induces a self-loop at $q$.\footnote{A sink state is sometimes defined to be a trap state reachable from every other state, but here it is just another name for a trap state.}

An automaton $\semiautomat{A}{Q}{\Sigma}{\delta}$ 
is called a \emph{unary} (resp.~\emph{binary}, \emph{ternary}) \emph{automaton} if $|\Sigma| = 1$ (resp.~$|\Sigma| = 2$, $|\Sigma| = 3$).
The automata from the \v{C}ern\'y
family (Figure~\ref{intro:fig:cerny}) are binary.
\IN{unary automaton}
\IN{binary automaton}
\IN{ternary automaton}

\IN{primitive automaton}
We call $\mathcal A$ \emph{primitive} if the trivial $\TS{\mathcal A}$-congruences
are the only $\TS{\mathcal A}$-congruences.

\IN{block!in automaton}
A subset $S \subseteq Q$ is said to be a \emph{block} in $\mathcal A$
if it is non-empty and $\delta(S, x) = S$
or $\delta(S, x) \cap S = \emptyset$
for all $x \in \Sigma$. This extends the notion of a block from permutation group theory
introduced in \cref{prelim:sec:permgrp} (see \pageref{def:block-perm-grp})
to automata. 
The context should make it clear whether we are referring to a block in a permutation group 
or in an automaton, so no confusion should arise.

\section{The Transformation Semigroup of an Automaton}
\label{prelim:subsec:transsemigrp}

%!TEX root = ./../../thesis-trier.tex

\label{def:rankword}
\label{def:defectword}
\label{def:exclword}
\label{def:duplword}
\label{def:transitionmtx} 
\label{def:fwordu}
\IN{Boolean transition matrix!automaton}
\IN{duplicate set!of a word}
\IN{permutational letter/word}
\IN{simple idempotent letter/word} 
\IN{idempotent letter/word}
\IN{rank!of a word}
\IN{defect!of a word}
\IN{excluded set!of a word}
\IN{function induced by word}

Let $\semiautomat{A}{Q}{\Sigma}{\delta}$. Each word $u \in \Sigma^*$  \emph{induces a transformation} 
$\fword{\delta}{u} \colon Q \to Q$ by $\applyfunc{q}{\fword{\delta}{u}} = \applytrans{\delta}{q}{u}$ for all $q \in Q$. 
For $u, v \in \Sigma^*$ we have $\fword{\delta}{uv} = \concat{\fword{\delta}{u}}{\fword{\delta}{v}}$
and $\fword{\delta}{{uv}}^{-1} = \concat{\fword{\delta}{v}^{-1}}{\fword{\delta}{u}^{-1}}$.
Compare the latter equation with Equation~\eqref{prelim:eqn:delta_inverse}
wheres $u$ and $v$ are ``swapped'' in the expression on the right hand side.
This follows as in the $\delta$-notation the left-most word 
is applied first whereas in the composition of two functions, by our definition,
the right-most function is applied first.

We call a letter $x \in \Sigma$ (resp.~a word $u \in \Sigma^*$)
a \emph{permutational / idempotent / simple idempotent} letter (resp.~word) 
if $\fword{\delta}{x}$ (resp.~$\fword{\delta}{u}$) is a permutation/idempotent/simple idempotent.
The \emph{rank} (resp.~\emph{defect}, \emph{excluded set}, \emph{duplicate set}) of a word $u \in \Sigma^*$, denoted by $\rank{u}$ (resp.~$\defect{u}$, $\excl{u}$, $\dupl{u}$) 
is the rank (resp.~defect, excluded set, duplicate set) of~$\fword{\delta}{u}$. 
Furthermore, when using terminology from functions, transformations, and permutations, as cyclically permuting the states, the image of a set, transition matrix,
being a cycle, having a fixed point etc., in the context of a word (or letter) $u \in \Sigma^*$,
we refer to the induced transformation~$\fword{\delta}{u}$.
We also say that~$\fword{\delta}{u}$ is the \emph{action of $u$} on the states
or that~$u$ \emph{acts as $\fword{\delta}{u}$}. The same terminology transfers when speaking about properties of $\fword{\delta}{u}$,
for example, we say that~$u$ acts as a cycle if~$\fword{\delta}{u}$ is a cycle.

For $u \in \Sigma^*$, we write $[u]_{\mathcal A}$ (or just $[u]$ when there is no ambiguity to which automaton we refer)
for the transition matrix of $\fword{\delta}{u}$.
\Cref{prelim:fig:transitionmtxCerny} depicts transition matrices for the \Cerny automata $\mathcal C_n$
from \cref{intro:sec:cerny} w.r.t.~the natural order of its state set~$[n]$.
When  the specific ordering of the states is not important, we
assume some order to be taken without explicitly mentioning it.\footnote{The maps $u \mapsto \fword{\delta}{u}$
	and $u \mapsto [u]$ are morphisms from $\Sigma^*$ into $\mathcal T_Q$ and the set of $|Q| \times |Q|$ matrices, respectively.
	Thus, by \cref{prelim:prop:UMP}, they
	are uniquely determined by their images on the letters~$\Sigma$.}

\begin{figure}[th]
	\centering
	\[
	[b]_{\mathcal C_n} =
	\left( {\begin{array}{ccccc}
			0 & 1 &  0 & \cdots & 0\\
			0 & 0 &  1 & \cdots & 0\\
			\vdots & \vdots & \vdots & \ddots & \vdots\\
			0 & 0 & 0 & \cdots & 1 \\
			1 & 0 & 0 & \cdots & 0
	\end{array} } \right)
	\qquad 
	[a]_{\mathcal C_n} =
	\left( {\begin{array}{cccccc}
			0 & 1 &  0 & \cdots & 0 & 0 \\
			0 & 1 &  0 & \cdots & 0 & 0 \\
			\vdots & \vdots & \vdots & \ddots & \vdots\\
			0 & 0 & 0 & \cdots & 1 & 0 \\
			0 & 0 & 0 & \cdots & 0 & 1
	\end{array} } \right)
	\]
	\caption[The transition matrices for the \Cerny automata $\mathcal C_n$]{
		The transition matrices for the \Cerny automata $\mathcal C_n$
		from \cref{intro:sec:cerny} (\cref{intro:fig:cerny})
		where the states are ordered w.r.t.~the natural order on $[n]$.
	}
	\label{prelim:fig:transitionmtxCerny}
\end{figure} 

\label{def:transformation_semigroup}
The \emph{transformation semigroup of $\mathcal A$} is 
\[ \TS{\mathcal A} = \{ \fword{\delta}{u} \mid u \in \Sigma^+ \}
\] 
and the \emph{transformation monoid} is $\TM{\mathcal A} =  \{ \fword{\delta}{u} \mid u \in \Sigma^* \}$. 
\label{def:transformation_monoid}
\index{transformation semigroup!of automaton}
\index{transformation monoid!of automaton}

\begin{proposition}
	\label{sync_max:prop:grp_of_units_inv}	
	Let $\semiautomat{A}{Q}{\Sigma}{\delta}$
	and $\fword{\delta}{u}$ be a permutation for $u \in \Sigma^*$.
	Then $\fword{\delta}{u}^{-1} \in \TM{\mathcal A}$,
	i.e., the inverse permutation
	is induced by some word over $\Sigma^*$.
\end{proposition}
\begin{proof}
	By finiteness, 
    we have $\fword{\delta}{u^i} = \operatorname{id}_Q$ for some $i > 0$.
	Then $\fword{\delta}{u}^{-1} = \fword{\delta}{u^{i-1}} \in \TM{\mathcal A}$.	
\end{proof}

Next, we determine the group of units in $\TM{\mathcal A}$.

\begin{proposition}
	\label{sync_max:prop:grp_of_units_TA}	
	Let $\semiautomat{A}{Q}{\Sigma}{\delta}$ and 
	$\Gamma = \{ x \in \Sigma \mid x \text{ is a permutational letter } \}$.
	Then for the group of units $G$ of $\TM{\mathcal A}$ we have
	\[
	G = \{ \delta_u \mid \text{$\delta_u$ is a permutation and $u \in \Sigma^*$} \}
	= \{ \fword{\delta}{u} \mid u \in \Gamma^* \},
	\]
	and if~$G$ is non-trivial, equivalently $\Gamma \ne \emptyset$,
	we also have $G =  \{ \fword{\delta}{u} \mid u \in \Gamma^+ \}$.
\end{proposition}
\begin{proof}
	We have	$G \subseteq \{ \delta_u \mid \text{$\delta_u$ is a permutation and $u \in \Sigma^*$} \}$ 
	as every transformation in $G$, having an inverse, is a permutation.
	Next, by~\crefitem{prelim:prop:rank-defect}{prelim:prop:item:rank-eqn} 
	the composition of a non-permutation with a transformation
	gives a non-permutation.  So, for $u \in \Sigma^*$,
	$\fword{\delta}{u}$ is a permutation if and only if $u \in \Gamma^*$, i.e.,
	\[ 
	\{ \delta_u \mid \text{$\delta_u$ is a permutation and $u \in \Sigma^*$} \}
	= \{ \fword{\delta}{u} \mid u \in \Gamma^* \}.
	\]	
	By \cref{sync_max:prop:grp_of_units_inv},
	for $u \in \Gamma^*$
	we have $\fword{\delta}{u}^{-1} = \fword{\delta}{v}$ for some $v \in \Sigma^*$, i.e.,
	$\fword{\delta}{u}$ has an inverse in $\TM{\mathcal A}$
	and so $\fword{\delta}{u} \in G$.
	This gives $\{ \fword{\delta}{u} \mid u \in \Gamma^* \} \subseteq G$.
	Combining everything,
	we find $G = \{ \delta_u \mid \text{$\delta_u$ is a permutation and $u \in \Sigma^*$} \}
	= \{ \fword{\delta}{u} \mid u \in \Gamma^* \}$.
	
	Lastly, if $\Gamma \ne \emptyset$,
	then $\fword{\delta}{u}$ equals the identity transformation
	for some $u \in \Gamma^+$.
	This gives $\{ \fword{\delta}{u} \mid u \in \Gamma^* \}
	= \{ \fword{\delta}{u} \mid u \in \Gamma^+ \}$.
\end{proof}

Thus the group of units of $\TM{\mathcal A}$ is precisely the set of permutations in~$\TM{\mathcal A}$
and every permutation group in $\TS{\mathcal A}$ is contained in the group of units of $\TS{\mathcal A}$ (but not every
subset $S \subseteq \TM{\mathcal A}$ that forms a group in itself; for example
for every non-permutational word $u$ permuting its image $\delta(Q, u)$ 
the set $\{ \fword{\delta}{u^i} \mid i \ge 1 \}$
forms a group, but contains no permutation).

\section{The Reachability Order}
\label{prelim:sec:reach} 
\label{prelim:subsec:reachorder}

First, in \cref{prelim:sec:reach:def} we introduce the reachability order.
Then, \cref{prelim:sec:reach:conn} investigates connectivity
notions and the notion of an upper set derived
from the reachability order.

\subsection{Definition and a Characterization}
\label{prelim:sec:reach:def}

%!TEX root = ./../../thesis-trier.tex

\IN{reachability order!automaton}
\index{reachability (quasi-)order|see{reachability order}}
Let $\semiautomat{A}{Q}{\Sigma}{\delta}$ and $s,t \in Q$. We define the \emph{reachability (quasi)order} 
\[
s \lereach^{\mathcal A} t 
\Leftrightarrow 
\exists u \in \Sigma^* \colon \applytrans{\delta}{s}{u} = t. \label{def:reach_order}
\]
We say that $t$ is \emph{reachable} from $s$ (in $\mathcal A$) if $s \lereach^{\mathcal A} t$. 
When the automaton we refer to is clear from context, we simply write~$\lereach$.
The reachability order is characterized as the finest quasiorder such that $s \lereach \applytrans{\delta}{s}{u}$
for all $s \in Q$ and~$u \in \Sigma^*$, as we formally state next.

\begin{proposition}
	\label{prelim:prop:reachability-finest}
	Let $\semiautomat{A}{Q}{\Sigma}{\delta}$. We have $s \lereach \applytrans{\delta}{s}{u}$ for all $s \in Q$ and $u \in \Sigma^*$.
	Further, if $\lesssim$ is a quasiorder on $Q$
	such that $s \lesssim \applytrans{\delta}{s}{u}$
	for all $s \in Q$ and $u \in \Sigma^*$,
	then $s \lereach t$ implies $s \lesssim t$ for all $s,t \in Q$.
\end{proposition} 
\begin{proof}
	Clearly, we have $s \lereach \applytrans{\delta}{s}{u}$ for $s \in S$ and $u \in \Sigma^*$.
	Next, suppose $\lesssim$
	is a quasiorder such that $s \lesssim \delta(s, u)$ for all $s \in Q$ and $u \in \Sigma^*$. 
	Now, for $s,t \in Q$, $s \lereach t$ is equivalent to $\applytrans{\delta}{s}{u} = t$
	for some $u \in \Sigma^*$. So $s \lesssim \applytrans{\delta}{s}{u} = t$.  
\end{proof}

\label{def:induced_reach_order} 
\IN{reachability equivalence}
By $\eqreach$ we denote the \emph{(reachability) equivalence relation} induced by $\lereach$ 
according to~\cref{prelim:eqn:induced_eqrel}.

%\subsection{Strongly Connected Components and Upper Sets}
\subsection{Connectivity and Upper Sets}
\label{prelim:sec:reach:conn} 

%!TEX root = ./../../thesis-trier.tex

\IN{unilaterally connected!automaton} 
\IN{strongly connected!in automaton}
\IN{strongly connected component!automaton}
\IN{upper set!in automaton}
Let Let $\semiautomat{A}{Q}{\Sigma}{\delta}$ and $S \subseteq Q$. If nothing else is said, 
order-theoretical notions in the context of automata
are meant w.r.t.~the reachability order. For example,
by calling~$S$ an \emph{upper set}, we mean 
it to be an upper set w.r.t.~the reachability quasiorder. This is equivalent
to $\applytrans{\delta}{S}{x} \subseteq S$ for all $x \in \Sigma$. Note that the empty set $\emptyset$
is an upper set in each automaton.
Since then $\applytrans{\delta}{S}{u} \subseteq \applytrans{\delta}{Q}{u} \cap S$ for all
$u \in \Sigma^*$, we obtain
the following equation:
\begin{equation} 
	\text{$S$ is an upper set} \Rightarrow \forall u \in \Sigma^* \colon \applytrans{\delta}{Q}{u} \cap S \ne \emptyset. \label{sync_max:eqn:upper_set} 
\end{equation}
The automaton $\semiautomat{A}{\{q_1,q_2\}}{\{a\}}{\delta}$
with $\applytrans{\delta}{q_1}{a} = q_2$
and $\applytrans{\delta}{q_2}{a} = q_1$ 
shows, by considering $S = \{q_1\}$,
that the condition on the right hand side can hold 
for sets that are not upper sets. A state $t \in Q$
is a trap state if and only if $\{t\}$ is an upper set.
We also have
%For $\semiautomat{A}{Q}{\Sigma}{\delta}$ and for $S \subseteq Q$ we have
\begin{equation}\label{eqn:closed-block}
	\text{$S$ is upper set and block} \Leftrightarrow  \forall x \in \Sigma \colon \delta(S, x) = S \Leftrightarrow \forall u \in \Sigma^* \colon \delta(S, u) = S.
\end{equation}

\IN{sink component}
A subset $S \subseteq Q$ is called \emph{strongly connected} (resp.~\emph{unilaterally connected}) 
if, for all $s, t \in S$, we have $s \lereach t$
%and $t \lereach s$ 
(resp.~$s \lereach t$
or $t \lereach s$). 
A maximal strongly connected subset is called a \emph{strongly connected component}
and equals an equivalence class w.r.t.~$\eqreach$.
A strongly connected component that is an upper set is called a \emph{sink (strongly connected) component}. Equivalently,
it is an upper set that is strongly connected.
By \cref{prelim:general:order:prop:min-equiv-class}, 
a sink component is precisely a minimal non-empty upper set w.r.t.~set inclusion.
In the automaton graph, the sink components are the strongly connected components
without transitions leading out of it.

\begin{proposition} \label{prelim:autsyncw:lem:minrank}
	In an automaton, the minimal possible rank of a word
	is bounded from below by the number of distinct sink components.
\end{proposition}
\begin{proof}
	This follows from \cref{sync_max:eqn:upper_set} 	
	and the fact that distinct sink components are disjoint.
\end{proof}

\begin{proposition} \label{prelim:autsyncw:lem:char-union-of-sink}
	Let $\semiautomat{A}{Q}{\Sigma}{\delta}$ and
	$S \subseteq Q$ be an upper set.
	Then~$S$ is a union of sink components
	if and only if $\delta(q, u) \eqreach q$ for all $q \in Q$
	and $u \in \Sigma^*$.
\end{proposition}
\begin{proof}
	First, suppose $\delta(q, u) \eqreach q$ for all $q \in Q$
	and $u \in \Sigma^*$. Consider some $q \in S$.
	By assumption $\upperset{\{q\}} \subseteq [q]_{\eqreach}$,
	and so $[q]_{\eqreach}$ is a sink component
	by \cref{prelim:general:order:cor:min-equiv-class}. 
	Since~$S$ is an upper set, we get $[q]_{\eqreach} \subseteq S$
	and so~$S$ equals the union of all $[q]_{\eqreach}$
	with $q \in S$.
	
	Conversely, suppose~$S$ is a union of sink components.
	Let $q\in S$. Then~$q$ is contained
	in some sink component, which yields that $[q]_{\eqreach}$
	is a sink component and so $\delta(q, u) \eqreach q$.	
\end{proof}

\begin{corollary} \label{prelim:autsyncw:lem:perm-each-letter-sink}
	A set permuted by each letter equals a union
	of sink components.	
\end{corollary}
\begin{proof}
	Let $S \subseteq Q$ be a set in $\semiautomat{A}{Q}{\Sigma}{\delta}$
	such that $\delta(S, a) = S$ for each $a \in \Sigma$.
	Then $S$ is an upper set.	
	Consider $q \in S$ and $u \in \Sigma^*$.	
	Since the word~$u$ permutes~$S$,
	there exists $i > 0$
	such that $\delta(p, u^i) = p$
	for all $p \in Q$.
	Then $\delta(q, u u^{i- 1}) = q$
	gives $\delta(q, u) \lereach q$
	and since $q \lereach \delta(q, u)$
	the claim follows from~\cref{prelim:autsyncw:lem:char-union-of-sink}.
\end{proof}

\IN{strongly connected automaton}
\IN{unilaterally connected automaton}
\IN{cyclic automaton}
We call~$\semiautomat{A}{Q}{\Sigma}{\delta}$ \emph{strongly connected} (resp.~\emph{unilaterally connected})
if~$Q$ is strongly connected (resp.~unilaterally connected).
For $Q \ne \emptyset$,~$\mathcal A$ is strongly connected if and only if $Q$
is a strongly connected component. 
We call~$\mathcal A$ \emph{cyclic}~\cite{Ito04} if $Q = \upperset{\{q\}}$
for some $q \in Q$.

\begin{proposition}\label{prelim:automata:prop:char-strongly-conn}
	An automaton~$\semiautomat{A}{Q}{\Sigma}{\delta}$ 
	is strongly connected 
	if and only if there exists no proper non-empty upper set, i.e., $\emptyset$
	and~$Q$ are the only upper sets.
\end{proposition}
\begin{proof}
	This follows from \cref{prelim:automata:prop:TsetminusS-strong-conn}.	
\end{proof}

\begin{proposition}\label{cor:maximal_upper_set_trap}
	Let $\semiautomat{A}{Q}{\Sigma}{\delta}$ with a trap state $t \in Q$ and $|Q| \ge 2$. Then $Q \setminus \{t\}$
	is a strongly connected component
	if and only if $Q$ and $\{t\}$
	are the only non-empty upper sets in $\mathcal A$.		
\end{proposition} 
\begin{proof}
	This follows from \cref{prelim:automata:prop:TsetminusS-strong-conn} with $S = \{t\}$.
\end{proof}

\begin{proposition}\label{prelim:automata:unilateral-implies-cyclic}
	An unilaterally connected
	automaton $\semiautomat{A}{Q}{\Sigma}{\delta}$
	is cyclic.
\end{proposition}
\begin{proof}
	Let $q \in Q$ be such that $\upperset{\{q\}}$ has maximal size.
	Assume $p \notin \upperset{\{q\}}$ for some $p \in Q$.
	Since~$\mathcal A$ is unilaterally connected, 
	we have $q \in \upperset{\{p\}}$
	and so $\upperset{\{q\}} \subsetneq \upperset{\{p\}}$,
	contradicting the maximal choice of~$q$.	
\end{proof}

The converse of \cref{prelim:automata:unilateral-implies-cyclic} does not hold. For example,
for $\semiautomat{A}{\{q_1, q_2, q_3\}}{\{a,b\}}{\delta}$
with $\applytrans{\delta}{q_1}{a} = q_2$, $\applytrans{\delta}{q_1}{b} = q_3$,
and $\applytrans{\delta}{q_i}{x} = q_i$ for $i \in \{2,3\}$
and $x \in \{a,b\}$
we have $\upperset{\{q_1\}} = \{q_1,q_2,q_3\}$, but
it is not unilaterally connected.

By the results from \cref{prelim:sec:order},
a subset of states is a strongly connected component
if and only if for each upper set it is contained in this set
or its complement (\cref{prelim:general:order:subset-in-equiv}), 
every upper set is a union of strongly connected components (\cref{prelim:general:order:upper-union-equiv-class}),
and for the set difference $T \setminus S$ of two upper sets~$S,T$
the following are equivalent: strong connectivity,
being a strongly connected component, and maximality relative to~$T$ 
and set inclusion (\cref{prelim:automata:cor:SsetminusT,prelim:automata:prop:TsetminusS-strong-conn}).

\IN{strongly trap-connected automaton}
An automaton $\semiautomat{A}{Q}{\Sigma}{\delta}$
is said to be \emph{strongly trap-connected} if there exists a trap state~$t \in Q$
such that for all $p \in Q \setminus \{t\}$ and $q \in Q$
we have $p \lereach q$. Clearly, a strongly trap-connected
automaton has a unique trap-state, is synchronizing,
and unilaterally connected.

\begin{proposition}\label{prelim:automata:prop:char-strongly-trap}
	Let $\semiautomat{A}{Q}{\Sigma}{\delta}$ with $|Q| \ge 2$.
	Then the following are equivalent.
	
	\begin{enumerate}
		\item \label{prelim:automata:prop:char-strongly-trap:1} $\mathcal A$ is strongly trap-connected.
		
		\item \label{prelim:automata:prop:char-strongly-trap:1b} $Q = S \cup \{t\}$ for a trap state $t \in Q$ and a non-sink strongly connected component~$S$.
		
		\item \label{prelim:automata:prop:char-strongly-trap:2} There exists $t \in Q$ such that $\{t\}$ is the unique proper and non-empty upper set.
		
		\item \label{prelim:automata:prop:char-strongly-trap:3} $\mathcal A$ is not strongly connected and has no proper upper set of size at least two.
	\end{enumerate}
\end{proposition}
\begin{proof}
	$\ref{prelim:automata:prop:char-strongly-trap:1} \Rightarrow~\ref{prelim:automata:prop:char-strongly-trap:1b}$:
	Let~$t$ be the unique trap state.
	Suppose~$R$ is a non-empty upper set distinct from $\{t\}$.
	Then there exists $q \in R \setminus\{t\}$.
	Since~$\mathcal A$ is strongly trap-connected,
	we have $\upperset{\{q\}} = Q$. As an upper set, 
	it is $\upperset{\{q\}}\subseteq R$
	and so $R = Q$. Thus $Q$ and $\{t\}$ are the only non-empty upper sets.
	Since~$\{t\}$ is a sink component, \cref{prelim:general:order:XS-only-uppersets} 
	gives the claim.	
	
	$\ref{prelim:automata:prop:char-strongly-trap:1b} \Rightarrow~\ref{prelim:automata:prop:char-strongly-trap:2}$:
	This follows from \cref{prelim:general:order:XS-only-uppersets}.
	
	$\ref{prelim:automata:prop:char-strongly-trap:2} \Rightarrow~\ref{prelim:automata:prop:char-strongly-trap:3}$: Since~$\mathcal A$ has at least two states
	and an upper set~$\{t\}$ of size one, by \cref{prelim:automata:prop:char-strongly-conn}
	it is not strongly connected. By assumption~$\{t\}$
	is a unique proper and non-empty upper set, so that~$\mathcal A$
	has no proper upper set of size at least two.
	
	$\ref{prelim:automata:prop:char-strongly-trap:3} \Rightarrow~\ref{prelim:automata:prop:char-strongly-trap:1}$: Since~$\mathcal A$
	is not strongly connected, by \cref{prelim:automata:prop:char-strongly-conn}
	there exist some proper and non-empty upper set~$S \subseteq Q$.
	By assumption, $S = \{t\}$ for some $t \in Q$.
	Let $q \in Q \setminus\{t\}$.
	Then $\upperset{\{q\}} \cup \{t\}$ is an upper set
	of size at least two. Hence, we have $\upperset{\{q\}} \cup \{t\} = Q$.
	This gives $|\upperset{\{q\}}| \ge 2$. Now, we have $t \in \upperset{\{q\}}$
	as otherwise, $\upperset{\{q\}}$
	is a proper upper set of size at least two,
	and so $\upperset{\{q\}} = Q$.
\end{proof}

\begin{corollary}
	\label{prelim:fact:no-proper-upper-sets-size-two}
	In automata with at least three states, there exists no proper upper set of size at least two in $\mathcal A$
	if and only if~$\mathcal A$ is strongly connected
	or strongly trap-connected.	  
\end{corollary}
\begin{proof}
	By \cref{prelim:automata:prop:char-strongly-conn,prelim:automata:prop:char-strongly-trap},
	a strongly connected or strongly trap-connected
	automaton has no proper upper set of size at least two.
	Conversely, the automaton is either strongly connected
	or, by \crefitem{prelim:automata:prop:char-strongly-trap}{prelim:automata:prop:char-strongly-trap:3}
	it has no proper upper set of size at least two.
\end{proof}

A two-state automaton where every letter induces a self-loop
also has no proper upper set of size at least two, but is neither strongly connected
nor strongly trap-connected.

\section{Morphisms} 
\label{prelim:subsubsec:endo_aut}

%!TEX root = ./../../springer-mono-format/book.tex
%!TEX root = ./../../thesis-trier.tex

Here, we introduce morphisms, endomorphisms, isomorphisms, and automorphisms
of automata and collect basic facts about them.

\subsection{Definitions and Basic Facts}

Let $\semiautomat{A}{Q}{\Sigma}{\delta}$ and $\semiautomat{B}{S}{\Sigma}{\mu}$
be two semiautomata. 

\IN{morphism!of (semi-)automata}
\IN{endomorphism!semiautomaton}
\IN{isomorphism!semiautomaton}
\IN{automorphism!semiautomaton}
\IN{isomorphic automata}
\begin{enumerate}
	\item A \emph{(semiautomaton) (homo-)morphism} from $\mathcal A$ to $\mathcal B$
	(or between $\mathcal A$ and $\mathcal B$) is a function
	$\varphi \colon Q \to S$
	such that for each $q \in Q$ and $x \in \Sigma$ we have
	\[
	\applyfunc{\applytrans{\delta}{q}{x}}{\varphi}
	=
	\applytrans{\mu}{\applyfunc{q}{\varphi}}{x}.
	\]
	Equivalently, $\varphi$ is a morphism if and only if
	$\concat{\delta_x}{\varphi} = \concat{\varphi}{\mu_x}$
	for all $x \in \Sigma$.
	Thus, regardless of which of the two paths we follow in the diagram below,
	we end up at the same state.
	\[
	\begin{tikzcd}[row sep=1cm,column sep=1.25cm]
		Q \ar{r}{\delta(\cdot, x)} \ar[swap]{d}{\varphi}
		& Q \ar{d}{\varphi} \\
		S \ar[swap]{r}{\mu(\cdot, x)}
		& S
	\end{tikzcd}
	\]
	If $\varphi$ is surjective, we say that it maps $\mathcal A$ \emph{onto}
	$\mathcal B$.
	
	\item  A \emph{(semiautomaton) endomorphism}
	is a morphism with $\mathcal A = \mathcal B$.
	
	\item A \emph{(semiautomaton) isomorphism} is a morphism $\varphi \colon Q \to S$ such that
	there exists a morphism $\psi \colon S \to Q$ with $\concat{\varphi}{\psi} = \operatorname{id}_Q$
	and $\concat{\psi}{\varphi} = \operatorname{id}_S$.
	The automata~$\mathcal A$ 
	and~$\mathcal B$ are called \emph{isomorphic}, written \[ \mathcal A \cong \mathcal B, \] \label{def:isomorphic}
	if there exists an isomorphism between them.
	
	\item A \emph{(semiautomaton) automorphism} is an endomorphism that is an isomorphism.
\end{enumerate}

The set of endomorphisms of a given automaton forms a transformation semigroup
and the set of automorphisms forms a permutation group.

\begin{example} \label{prelim:aut:ex:mappings}
	\label{prelim:ex:SCC} \label{prelim:ex:SCC2}
	Consider the automata from \cref{prelim:aut:fig:mappings}, where the dashed lines indicate a 
	mapping of their state sets.
		
	In the automaton from \cref{prelim:aut:fig:mappingmorph}
	the mapping is an endomorphism that does not map the strongly connected
	component $\{q_0\}$ to a strongly connected component, but to $\{q_1\}$.	
	Another endomorphism of the same automaton is given in \cref{prelim:aut:fig:mappingmorph0}.
	Together with the identity transformation, this automaton has exactly three endomorphisms.
	
	In \cref{prelim:aut:fig:mappingnotmorph0,prelim:aut:fig:mappingotmorph},
	the mapping $\varphi$ is not a morphism since in both cases we have
	$\varphi(\delta(q_0, a)) = q_1 \ne q_2 = \delta(\varphi(q_0), a)$.
	
	\begin{figure}[ht]
	\centering
	\begin{subfigure}{0.25\textwidth} \centering
		\begin{tikzpicture}[node distance = 2cm, on grid, auto, ->]
			\node (q0) [state] {$q_0$};
			\node (q1) [state, below of = q0] {$q_1$};
			\node (q2) [state, below of = q1] {$q_2$};
			
			\path (q0) edge node {$a$} (q1) 
			(q1) edge [bend left] node {$a$} (q2)
			(q2) edge [bend left] node {$a$} (q1);
			
			\path (q0) edge [dashed,bend left=2cm] (q1);
			\path (q1) edge [dashed, bend left=3cm] (q2)
			(q2) edge [dashed, bend left=3cm] (q1);
		\end{tikzpicture}
		\caption{An endomorphism.}
		\label{prelim:aut:fig:mappingmorph}
	\end{subfigure}%
	\hfill
	\begin{subfigure}{0.25\textwidth} \centering
		\begin{tikzpicture}[node distance = 2cm, on grid, auto, ->]
			\node (q0) [state] {$q_0$};
			\node (q1) [state, below of = q0] {$q_1$};
			\node (q2) [state, below of = q1] {$q_2$};
			
			\path (q0) edge node {$a$} (q1) 
			(q1) edge [bend left] node {$a$} (q2)
			(q2) edge [bend left] node {$a$} (q1);
			
			\path (q0) edge [dashed,bend left=2cm] (q2);
			\path (q1) edge [dashed, loop left] (q1)
			(q2) edge [dashed, loop left] (q2);
		\end{tikzpicture}
		\caption{An endomorphism.}
		\label{prelim:aut:fig:mappingmorph0}
	\end{subfigure}%
	\hfill
	\begin{subfigure}{0.25\textwidth} \centering
		\begin{tikzpicture}[node distance = 2cm, on grid, auto, ->]
			\node (q0) [state] {$q_0$};
			\node (q1) [state, below of = q0] {$q_1$};
			\node (q2) [state, below of = q1] {$q_2$};
			
			\path (q0) edge node {$a$} (q1) 
			(q1) edge [bend left] node {$a$} (q2)
			(q2) edge [bend left] node {$a$} (q1);
			
			\path (q0) edge [dashed,bend left=2cm] (q1);
			\path (q1) edge [dashed, loop left] (q1)
			(q2) edge [dashed, loop left] (q2);
		\end{tikzpicture}
		\caption{Not a morphism.}
		\label{prelim:aut:fig:mappingnotmorph0}
	\end{subfigure}%
	\hfill
	\begin{subfigure}{0.25\textwidth} \centering
		\begin{tikzpicture}[node distance = 2cm, on grid, auto, ->]
			\node (q0) [state] {$q_0$};
			\node (q1) [state, below of = q0] {$q_1$};
			\node (q2) [state, below of = q1] {$q_2$};
			
			\path (q0) edge [loop right ]node {$a$} (q0) 
			(q1) edge [bend left] node {$a$} (q2)
			(q2) edge [bend left] node {$a$} (q1);
			
			\path (q0) edge [dashed,bend left=2cm] (q1);
			\path (q1) edge [dashed, bend left=3cm] (q2)
			(q2) edge [dashed, bend left=3cm] (q1);
		\end{tikzpicture}
		\caption{Not a morphism.}
		\label{prelim:aut:fig:mappingotmorph}
	\end{subfigure}
	\caption[Example morphisms and non-morphisms on automata]{Automata and mappings of their states indicated by dashed lines, see \cref{prelim:aut:ex:mappings}.}
	\label{prelim:aut:fig:mappings}
\end{figure}
\end{example}

\begin{example} \label{prelim:aut:ex:autnomorph0}
	Let $\semiautomat{A}{\{q_1, q_2\}}{\{a_1,a_2\}}{\delta}$
	with $\delta(q_1, a_i) = \delta(q_2, a_i) = q_i$
	for $i \in \{1,2\}$, see \cref{prelim:aut:fig:autnomorph}.
	If~$\varphi \colon Q \to Q$ is an endomorphism,
	then $\varphi(q_i) = \varphi(\delta(q_i, a_i)) = \delta(\varphi(q_i), a_i) = q_i$
	for $i \in \{1,2\}$.
	Thus the identity transformation is the only endomorphism of~$\mathcal A$.	
	
	Later, in \cref{prelim:automata:morphism:cor:endo-sync-SCC}, we will show that, in general, a 
	strongly connected and synchronizing automaton only admits the identity transformation as an endomorphism.
	
	\begin{figure}[ht]
		\centering	
		\begin{tikzpicture}[node distance = 2cm, on grid, auto, ->]
			\node (q0) [state] {$q_1$};
			\node (q1) [state, right of = q0] {$q_2$};
			
			\path (q0) edge [bend left=10] node {$a_2$} (q1) 
			(q1) edge [loop right] node {$a_2$} (q1);
			
			\path (q1) edge [bend left=10] node {$a_1$} (q0)
			(q0) edge [loop left] node {$a_1$} (q0);
		\end{tikzpicture}
		\caption[Automaton not admitting any non-identity endomorphism.]{Automaton not admitting any non-identity endomorphism, see \cref{prelim:aut:ex:autnomorph0}.}
		\label{prelim:aut:fig:autnomorph}
	\end{figure}
\end{example}

Next, we give some alternative characterizations for morphisms.

\begin{proposition} \label{prelim:automata:morphism:prop:principal-upper-set}
	Let $\varphi \colon Q \to S$ be a function. Then the following are equivalent.
	
	\begin{enumerate}
		\item \label{prelim:automata:morphism:prop:principal-upper-set:1} $\varphi$ is a morphism.
		
		\item \label{prelim:automata:morphism:prop:principal-upper-set:2} For all $p \in Q$ and all $u \in \Sigma^*$
		we have
		\[
		\varphi(\delta(p, u)) = \mu(\varphi(p), u)
		\]
		
		\item \label{prelim:automata:morphism:prop:principal-upper-set:3} There exists $T \subseteq Q$ with $Q = \bigcup_{p \in T} \upperset{\{p\}}$
		such that for all $p \in T$ and all $u \in \Sigma^*$ we have
		\[
		\varphi(\delta(p, u)) = \mu(\varphi(p), u)
		\]
	\end{enumerate}
\end{proposition}
\begin{proof}
	$\ref{prelim:automata:morphism:prop:principal-upper-set:1}\Rightarrow\ref{prelim:automata:morphism:prop:principal-upper-set:2}$: This clearly holds for $u = \varepsilon$.
	Next, let $u = va$ for some $v \in \Sigma^*$ and $a \in \Sigma$.
	Inductively, we have $\varphi(\delta(p, v)) = \mu(\varphi(p), v)$.
	Using that~$\varphi$ is a morphism, this gives $\varphi(\delta(p, u)) 
	= \varphi(\delta(\delta(p, v), a))
	= \mu(\varphi(\delta(p, v)), a)
	= \mu(\mu(\varphi(p), v), a)
	= \mu(\varphi(p), u)$.
	Then the claim follows by induction.
	
	$\ref{prelim:automata:morphism:prop:principal-upper-set:2}\Rightarrow\ref{prelim:automata:morphism:prop:principal-upper-set:3}$: By assumption, this holds for $T = Q$.
	
	$\ref{prelim:automata:morphism:prop:principal-upper-set:3}\Rightarrow\ref{prelim:automata:morphism:prop:principal-upper-set:1}$: Let $a \in \Sigma$ and $q \in Q$.
	By assumption, there exists $u \in \Sigma^*$ such that $q = \delta(p, u)$
	for some $p \in T$. Set $s = \varphi(p)$.
	Then we have
	\begin{align*}
		\varphi(\delta(q, a))
		& = \varphi(\delta(p, ua)) \tag{$q = \delta(p,u)$} \\
		& = \mu(s, ua)\tag{$\varphi$ morphism}  \\
		& = \mu(\mu(s, u), a) \tag{Def.~transition function} \\
		& = \mu(\varphi(\delta(p, u)), a) \tag{$\varphi(\delta(p, u)) = \mu(s, u)$} \\
		& = \mu(\varphi(q), a). \tag{$q = \delta(p,u)$}
	\end{align*}
	Thus, $\varphi$ is a morphism.
\end{proof}

By \cref{prelim:automata:morphism:prop:principal-upper-set:3} of \cref{prelim:automata:morphism:prop:principal-upper-set}, a morphism
is completely specified by the images of $\varphi(p_1), \ldots, \varphi(p_m)$
for some $p_1, \ldots, p_m \in Q$ such that 
$Q = \upperset{\{p_1\}} \cup \cdots \cup \upperset{\{p_m\}}$.

\begin{corollary}
	If~$\mathcal A$ is cyclic, there are at most~$|S|$ morphisms
	from~$\mathcal A$ into~$\mathcal B$.
	In particular, there are at most~$|Q|$ endomorphisms of~$\mathcal A$.
\end{corollary}
\begin{proof}
	By \cref{prelim:automata:morphism:cor:two-morp-upper-set} 
	two morphisms having the same image for a state~$p$ with $Q = \upperset{\{p\}}$
	are equal. Since we have at most $|S|$ distinct images for this state,
	the claim follows.
\end{proof}

\begin{example}
	The might be strictly less than $|Q|$ endomorphisms in a cyclic automaton.
	For example, the automaton from \cref{prelim:aut:ex:autnomorph0}
	is cyclic and admits no non-identity morphism.
\end{example}

\IN{morphic image}
In the rest of this section, let  $\varphi \colon Q \to S$ 
be a morphism. For $T \subseteq Q$, we call $\applyfunc{T}{\varphi}$ the \emph{morphic image (of $T$)}.

By our next statement, isomorphisms are precisely the bijective morphisms. 
Note that a similar statement does not hold for order morphisms, see \cref{prelim:ex:order:bijmorph}.

\begin{proposition}\label{prelim:prop:iso-bijmorph}
	The morphism $\varphi \colon Q \to S$ is an isomorphism if and only if it is bijective.
\end{proposition}
\begin{proof}
	By \cref{prelim:prop:surj-inj-inverses}, existence of a map $\psi \colon S \to Q$ 
	such that $\concat{\varphi}{\psi} = \operatorname{id}_Q$
	and $\concat{\psi}{\varphi} = \operatorname{id}_S$
	is equivalent to bijectivity of $\varphi$. So, an isomorphism is bijective.
	
	Next, assume~$\varphi$ is bijective. 
	Since~$\varphi$ is surjective,
	we have $\concat{\psi}{\varphi} = \operatorname{id}_S$
	for some $\psi \colon S \to Q$ 
	by \cref{prelim:prop:surj-inj-inverses}.
	Let $s \in S$ and $x \in \Sigma$.
	Then we find
	\begin{align*}
		\varphi(\delta(\psi(s), x))
		& = \mu(\varphi(\psi(s)), x) \tag{$\varphi$ morphism} \\
		& = \mu(s, x) \tag{$\concat{\psi}{\varphi} = \operatorname{id}_S$} \\
		& = \varphi(\psi(\mu(s, x))). \tag{$\concat{\psi}{\varphi} = \operatorname{id}_S$}
	\end{align*}
	Since~$\varphi$ is injective, $\varphi(\delta(\psi(s), x)) = 
	\varphi(\psi(\mu(s, x)))$
	implies $\delta(\psi(s), x) = \psi(\mu(s, x))$.
	Thus $\psi$ is a morphism.
\end{proof}

\begin{remark}
	By \cref{prelim:prop:surj-inj-inverses} 
	a morphism $\varphi \colon Q \to S$
	is injective (resp.~surjective) if
	there exists a morphim $\psi \colon S \to Q$
	such that $\concat{\varphi}{\psi}$ (resp.~$\concat{\psi}{\varphi}$).
	However, unlike for functions,
	the converse implication does not hold in general.	

	To show this, for $i \in \{1,2,3\}$, let $\mathcal A_i = (Q_i, \{a\}, \delta_i)$
	where $Q_1 = \{ q_1, q_2 \}$, $Q_2 = \{ q_1, q_2, t \}$, $Q_3 = \{t\}$,
	and where the letter~$a$ swaps the states $q_1$ and~$q_2$ in $\mathcal A_1$ and~$\mathcal A_2$
	and fixes the state~$t$ in $\mathcal A_2$ and~$\mathcal A_3$.	
	Define $\varphi \colon Q_1 \to Q_2$ by $\varphi(q_1)=q_1$, $\varphi(q_2)=q_2$,
	and $\psi \colon Q_1 \to Q_3$ by $\psi(q_1)=\psi(q_2)=t$.
	Then $\varphi$ is injective and $\psi$ is surjective
	and both maps are morphisms.	
	However, there exists no morphism from $\mathcal A_2$ to $\mathcal A_1$,
	and no morphism from $\mathcal A_3$ to $\mathcal A_1$.
\end{remark}

\begin{proposition} \label{prelim:aut:prop:morph:trap-to-trap}
	If~$t \in Q$ is a trap state, then $\varphi(t)$ is a trap state.
\end{proposition}
\begin{proof}
	Let $x \in \Sigma$.
	Then $\mu(\varphi(t), x) = \varphi(\delta(t,x)) = \varphi(t)$
	and so~$\varphi(t)$ is a trap state.
\end{proof}

\begin{proposition} \label{prelim:automata:morphism:prop:surj-morphism-preserves-sync}
	If~$\varphi$ is surjective, the image of a synchronizing state is a synchronizing state.
\end{proposition} 
\begin{proof}
	Let $\delta(Q, u) = \{q\}$ for $u \in \Sigma^*$ and $q \in Q$.
	Then $\mu(S, u) = \mu(\varphi(Q), u) = \varphi(\delta(Q,u)) = \{\varphi(q)\}$
\end{proof}

Next, a morphism of automata induces a morphism of the underlying digraphs.

\begin{proposition} \label{prelim:aut:prop:morph-digraph}
	$\varphi \colon Q \to S$ is also a graph morphism
	of the automaton digraphs of~$\mathcal A$ and~$\mathcal B$.
\end{proposition}
\begin{proof}
	For $p \in Q$ and $x \in \Sigma$, the edge $(p, \delta(p,x))$
	is mapped to the edge $(\varphi(p), \mu(\varphi(p), x))$
	by~$\varphi$ in the respective automaton digraphs,
	which gives the claim.
\end{proof}

The following two examples gives mappings that are graph morphisms of the automaton digraphs,
but not automaton morphisms.

\begin{example}
Let $\semiautomat{A}{\{p,q\}}{\{a,b\}}{\delta}$
with $\delta(p, x) = q$ for $x \in \{a,b\}$, $\delta(q, a) = p$, and $\delta(q, b) = q$.
Then the transformation~$\varphi$ that maps all states to~$q$
is a graph morphism of the automaton digraph of~$\mathcal A$,
but it is not an endomorphism
since $\delta(\varphi(q), a) = p$ and $\varphi(\delta(q, a)) = q$.
\end{example}

\begin{example}
Let $\semiautomat{A}{\{q_1,\ldots, q_n\}}{\Sigma}{\delta}$ with $n \ge 2$
be such that for each pair of distinct states $q_i, q_j$ there exists 
a letter $a_{ij}$ inducing a simple idempotent mapping $q_i$ to~$q_j$.
Every permutation of $\{q_1, \ldots, q_n\}$
is a graph automorphism of the automaton digraph of~$\mathcal A$.
However, $\mathcal A$ has no non-identity
automaton morphism
since if $\varphi(q_i) = q_j$ with $i \ne j$,
then $\varphi(\delta(q_i, a_{ji})) = q_j$,
but $\delta(\varphi(q_i), a_{ji}) = q_i$.
\end{example}



\subsection{Morphisms and the Reachability Order}

In this section, let $\mathcal A = (Q, \Sigma, \delta)$ and $\mathcal B = (S, \Sigma, \mu)$
be two automata.

First, we show that two morphisms agree on an upper set of the form $\upperset{\{p\}}$
if and only if they agree at~$p$.

\begin{proposition} \label{prelim:automata:morphism:cor:two-morp-upper-set}
For morphisms $\varphi, \psi \colon Q \to S$ and $p \in Q$, we have
  \begin{equation}
	\varphi(p) = \psi(p) \Rightarrow \forall q \in \upperset{\{p\}} \colon \varphi(q) = \psi(q) \label{prelim:eqn:two-morph-equal} 
  \end{equation}
 In particular, two morphism $\varphi, \psi$ agree on a strongly connected component
 or have different images for each state of the strongly connected component.
\end{proposition}
\begin{proof}
	Let $q \in \upperset{\{p\}}$.
	Then we have $q = \delta(p, u)$ for some $u \in \Sigma^*$.
	By \cref{prelim:automata:morphism:prop:principal-upper-set}
	we find $\varphi(\delta(p, u)) = \mu(\varphi(p), u) =  \mu(\psi(p), u) = \psi(\delta(p, u))$.
\end{proof}

For the remainder of this section, suppose $\varphi \colon Q \to S$ is a morphism.
Let $p,q \in Q$. 
Then $\varphi$ is an order morphism with respect to~$\lereach$.
Indeed, if $p \lereach q$, then $q = \delta(p,u)$ for some $u \in \Sigma^*$.
Consequently, 
$
\varphi(q) = \varphi(\delta(p,u)) = \mu(\varphi(p),u),
$
so $\varphi(p) \lereach \varphi(q)$.
Thus, we have shown that\footnote{This also 
	follows from \cref{prelim:aut:prop:morph-digraph}
	since a morphism maps paths to paths in the automaton digraphs.}
\begin{equation}
	p \lereach q \Rightarrow \varphi(p) \lereach \varphi(q). \label{prelim:eqn:morph-reachability-order}	
\end{equation}

\begin{proposition} \label{prelim:automata:morpisms:reachability:prop:images}
Let $\varphi \colon Q \to S$ be a morphism. %Then we have the following:

\begin{enumerate}
\item \label{prelim:automata:morpisms:reachability:prop:images:principal-upper-set} Let $T \subseteq Q$.
Then we have
\begin{equation}  \label{prelim:eqn:morph-principal-upper-set}
\varphi(\upperset{T}) = \upperset{\varphi(T)}
\end{equation}

\item  \label{prelim:fact:morphism-sink} \label{prelim:automata:morpisms:reachability:prop:images:preservation}
$\varphi(T)$ is an upper set (resp.~unilaterally connected, strongly connected, a sink component) for an upper set (resp.~unilaterally connected, strongly connected, sink component) $T \subseteq Q$.

\item \label{prelim:automata:morpisms:reachability:prop:images:twoSCC} For two strongly connected components $T \subseteq Q$ and $R \subseteq S$,
 we have $\varphi(T) \cap R = \emptyset$ or $\varphi(T) \subseteq R$.
 
\item \label{prelim:automata:morpisms:reachability:prop:images:inverse} Let $T \subseteq Q$
 be an upper set. Then $\varphi^{-1}(T)$ is also an upper set.
\end{enumerate}
\end{proposition}
\begin{proof}
\ref{prelim:automata:morpisms:reachability:prop:images:principal-upper-set}: 
We have
\begin{align*}
	\varphi(\upperset{T}) 
	& = \varphi(\{ \delta(p, u) \mid p \in T \land u \in \Sigma^* \}) \\
	& = \{ \varphi(\delta(p, u)) \mid p \in T \land u \in \Sigma^* \} \\
	& = \{ \mu(\varphi(p), u) \mid p \in T \land u \in \Sigma^* \} \\
	& = \upperset{\varphi(T)}.
\end{align*}

\ref{prelim:automata:morpisms:reachability:prop:images:preservation}: Let $T \subseteq Q$ be an upper set. 
Then $T = \upperset{T}$
and, with \cref{prelim:automata:morpisms:reachability:prop:images:principal-upper-set}, we find 
$\varphi(T) = \varphi(\upperset{T}) = \upperset{\varphi(T)}$.
Since $\upperset{\varphi(T)} \subseteq \varphi(T)$, $\varphi(T)$ is an upper set.

For unilaterally connected~$T$, consider $\varphi(p), \varphi(q)$ for $p,q \in T$.
Since $p \lereach q$ or $q \lereach p$, we have $\varphi(p) \lereach \varphi(q)$
or $\varphi(q) \lereach \varphi(p)$ and so~$\varphi(T)$
is unilaterally connected.

By \cref{prelim:prop:morphism_equiv_class}, $\varphi(T)$ is strongly connected
if~$T$ is strongly connected.

Lastly, suppose~$T$ is a sink component.
Then $\varphi(T)$ is strongly connected and an upper set
since both properties are preserved as shown above.

\ref{prelim:automata:morpisms:reachability:prop:images:twoSCC}: 
This follows from \cref{prelim:general:order:prop:equiv-class-image-disjoint-or-subset}.
\end{proof}

For general order morphisms, \Cref{prelim:eqn:morph-principal-upper-set} may fail and 
the images of upper sets need not be upper sets, see \cref{prelim:general:order:rem:morphism-upperset}.

\begin{corollary} \label{prelim:automata:morpisms:reachability:cor:images}
	Let $T \subseteq Q$.
	If $\varphi(T) \subseteq R$ for a sink component~$R$,
	then $\varphi(\upperset{T}) = R$.
\end{corollary}
\begin{proof}
	Since $\varphi(T) \subseteq R$
	we have $\upperset{\varphi(T)} \subseteq \upperset{R} = R$.
	Since~$R$ is strongly connected and $\upperset{\varphi(T)} \cap R \ne \emptyset$,
	we have $R \subseteq \upperset{\varphi(T)}$.
	Thus, with \cref{prelim:automata:morpisms:reachability:prop:images},
	we find $\varphi(\upperset{T}) = \upperset{\varphi(T)} = R$.
\end{proof}

\begin{remark}
	Consider the automata in \cref{prelim:aut:fig:mappingmorph,prelim:aut:fig:mappingmorph0} 
	from \cref{prelim:ex:SCC2}. 
	The image of $\{q_0\}$ is $\{q_1\}$, which is strictly
	contained in the sink component $\{q_1, q_2\}$,
	and the image of $\upperset{\{q_0\}}$ equals $\{q_1, q_2\}$.
	Note that $\{q_0\}$ is a strongly connected component,
	but its image is not a strongly connected component in both automata.
	
	In \cref{prelim:aut:fig:mappingotmorph}, $\{q_0\}$
	is a sink component and its image is not a sink component.
	Thus, in this case, we also obtain that the mapping
	is not a morphism from \cref{prelim:automata:morpisms:reachability:cor:images}.
\end{remark}

In general, bijective order morphisms
do not have to be order isomorphisms, see \cref{prelim:ex:order:bijmorph}.
By the next statement, with \cref{prelim:prop:iso-bijmorph}, bijective order 
morphisms arising from automata morphisms are in fact order isomorphisms.

\begin{proposition} \label{prelim:automata:morphism:prop:injective}
	If $\varphi \colon Q \to S$ is an isomorphism,
	then it is also an order isomorphism.
\end{proposition}
\begin{proof}
	By assumption, there exists a morphism $\psi \colon S \to Q$
	such that $\concat{\varphi}{\psi} = \operatorname{id}_Q$
	and $\concat{\psi}{\varphi} = \operatorname{id}_S$.
	By \cref{prelim:eqn:morph-reachability-order}, $\psi$
	is also an order morphism.
	Thus, $\varphi$ is an order isomorphism.	
\end{proof}

\begin{corollary} \label{prelim:automata:morphism:prop:injective:SCC-to-SCC}
	Suppose $\varphi \colon Q \to S$ is an isomorphism.
	If $T \subseteq Q$ is a strongly connected component,
	then $\varphi(T)$ is a strongly connected component. 
\end{corollary}
\begin{proof}
	This follows from \cref{prelim:prop:morphism_equiv_class,prelim:automata:morphism:prop:injective}.
\end{proof}


\begin{proposition} \label{prelim:automata:morphism:prop:image-SCC}
\begin{enumerate}
\item \label{prelim:automata:morphism:prop:image-SCC:1} If $Q \ne \emptyset$ and $\mathcal B$ is strongly connected, or strongly trap-connected where $\varphi(Q)$ has size at least two,
then $\varphi$ is surjective, i.e., $\varphi(Q) = S$.

\item \label{prelim:automata:morphism:prop:image-SCC:2} If $\mathcal A$ is strongly connected (resp.~strongly trap-connected)
and $\varphi$ is surjective, 
then~$\mathcal B$ is strongly connected (resp.~strongly trap-connected). 
\end{enumerate}
\end{proposition}
\begin{proof}
\ref{prelim:automata:morphism:prop:image-SCC:1}: By \cref{prelim:automata:morpisms:reachability:prop:images:preservation} of \cref{prelim:automata:morpisms:reachability:prop:images}, $\varphi(Q)$ is an upper set.
 First, suppose $\mathcal B$ is strongly connected.
 By \cref{prelim:automata:prop:char-strongly-conn}, a strongly connected automaton has no proper non-empty upper
 set. Thus $\varphi(Q) = S$.
 Next, suppose~$\mathcal B$ is strongly trap-connected and $|\varphi(Q)| \ge 2$.
 By \cref{prelim:automata:prop:char-strongly-trap}, $\mathcal B$
 has no proper upper set of size at least two.
 Thus $\varphi(Q) = S$.

\ref{prelim:automata:morphism:prop:image-SCC:2}: We can assume $|Q|, |\varphi(Q)| \ge 2$.

If~$\mathcal A$ is strongly connected, then~$Q$ is a sink component.
Thus, by \cref{prelim:automata:morpisms:reachability:prop:images:preservation} from \cref{prelim:automata:morpisms:reachability:prop:images},
$\varphi(Q)$ is also a sink component. By surjectivity, $\mathcal B$ is strongly connected.

Next, suppose~$\mathcal A$ is strongly trap-connected.
Let $T \subseteq S$ be an upper set in $\mathcal B$
of size at least two.
By \cref{prelim:general:order:eqn:morphism-upperset-inverse}, $\varphi^{-1}(T)$
is an upper set. Since~$\varphi$ is surjective, $|\varphi^{-1}(T)| \ge 2$.
By \cref{prelim:automata:prop:char-strongly-trap} we find $Q = \varphi^{-1}(T)$.
Thus $S = \varphi(Q) \subseteq T$.
Since the trap state in $\mathcal A$ is mapped to a trap state in~$\mathcal B$,
$\mathcal B$ is not strongly connected.
Then, \cref{prelim:automata:prop:char-strongly-trap} gives
that $\mathcal B$ is strongly trap-connected.
\end{proof}

\begin{corollary} \label{prelim:automata:morphism:cor:image-SCC}
	If $|Q| = |S|$ and~$\mathcal B$ is strongly connected,
	or strongly trap-connected with $|\varphi(Q)| \ge 2$,
	then~$\varphi$ is an isomorphism (and so~$\mathcal A$ is also strongly connected or strongly trap-connected).
\end{corollary}
\begin{proof}
	By \cref{prelim:automata:morphism:prop:image-SCC:1} from \cref{prelim:automata:morphism:prop:image-SCC}, 
	we have $\varphi(Q) = S$.
	Since $\rank{\varphi} = |Q|$,
	by \cref{prelim:prop:item:rank-defect-eqn}
	from \cref{prelim:prop:rank-defect} we find $\defect{\varphi} = \emptyset$.
	Thus, $\varphi$ is bijective.
	By \cref{prelim:prop:iso-bijmorph}, $\varphi$ is an isomorphism.
	
	The inverse map of~$\varphi$ is an isomorphism.
	Since~$\mathcal B$ is strongly connected or strongly trap-connected,
	by \cref{prelim:automata:morphism:prop:image-SCC:2} from \cref{prelim:automata:morphism:prop:image-SCC},
	we find that $\mathcal A$ is strongly connected
	or strongly trap-connected.		
\end{proof}

\subsection{Endomorphisms}
\label{prelim:sec:endo}

By the first result of this section, \cref{prelim:automata:prop:endo-reach-order},
an endomorphism cannot map a state to one that strictly precedes it in the
reachability order.
This property
does not hold for general quasiorders, see \cref{prelim:general:order:rem:morphism-upperset}.
We derive several consequences of this fact. For example,
\cref{prelim:automata:prop:endo:SCC}, showing that an endomorphism permutes a
strongly connected component~$T$ if and only if it permutes the upper set
generated by~$T$, or
\cref{prelim:automata:cor:SCC-disjoint-or-permutates} showing that an
endomorphism either permutes a strongly connected component or maps all of
its states to states outside the component.
Lastly, we give several results on endomorphisms of strongly connected
and synchronizing automata.

In this section, let $\semiautomat{A}{Q}{\Sigma}{\delta}$ be an automaton and $\varphi \colon Q \to Q$ be an endomorphism. 

\begin{proposition} \label{prelim:automata:prop:endo-reach-order}
 Let $p \in Q$.
 Then $\varphi(p) \lereach p$ implies $p \lereach \varphi(p)$.	
\end{proposition}
\begin{proof}
	Suppose $\varphi(p) \lereach p$.
	By \cref{prelim:eqn:reach-upperset} 
	we have $\upperset{\{p\}}\subseteq \upperset{\{\varphi(p)\}}$.
	By \cref{prelim:eqn:morph-principal-upper-set} 
	we have $\upperset{\{\varphi(p)\}} = \varphi(\upperset{\{p\}})$.
	Thus $\upperset{\{p\}}\subseteq \varphi(\upperset{\{p\}})$.
	By finiteness, this gives 
	$\upperset{\{p\}} = \varphi(\upperset{\{p\}}) = \upperset{\{\varphi(p)\}}$.
	\Cref{prelim:general:order:eqn:equiv-uppersets}
	we find $p \eqreach \varphi(p)$.
\end{proof}

\begin{corollary}
	Either $p$ and $\varphi(p)$ are incomparable or we have $p \lereach \varphi(p)$
\end{corollary}

\begin{corollary} \label{prelim:automata:prop:endo:SCC}
	For a strongly connected component $T \subseteq Q$,
	the following are equivalent:
	\begin{enumerate}
		\item \label{prelim:facts:permutes_SCC:1a} $T \subseteq \varphi(\upperset{T})$,
		
		\item \label{prelim:facts:permutes_SCC:3} $\varphi(T) = T$,

		\item \label{prelim:facts:permutes_SCC:2} $\varphi(T) \cap T \ne \emptyset$,
		
		\item \label{prelim:facts:permutes_SCC:1} $\varphi(\upperset{T}) = \upperset{T}$.
	\end{enumerate}	
\end{corollary}
\begin{proof}

$\ref{prelim:facts:permutes_SCC:1a} \Rightarrow \ref{prelim:facts:permutes_SCC:3}$:
Let $p \in T$. By assumption, we have $p = \varphi(q)$ for some $q \in \upperset{T}$.
Since~$T$ is a strongly connected component,
we have $\upperset{T} = \upperset{\{p\}}$.
Thus, there exists $u \in \Sigma^*$ such that $q = \delta(p, u)$.
This gives $p = \varphi(\delta(p, u))$.
By \cref{prelim:prop:reachability-finest} we have 
$p \lereach \delta(p, u)$.
Hence, $\varphi(\delta(p, u)) \lereach \delta(p, u)$.
By \cref{prelim:automata:prop:endo-reach-order}
this gives $\varphi(\delta(p, u)) \eqreach \delta(p, u)$.
Thus $p \eqreach \delta(p, u)$
and so $q \in T$. Hence $p \in \varphi(T)$.
This gives $T \subseteq \varphi(T)$.
Since $|\varphi(T)| \le |T|$ by \cref{eqn:subset-mapped-to-smaller-set},
we find $T = \varphi(T)$.	

$\ref{prelim:facts:permutes_SCC:3} \Rightarrow \ref{prelim:facts:permutes_SCC:2}$: This is clear.

$\ref{prelim:facts:permutes_SCC:2} \Rightarrow \ref{prelim:facts:permutes_SCC:1}$: Let $p \in \varphi(T) \cap T$.
By \cref{prelim:general:order:prop:equiv-class-image-disjoint-or-subset} 
we have $\varphi(T) \subseteq T$.
Thus $\varphi(p) \in T$.
Since~$T$ is a strongly connected
component, we have $\upperset{T} = \upperset{\{p\}} = \upperset{\{\varphi(p)\}}$. So
\[
\varphi(\upperset{T})
= \varphi(\upperset{\{p\}})
= \upperset{\{\varphi(p)\}}
= \upperset{T}.
\]
This gives the claim.\footnote{Alternatively,
	since $\varphi(\upperset{T})$ is an upper set having non-empty intersection with~$T$,
	we have $T \subseteq \varphi(\upperset{T})$.
	As $\upperset{T}$ is the smallest upper set containing~$T$,
	this gives $\upperset{T}\subseteq \varphi(\upperset{T})$.
	By finiteness, we then obtain 
	$\upperset{T} = \varphi(\upperset{T})$.}
	
$\ref{prelim:facts:permutes_SCC:1} \Rightarrow \ref{prelim:facts:permutes_SCC:1a}$:	This is clear.
\end{proof}

\begin{corollary} \label{prelim:automata:cor:SCC-disjoint-or-permutates}
	For a strongly connected component $T \subseteq Q$,
	we have $\varphi(T) \cap T = \emptyset$
	or $\varphi(T) = T$.
\end{corollary}
\begin{proof}
	This directly follows from Items~\ref{prelim:facts:permutes_SCC:3} and~\ref{prelim:facts:permutes_SCC:2}
	of the previous corollary.	
\end{proof}

\begin{remark}	
	In general, for a sink component $T \subseteq Q$, a morphism $\varphi$ need not induce an isomorphism between
	$\mathcal A_{|T}$ and $\mathcal A_{|\varphi(T)}$.
	For example, this fails when $|\varphi(T)| = 1$ and $|T| > 1$.
	However, by \cref{prelim:automata:cor:SCC-disjoint-or-permutates}, if
	$\varphi(T) \cap T \neq \emptyset$, then $\varphi(T)=T$.
	Consequently, the restriction $\varphi_{|T}$ is an automorphism of
	$\mathcal A_{|T}$.
\end{remark}

\begin{corollary} \label{prelim:automata:morphism:prop:endo-SCC-is-iso}
	An endomorphism of a strongly connected finite automaton is an automorphism.
\end{corollary}
\begin{proof}
	By \cref{prelim:automata:cor:SCC-disjoint-or-permutates},
	it follows that every endomorphism is bijective, and so an automorphism by \cref{prelim:prop:iso-bijmorph}.  
	
	Alternatively, this also follows from \cref{prelim:automata:morphism:prop:image-SCC:1} 
	of \cref{prelim:automata:morphism:prop:image-SCC}.
\end{proof}

There exist automata that are not strongly connected and whose non-identity endomorphisms are all automorphisms. For example, a unary automaton in which the single letter induces two cycles of coprime lengths.

\begin{corollary} \label{prelim:automata:morphim:cor:strongly-conn}
	Let $S \subseteq Q$ be strongly connected.
	Then $\varphi(S) \subseteq S$
	implies $\varphi(S) = S$.
\end{corollary}
\begin{proof}
	We have $S \subseteq T$ for some strongly connected component $T \subseteq Q$.
	Let $p \in S$. By assumption, we have $\varphi(p) \in T$.
	Thus $\varphi(p) \in T \cap \varphi(T)$.
	By \cref{prelim:automata:cor:SCC-disjoint-or-permutates}
	we find $\varphi(T) = T$.
	Hence~$\varphi$ permutes the states in~$T$
	and so $\varphi_{|T}$ is injective.
	Then $\varphi(S) \subseteq S$ gives $\varphi(S) = S$.
\end{proof}

\begin{corollary}
	Let $\psi \colon Q \to Q$ be another endomorphism
	and $S \subseteq Q$ be strongly connected.
	If $\psi(\varphi(S)) \subseteq \varphi(S)$,
	then $\psi(\varphi(S)) = \varphi(S)$.	
\end{corollary}
\begin{proof}
	By \cref{prelim:automata:morpisms:reachability:prop:images:twoSCC} of \cref{prelim:automata:morpisms:reachability:prop:images}
	the set $\varphi(S)$ is strongly connected.
	Then the claim follows from \cref{prelim:automata:morphim:cor:strongly-conn}.	
\end{proof}

\begin{corollary}
	For $\varphi(S) \subseteq S$ and $\delta(S, u) \subseteq T$ with strongly connected components~$S, T \subseteq Q$ 
	and $u \in \Sigma^*$ we have $\varphi(T) = T$.	
\end{corollary}
\begin{proof}
	Let $q \in S$. Then $\varphi(q) \in S$.
	Since $\delta(S, u) \subseteq T$,
	we find $\delta(q, u) \in T$ and $\varphi(\delta(q, u)) = \delta(\varphi(q), u) \in T$.	
	Thus $\varphi(\delta(q, u)) \in \varphi(T) \cap T$.
	By \cref{prelim:automata:prop:endo:SCC},
	this gives $\varphi(T) = T$.
\end{proof}

By \cref{prelim:automata:prop:endo:SCC},
for $\varphi(\upperset{\{p\}}) = \upperset{\{p\}}$, $\varphi$ permutes the strongly connected component
$[p]_{\eqreach}$.
However, this does not imply
that $\varphi$ permutes each strongly connected component in $\upperset{\{p\}}$,
as \cref{prelim:ex:endo} shows next.

\begin{example} \label{prelim:ex:endo}
	Let $\semiautomat{A}{Q}{\{a,b\}}{\delta}$
	with $Q = \{q_1, q_2, q_3, q_4\}$ and
	\[
	\delta(q_i, a) = 
	\begin{cases}
		q_4 & \text{if } i = 1, \\
		q_3 & \text{if } i = 2, \\
		q_i & \text{otherwise,} 
	\end{cases}
	\qquad 
	\delta(q_i, b) =
	\begin{cases}
		q_1  & \text{if } i = 2, \\
		q_2  & \text{if } i = 1, \\
		q_i  & \text{otherwise.} 
	\end{cases}
	\]
	The strongly connected components are $\{q_1, q_2\}$, $\{q_3\}$ and $\{q_4\}$.
	We have $\upperset{\{q_1\}} = Q$.
	The map $\varphi$ given by $\varphi(q_1) = q_2$, $\varphi(q_2) = q_1$, $\varphi(q_3) = q_4$
	and $\varphi(q_4) = q_3$ is a morphism
	that maps $\{q_1, q_2\}$ to itself and interchanges $\{q_3\}$ and $\{q_4\}$.
	
	\begin{figure}[ht]
		\centering	
		\begin{tikzpicture}[node distance = 2cm, on grid, auto, ->]
			\node (q1) [state] {$q_1$};
			\node (q2) [state, right of = q1] {$q_2$};
			\node (q3) [state, right of = q2] {$q_3$};
			\node (q4) [state, left of  = q1] {$q_4$};
			
			\path (q1) edge [above] node {$a$} (q4)
			      (q2) edge node {$a$} (q3)
			      (q3) edge [loop right] node {$a,b$} (q3)
			      (q4) edge [loop left]  node {$a,b$} (q4);
			\path (q1) edge [bend left] node {$b$} (q2)
			      (q2) edge [bend left] node {$b$} (q1); 
		\end{tikzpicture}
		\caption[Some strongly connected component fixed and others swapped]{Automaton admitting a morphism that leaves $\{q_1, q_2\}$ invariant and swaps $\{q_4\}$ and $\{q_3\}$, see \cref{prelim:ex:endo}.}
		\label{prelim:aut:fig:autnomorph2}
	\end{figure}
\end{example} 

\begin{proposition} \label{prelim:automata:prop:endo-SCC-fixed-all-or-none}
	For a strongly connected component $T \subseteq Q$, 
	either we have $\varphi(p) = p$ for all $p \in T$
	or $\varphi(p) \ne p$ for all $p \in T$.
\end{proposition}
\begin{proof}
	Suppose we have $\varphi(p) = p$ for some $p \in T$.
	Since $T \subseteq \upperset{\{p\}}$,
	from \cref{prelim:automata:morphism:cor:two-morp-upper-set} 
	applied to~$\varphi$ 
	and the identity transformation~$\operatorname{id}_Q$
	we get $\varphi(q) = q$ for all $q\in \upperset\{p\}$,
	in particular~$\varphi$ fixes all states of~$T$.	
\end{proof}

\begin{corollary} \label{prelim:automata:morphism:prop:endo-SCC-is-iso-b}
    The automorphism group of a strongly connected automaton is semi-regular (and so has size at most~$|Q|$).
\end{corollary}
\begin{proof}
	By  \cref{prelim:automata:prop:endo-SCC-fixed-all-or-none}, every non-identity isomorphism
	fixes no state.
	Thus, the automorphism group is semiregular.
	By \cref{prelim:groups:prop:semiregular-size}, 
	its size is at most~$|Q|$.	
\end{proof}

\begin{proposition} \label{prelim:automata:prop:endo-SCC-fixed-sync-state}
	$\varphi \colon Q \to Q$ maps a synchronizing state to itself.
\end{proposition}
\begin{proof}
	Let $u \in \Sigma^*$ be such that $\delta(Q, u) = \{q\}$.
	Then we have $\varphi(q) = \varphi(\delta(q, u)) = \delta(\varphi(q), u) = q$.
\end{proof}

\begin{corollary} \label{prelim:automata:morphism:cor:endo-sync-SCC}
	The only endomorphism of a strongly connected synchronizing automaton
	is the identity transformation.
\end{corollary}
\begin{proof}
	By \cref{prelim:automata:prop:endo-SCC-fixed-sync-state}
	an endomorphism fixes the synchronizing state.
	Thus, by \cref{prelim:automata:prop:endo-SCC-fixed-all-or-none},
	it equals the identity transformation.\footnote{Alternatively,
	since in this case every state is a synchronizing state
	as we show later in \cref{prelim:cor:SCC_sync},
	from \cref{prelim:automata:prop:endo-SCC-fixed-sync-state}
	we also get that the endomorphism fixes every state.}		
\end{proof}

\begin{corollary}
	For a strongly connected synchronizing automaton
	having the same number of states as~$\mathcal A$,
	there exists at most one morphism from~$\mathcal A$ to that automaton.
\end{corollary}
\begin{proof}
	Let $\semiautomat{B}{S}{\Sigma}{\mu}$ be strongly connected
	and synchronizing with $|Q| = |S|$.
	Suppose $\psi, \xi \colon Q \to S$ are two morphisms.
	Since~$\mathcal B$ is strongly connected,
	by \cref{prelim:automata:morphism:cor:image-SCC}, 
	they are both isomorphisms.
	Let $\psi^{-1} \colon S \to Q$ be the inverse morphism of~$\psi$.
	Then $\concat{\psi^{-1}}{\xi}$ is an endomorphism of~$\mathcal B$.
	By \cref{prelim:automata:morphism:cor:endo-sync-SCC} 
	it equals the identity transformation, which implies $\psi = \xi$.
\end{proof}

\subsection{Cyclic Unary Automata}

As a simple example, here we classify all morphisms
between cyclic unary automata. First, \cref{prelim:automata:lem:morphism-unary-automata,prelim:automata:prop:morphism-unary-automata}
are preliminary statements, and \cref{prelim:automata:cor:morphism-unary-automata} 
gives the classification.

\begin{lemma} \label{prelim:automata:lem:morphism-unary-automata}
	Let $\semiautomat{A}{Q}{\{a\}}{\delta}$
	and $\semiautomat{B}{S}{\{a\}}{\mu}$
	where $Q = \upperset{\{p\}}$ for some $p \in Q$.
	Suppose $\varphi \colon Q \to S$ is a morphism.
	Let $i_2 \ge 0$ be the least number such that,
	with a corresponding least number $p_2 > 0$,
	we have $\mu(\delta(p), a^{i_2 + p_2}) = \mu(\delta(p), a^{i_2})$.
	Then for all  $i_1 \ge 0$ and $p_1 > 0$
	with $\delta(p, a^{i_1 + p_1}) = \delta(p, a^{i_1})$ 
	we have  $i_2 \le i_1$ and that $p_1$ is a multiple of~$p_2$. 
\end{lemma}
\begin{proof}
	First, $\delta(p, a^{i_1 + p_1}) = \delta(p, a^{i_1})$
	implies $\mu(\varphi(p), a^{i_1 + p_1}) = 
	\varphi(\delta(p, a^{i_1 + p_1})) = \varphi(\delta(p, a^{i_1})) = \mu(\varphi(p), a^{i_1})$.
	Hence, by the minimal choice of~$i_2$ and~$p_2$,
	we have $i_2 \le i_1$ and $p_2 \le p_1$.
	Next, write $p_1 = l p_2 + r$
	for some $l \ge 0$ and $0 \le r < p_2$.
	Then 
	we have $\mu(\delta(p), a^{i_2 + p_1}) = \mu(\delta(p), a^{i_2 + r})$
	and, by minimality, we must have $r = 0$, i.e., $p_1$
	is a multiple of~$p_2$.	
\end{proof}

\begin{example}
	Let $\mathcal A = (Q, \{a\}, \delta)$ consist of two disjoint
	cycles for~$a$ of lengths $p$ and~$q$
	for two distinct prime numbers.
	Then endomorphisms mapping some state from one cycle
	to the other cycle are not possible by \cref{prelim:automata:lem:morphism-unary-automata}.
	With \cref{prelim:automata:cor:endo-unary-automata}, each endomorphism
	shifts each state on the cycle of length~$p$
	and on the cycle of length~$q$ by some
	amount~$0 \le i < p$ and $0 \le j < q$, respectively.
	So, $\mathcal A$
	has~$pq$ endomorphisms.
\end{example}

\begin{corollary} \label{prelim:automata:cor:morphism-unary-automata}
	Let $\semiautomat{A}{Q}{\{a\}}{\delta}$ 
	where $Q = \upperset{\{p\}}$ for some $p \in Q$
	and $i_1 \ge 0$ be the least number
	with a corresponding least number $p_1 > 0$
	such that $\delta(p, a^{i_1+p_1}) = \delta(p, a^{i_1})$.
	Then $\delta(p, a^{i_2 + p_2}) = \delta(p, a^{i_2})$
	implies $i_1 \le i_2$ and that $p_2$ is a multiple of~$p_1$.
\end{corollary}
\begin{proof}
	Apply \cref{prelim:automata:lem:morphism-unary-automata} with the identity transformation, which is a morphism.
\end{proof}

\begin{proposition} \label{prelim:automata:prop:morphism-unary-automata}
	Let $\semiautomat{A}{Q}{\{a\}}{\delta}$
	with $Q = \upperset{\{q\}}$ for some $q \in Q$
	and $\semiautomat{B}{S}{\{a\}}{\mu}$.
	Then $\varphi \colon Q \to S$ is a morphism if and only if $\mu(\varphi(q), a^{i_2 + p_2}) = \mu(\varphi(q), a^{i_2})$
	for some $0 \le i_2 \le i_1$ with $p_1$ being a multiple of~$p_2$
	and $\varphi(\delta(q, a^i)) = \mu(\varphi(q), a^i)$ for $i \in\{0,1,\ldots,n-1\}$
	with $n = |Q|$.
\end{proposition}
\begin{proof}
	First, let $\varphi$ be a morphism.
	By finiteness, we have  $\mu(\varphi(q), a^{i_2 + p_2}) = \mu(\varphi(q), a^{i_2})$
	for some minimal $i_2 \ge 0$ and some $p_2 > 0$.
	By \cref{prelim:automata:lem:morphism-unary-automata},
	we find $i_2 \le i_1$ and that $p_1$ is a multiple of~$p_2$.
	Since $\varphi$ is a morphism,
	we have $\varphi(\delta(q, a^i)) = \mu(\varphi(q), a^i)$ for all $i \ge 0$.
	
	Next, suppose we have $\mu(\varphi(q), a^{i_2 + p_2}) = \mu(\varphi(q), a^{i_2})$
	for some $0 \le i_2 \le i_1$ with $p_1$ being a multiple of~$p_2$
	and  $\varphi(\delta(q, a^i)) = \mu(\varphi(q), a^i)$ for $i \in\{0,1,\ldots,n-1\}$.
	We have $n = i_1 + p_1$.
	By assumption, there exists $k \ge 1$ such that $p_1 = k\cdot p_2$.
	Since $i_2 \le i_1$ 
	we have $\mu(\varphi(q), a^{i_1}) = \mu(\varphi(q), a^{i_1 + k\cdot p_2}) = \mu(\varphi(q), a^n)$.
	Then we find	
	\[
	\varphi(\delta(q, a^{n - 1}), a) = \varphi(\delta(q, a^{i_1}))
	= \mu(\varphi(q), a^{i_1})
	= \mu(\varphi(q), a^n)
	= \mu(\delta(q, a^{n - 1}), a).
	\]
	Since $Q = \{q, \delta(q, a), \ldots, \delta(q, a^{n-1})\}$, 
	we have $\varphi(\delta(p, a)) = \delta(\varphi(p), a)$
	for all $p \in Q$.
\end{proof}

\begin{corollary} \label{prelim:automata:cor:morphism-unary-automata}
	Let $\semiautomat{A}{Q}{\{a\}}{\delta}$
	and $\semiautomat{B}{S}{\{a\}}{\mu}$
	where $Q = \upperset{\{p\}}$ 
	and $S = \upperset{\{s\}}$
	for some $q \in Q$ and $s \in S$.
	Let  $i_1, i_2 \ge 0$ be the least numbers with corresponding least numbers $p_1,p_2 > 0$ 
	such that $\delta(q, a^{i_1 + p_1}) = \delta(q, a^{i_1})$
	and $\mu(s, a^{i_2 + p_2}) = \mu(s, a^{i_2})$.
	Then there exist $\min\{i_1, i_2\} + p_2$
	morphisms from~$\mathcal A$ to~$\mathcal B$ 
	if and only if $p_1$ is a multiple of~$p_2$,
	and there exists no morphism otherwise.
\end{corollary}
\begin{proof}
	If $p_1$ is not a multiple of~$p_2$,
	there exists no morphism by \cref{prelim:automata:cor:morphism-unary-automata}.
	Next, suppose $p_1$ is a multiple of~$p_2$.
	Then, for $i_2 - \min\{i_1, i_2\} \le i < |S|$,
	setting $\varphi(\delta(q, a^j)) = \mu(s, a^{i + j})$,
	which gives  $\varphi(\delta(q, a^j)) = \mu(\varphi(q), a^j)$,
	for $j \in \{0,1,\ldots,|Q|-1\}$
	yields a morphism by \cref{prelim:automata:cor:morphism-unary-automata}
	where $\mu(\varphi(q), a^{\min\{i_2 - i, 0\} + p_2}) = \mu(\varphi(q), a^{\min\{i_2 - i,0\}})$
	and $\min\{i_2 - i, 0\} \le \min\{i_1, i_2\} \le i_1$.
	
	Next, let $\varphi \colon Q \to S$ be a morphism
	and $0 \le i <|S|$ such that $\varphi(q) = \mu(s, a^i)$.
	We have $\mu(s, a^{i + j}) = \mu(s, a^{i + j + p_2})$
	for $j = \min\{i_2 - i, 0\}$.
	By \cref{prelim:automata:cor:morphism-unary-automata} we find $j \le i_1$.
	Thus, for $i_2 > i_1$ we must have $i_2 - i_1 \le i$
	since otherwise $i + j < i_2$.
	This gives $i_2 - \min\{i_1, i_2\} \le i < |S|$.
	
	So, each morphism is uniquely
	given by its image 
	$\varphi(q) = \mu(s, a^i)$
	for $i_2 - \min\{i_1, i_2\} \le i < |S|$.
	Thus for $i_2 < i_1$
	we have $|S| = i_2 + p_2 = \min\{i_1, i_2\} + p_2$ 
	possible morphisms.
	For $i_2 \ge i_1$
	we have $|\{ i \mid i_2 - i_1 \le i < i_2 + p_2 \}| =i_1 + p_2 = \min\{i_1, i_2\} + p_2$
	possible morphisms.	
\end{proof}

\begin{corollary} \label{prelim:automata:cor:endo-unary-automata}
	Let $\semiautomat{A}{Q}{\{a\}}{\delta}$
	with $Q = \upperset{\{q\}}$ for some $q \in Q$.
	Then $\varphi \colon Q \to Q$ is an endomorphism
	if and only if there exists $0 \le m < |Q|$ such that
	$\varphi(p) = \delta(p, a^m)$
	for all $p \in Q$.
\end{corollary}
\begin{proof}
	Let $m \ge 0$ such that $\varphi(q) = \delta(q, a^m)$.
	By \cref{prelim:automata:cor:morphism-unary-automata}, $\varphi$
	is a morphism if and only if $\varphi(\delta(q, a^i)) = \delta(\varphi(q), a^i)$
	for $0 \le i < |Q|$.
	This is equivalent to $\varphi(p) = \delta(p, a^m)$
	for all $p \in Q$ since $\delta(\varphi(q), a^i) = \delta(\delta(q, a^i), a^m)$.
\end{proof}

Since in unary automata, strongly connected components are cycles, by \cref{prelim:automata:lem:morphism-unary-automata}, 
the size of a strongly connected component in the image of a morphism divides the size of a strongly connected
component that maps to it. This might fail for non-unary automata, as shown next.

\begin{example} \label{prelim:aut:ex:morphism-not-divide}
	Let $\semiautomat{A}{\{q_0, q_1, q_2\}}{\{a,b\}}{\delta}$
	and
	$\semiautomat{B}{\{p_0, p_1\}}{\{a,b\}}{\mu}$
	with
	$\delta(q_0, a) = q_1$, $\delta(q_1, a) = q_2$,
	$\delta(q_1, b) = q_0$, $\delta(q_2, b) = q_0$,
	$\mu(p_0, a) = p_1$, $\delta(p_1, b) = p_0$
	and self-loops for all the remaining transitions, see \cref{prelim:aut:fig:morphism-not-divide}.
	The map~$\varphi$ with $\varphi(q_0) = p_0$, $\varphi(q_1) = p_1$,
	$\varphi(q_2) = p_1$
	is a morphism since $\mu(\varphi(q), b) = p_0 = \varphi(q_0) = \varphi(\delta(q, a))$
	for each $q \in \{q_0,q_1,q_2\}$ and
	\begin{align*}
		\mu(\varphi(q_0), a) & = \mu(p_0, a) = p_1 = \varphi(q_1) = \varphi(\delta(q_0, a)), \\
		\mu(\varphi(q_1), a) & = \mu(p_1, a) = p_1 = \varphi(q_2) = \varphi(\delta(q_1, a)), \\
		\mu(\varphi(q_2), a) & = \mu(p_1, a) = p_1 = \varphi(q_2) = \varphi(\delta(q_2, a)).
	\end{align*}	
	
	\begin{figure}[ht]
		\centering	
		\begin{tikzpicture}[->, >=Stealth, node distance=8em, on grid, auto]
			
			% --- Automaton A ---
			\node[state] (q0) {$q_0$};
			\node[state, right=of q0] (q1) {$q_1$};
			\node[state, right=of q1] (q2) {$q_2$};
			
			% Transitions in A
			\path
			(q0) edge[bend left=10] node {$a$} (q1)
			(q1) edge[bend left=10] node {$b$} (q0)
			(q1) edge[bend left=10] node {$a$} (q2)
			(q2) edge[bend right=32, above] node {$b$} (q0);
			
			%	% Self-loops in A
			%	(q0) edge[loop above] node {$b$} ()
			%	(q2) edge[loop above] node {$a$} ();
			
			% Label A
			\node[right=14em of q1] {Automaton $\mathcal A$};
			
			% --- Automaton B ---
			\node[state, below=7em of q0] (p0) {$p_0$};
			\node[state, right=of p0] (p1) {$p_1$};
			
			% Transitions in B
			\path
			(p0) edge[bend left=10] node {$a$} (p1)
			(p1) edge[bend left=10] node {$b$} (p0);
			
			%	% Self-loops in B
			%	\path
			%	(p0) edge[loop below] node {$b$} ()
			%	(p1) edge[loop below] node {$a$} ();
			
			% Label B
			\node[right=6em of p1] {Automaton $\mathcal B$};
			
			% --- Morphism phi ---
			\draw[dashed]
			(q0) -- node[left=0.6cm] {$\varphi$} (p0);
			\draw[dashed] (q1) -- (p1);
			\draw[dashed] (q2) -- (p1);
			
		\end{tikzpicture}
		\caption[Morphism between automata]{Morphism $\varphi$ between two automata $\mathcal A$ and~$\mathcal B$, see \cref{prelim:aut:ex:morphism-not-divide}.}
		\label{prelim:aut:fig:morphism-not-divide}
	\end{figure}		
\end{example}

%\section{Automata Constructions}
%\section{More Operations on Automata}
\section{Operations on Automata}

\subsection{Subautomata and Restricted Automata}
\label{prelim:sec:aut:sub}

%!TEX root = ./../../thesis-trier.tex

\todo[inline]{hier an einigen stellen befehl semiautomat nicht benutzt, weil probleme wie unten.}

Let $\semiautomat{A}{Q}{\Sigma}{\delta}$ and $\Gamma \subseteq \Sigma$.

\begin{enumerate}
	
	\IN{restriction on automaton}
	\item The \emph{restriction of $\mathcal A$ to $\Gamma$} is $\semiautomat{A_{|\Gamma}}{Q}{\Gamma}{\delta_{|\Gamma}}$
	where $\delta_{|\Gamma} \colon Q \times \Gamma \to Q$
	is defined by $\applytrans{\delta_{|\Gamma}}{q}{x} = \applytrans{\delta}{q}{x}$ for all $q \in Q$
	and all $x \in \Gamma$.	\label{def:restriction_Gamma}
	
	\label{def:restriction_aut}
	\item Let $S \subseteq Q$ be an upper set in $\mathcal A$.
	The \emph{restriction of $\mathcal A$ to the upper set $S$}
	is the 
	deterministic and complete semiautomaton $\semiautomat{A_{|S}}{S}{\Sigma}{\delta_{|S}}$ 
	with $\delta_{|S} \colon S \times \Sigma \to S$
	given by
	$\applytrans{\delta_{|S}}{q}{x} = \applytrans{\delta}{q}{x}$
	for $q \in S$ and $x \in \Sigma$.
	We also call~$\mathcal A_{|S}$ the \emph{sub(semi-)automaton} \IN{subautomaton}
	corresponding to (or induced/defined by) the upper set $S \subseteq Q$.
	We often just denote $\delta_{|S}$ by $\delta$ in this context.
		 
	\label{def:restriction_autSGamma}
	\item If $S \subseteq Q$ is an upper set in $\mathcal A_{|\Gamma}$,
	then by %$\semiautomat{A_{|S,\Gamma}}{S}{\Gamma}{\delta_{|S,\Gamma}}$
%	$\semiautomat{{\restrictAutGamS{\mathcal A}{\Gamma}{S}}{S}{\Gamma}{\delta_{|S,\Gamma}}$ %
	$\restrictAutGamS{\mathcal A}{\Gamma}{S} = (S, \Gamma, \delta_{|\Gamma, S})$	
	we denote the \emph{restriction of $\mathcal A$ to the the alphabet $\Gamma$ and the subset $S$}, i.e.,
	$(\mathcal A_{|\Gamma})_{|S}$. 
\end{enumerate}

Upper sets can be identified with subautomata. 
Let $\semiautomat{A}{Q}{\Sigma}{\delta}$
and $S, T \subseteq Q$ be upper sets. The following facts are readily seen.
We have
\begin{equation} \label{prelim:prop:automata:isos:item:subaut}
	(\mathcal A_{|S})_{|T} \cong \mathcal A_{|T}  \text{ for } T \subseteq S
\end{equation}
Clearly, $\mathcal A_{|S}$ is strongly connected if and only if~$S$
is a sink component.

\begin{proposition}
	\label{prelim:prop:automata:isos:item:corsub}
	The extension of the inclusion map $i \colon S \to Q$
	given by $i(q) = q$ for $q \in S$ to subsets induces
	a bijection between	the strongly connected components (resp.~upper sets) in~$\mathcal A_{|S}$ 
	and the strongly connected components (resp.~upper sets) $R$ in~$\mathcal A$ with~$R \subseteq S$.
\end{proposition}

\subsection{The Quotient Automaton}
\label{prelim:sec:aut:factor}

%!TEX root = ./../../thesis-trier.tex

\IN{factor (semi-)automaton} \label{def:factor_aut}
Let $\semiautomat{A}{Q}{\Sigma}{\delta}$. Suppose $\sim$ is a $\TS{\mathcal A}$-congruence. 
Equivalently, for every $p,q \in Q$ and $x \in \Sigma$,
we have
\begin{equation}\label{prelim:eqn:TS-congruence}
	p \sim q \Rightarrow \applytrans{\delta}{p}{x} \sim \applytrans{\delta}{q}{x}.
\end{equation}
The \emph{quotient automaton} is $\semiautomat{A / {\sim}}{Q / {\sim}}{\Sigma}{\delta_{\sim}}$ 
with the well-defined transition function $\applytrans{\delta_{\sim}}{[q]_{\sim}}{x} = [\applytrans{\delta}{q}{x}]_{\sim}$
for all $x \in \Sigma$ and~$q \in Q$.

The next two statements connect $\TS{\mathcal A}$-congruences, quotient automata, the kernel relation, and morphisms
in the following way: $\TS{\mathcal A}$-congruences
correspond to surjective morphisms onto quotient automata
via the canonical morphism (\cref{prelim:prop:congruence_induces_morphism}), and surjective morphisms
correspond to $\TS{\mathcal A}$-congruences via the kernel relation (\cref{prelim:thm:homtheorem}).

\IN{canonical morphism}
In view of~\cref{prelim:prop:congruence_induces_morphism}, in this context we also call $\pi_{\sim}$
the \emph{canonical morphism (for~$\sim$)} 
from $\mathcal A$ onto $\mathcal A / {\sim}$. 

\begin{theorem}
	\label{prelim:prop:congruence_induces_morphism}
	Let $\semiautomat{A}{Q}{\Sigma}{\delta}$
	with $\TS{\mathcal A}$-congruence $\sim$ on $\mathcal A$.
    Then the canonical projection map $\pi_{\sim} \colon Q \to Q / {\sim}$ from \cref{def:canprojmap}
    is a surjective morphism whose kernel relation equals~$\sim$.
\end{theorem}
\begin{proof}
	Clearly, $\pi_{\sim}$ is surjective and its kernel relation equals $\sim$.
	It is a morphism since by definition of $\mathcal A / {\sim}$
	we have $\delta_{\sim}(\pi_{\sim}(q), x) = \pi_{\sim}(\delta(q, x))$
	for $q \in Q$ and $x \in \Sigma$.
\end{proof}

\begin{theorem} 
	\label{prelim:thm:homtheorem}
	Let $\varphi \colon Q \to S$ be a morphism
	between the automata $\semiautomat{A}{Q}{\Sigma}{\delta}$
	and $\semiautomat{B}{S}{\Sigma}{\mu}$.
	Then the following is true:
	\begin{enumerate} 
	\item \label{prelim:thm:homtheorem:1} $\applyfunc{Q}{\varphi}$ is an upper set in $\mathcal B$,
	
	\item \label{prelim:thm:homtheorem:2} the kernel relation $\sim_{\varphi}$
	is a $\TS{\mathcal A}$-congruence,
	
	\item \label{prelim:thm:homtheorem:3} $[q]_{\sim_{\varphi}} \mapsto \applyfunc{q}{\varphi}$, $q \in Q$,
	is a well-defined isomorphism between
	$\mathcal A / {\sim_{\varphi}}$ and~$\mathcal B_{|\applyfunc{Q}{\varphi}}$.
	\end{enumerate}
\end{theorem}
\begin{proof}
\ref{prelim:thm:homtheorem:1}: This follows from \crefitem{prelim:automata:morpisms:reachability:prop:images}{prelim:automata:morpisms:reachability:prop:images:preservation}.

\ref{prelim:thm:homtheorem:2}: Suppose $p \sim_{\varphi} q$ for $p,q \in Q$.
 Then we have $\varphi(p) = \varphi(q)$.
 Since $\varphi$ is a morphism,
 we find $\varphi(\delta(p, x)) = \mu(\varphi(p), x) = \mu(\varphi(q), x) = \varphi(\delta(q, x))$,
 and so $\delta(p, x) \sim_{\varphi} \delta(q, x)$, i.e., $\sim_{\varphi}$
 is a $\TS{\mathcal A}$-congruence.
 
\ref{prelim:thm:homtheorem:3}: Well-definedness follows from the definition of the kernel relation.
 Now $\delta_{\sim_{\varphi}}([q]_{\sim_{\varphi}}, x) = [\delta(q, x)]_{\sim_{\varphi}}$ 
 is mapped to $\varphi(\delta(q, x))$
 and $[q]_{\sim_{\varphi}}$ is mapped to $\varphi(q)$.
 Since $\mu(\varphi(q), x) = \varphi(\delta(q, x))$,
 we obtain that the mapping is a morphism
 as we arrive at the same result when applying first a letter and then the mapping,
 or when first applying the mapping and then a letter.
 The mapping clearly maps onto $\mathcal B_{|\varphi(Q)}$.
\end{proof}

A similar statement as \cref{prelim:thm:homtheorem} in the theory of groups is often called the 
homomorphism theorem, for example, see~\cite[Sect.~1.2.5]{Kurzweil04}.
\Cref{prelim:thm:homtheorem} states that morphisms composed along paths in the following diagram
ending at the same point are equal.
\[
\begin{tikzcd}[row sep=1.3cm, column sep=2.5cm]
	Q \arrow{r}{q \mapsto [q]_{\sim_{\varphi}}}
	\arrow[swap,
	start anchor={[shift={(-3pt,3pt)}]south east}, % fine tune start of arrow
	end anchor={[shift={(2pt,-4pt)}]north west} % fine tune end of arrow
	]{dr}{\varphi}
	& Q / {\sim_{\varphi}} \arrow{d}{[q]_{\sim_{\varphi}} \mapsto \varphi(q)} \\
	& \varphi(Q) 
\end{tikzcd}
\]

\begin{corollary}	\label{prelim:cor:homtheorem}
	Let $\varphi \colon Q \to S$ be a morphism
	between the automata $\semiautomat{A}{Q}{\Sigma}{\delta}$
	and $\semiautomat{B}{S}{\Sigma}{\mu}$ with $Q \ne \emptyset$.
	If $\mathcal B$ is strongly connected, 
	then $\applyfunc{Q}{\varphi} = S$, i.e., $\varphi$ is surjective. So $\mathcal A / {\sim_{\varphi}} \cong \mathcal B$. 
	In particular, for $\mathcal A = \mathcal B$,
	by finiteness, $\varphi$ is an isomorphism.	
\end{corollary}

\begin{corollary}
  Suppose $\semiautomat{A}{Q}{\Sigma}{\delta}$ admits
  no non-trivial $\TS{\mathcal A}$-congruence.
  Then every morphism from $\mathcal A$ into another automaton
  is either injective or constant.  
\end{corollary}
\begin{proof}
  For a morphism into another automaton, its kernel relation
  is a $\TS{\mathcal A}$-congruence by \crefitem{prelim:thm:homtheorem}{prelim:thm:homtheorem:2}.
  So it either equals the universal relation, and so the morphism
  is a constant, or the equality relation, and so the morphism is injective.  
\end{proof}

Next, we show a condition to obtain a morphism from two morphisms
from the same automaton. The proof is straightforward and omitted.

\begin{theorem}
	\label{prelim:thm:homsfactorthrough}
	Let $\varphi \colon Q \to S$
	and $\psi \colon Q \to T$
	be morphisms
	from $\semiautomat{A}{Q}{\Sigma}{\delta}$
	to $\semiautomat{B}{S}{\Sigma}{\mu}$
	and $\semiautomat{C}{T}{\Sigma}{\xi}$, respectively.
	Suppose $\sim_{\varphi}$ refines $\sim_{\psi}$.
	Then $\varphi(q) \mapsto \psi(q)$
	is a well-defined surjective morphisms
	from $\mathcal B_{|\varphi(Q)}$
	onto $\mathcal C_{|\psi(Q)}$.
\end{theorem}

By \crefitem{prelim:thm:homtheorem}{prelim:thm:homtheorem:3},
$\mathcal A / {\sim_{\varphi}} \cong \mathcal B_{|\varphi(Q)}$
and $\mathcal A / {\sim_{\psi}} \cong \mathcal C_{|\psi(Q)}$
in \cref{prelim:thm:homsfactorthrough}. 
In the same way as for \cref{prelim:thm:homtheorem}, the situation from \cref{prelim:thm:homsfactorthrough} is illustrated in the following diagram.
\[
\begin{tikzcd}[column sep=1.3cm, row sep=1.3cm]
	& Q \arrow[dl, swap, "\varphi"] \arrow[dr, "\psi"] & \\
	S \arrow{rr}{\varphi(q) \mapsto \varphi(q)} & & T
\end{tikzcd}
\]

\begin{corollary}
 Let $\semiautomat{A}{Q}{\Sigma}{\delta}$ and $\semiautomat{B}{S}{\Sigma}{\mu}$
 with $\sim_{\mathcal A}$ and $\sim_{\mathcal B}$
 being $\TS{\mathcal A}$- and $\TS{\mathcal B}$-congruences, respectively.
 If $\sim_{\mathcal A}$ refines $\sim_{\mathcal B}$,
 then $[q]_{\sim_{\mathcal A}} \mapsto [q]_{\sim_{\mathcal B}}$
 is a morphism from $\mathcal A / {\sim_{\mathcal A}}$
 into $\mathcal B  / {\sim_{\mathcal B}}$.	
\end{corollary}

\subsection{The Rees Quotient Automaton}
\label{prelim:sec:rees} 

%!TEX root = ./../../springer-mono-format/book.tex
%!TEX root = ./../../thesis-trier.tex

Here, we introduce a special kind of quotient automata
whose construction resembles a similar one from the theory of semigroups
introduced by Rees~\cite{Rees40}.

\IN{Rees congruence}
\IN{Rees quotient automaton|quotient automaton}
Let $\semiautomat{A}{Q}{\Sigma}{\delta}$ and $S \subseteq Q$ be an upper set. 
The \emph{Rees congruence} is the  $\TS{\mathcal A}$-congruence on $Q$
given, for $p, q \in Q$, by \label{def:Rees_congruence}
\[ 
p \sim_S q \Leftrightarrow p,q \in S \lor (p = q).
\]
We call $\mathcal A / {\sim_S}$
the \emph{Rees quotient (semi-)automaton}
w.r.t.~the upper set $S \subseteq Q$.  For $S \ne \emptyset$, intuitively, 
$\mathcal A / {\sim_S}$
results from $\mathcal A$ by collapsing the states in $S$
to a single trap state in~$\mathcal A / {\sim_S}$.

\begin{remark}\label{prelim:rem:Rees}
	For $S \ne \emptyset$, unraveling the definition, $\mathcal A / {\sim_S}$ 
	has the state set $\{ \{p\} \mid p \in Q \setminus S \} \cup \{ S \}$,
	the transition function fulfills $\applytrans{\delta_{\sim_S}}{S}{x} = S$ and, 
	for all $p \notin S$ and all $x \in \Sigma$, 
	\begin{equation}\label{prelim:eqn:reestrans}
		\applytrans{\delta_{\sim_S}}{\{p\}}{x} = \begin{cases}
			\{ \applytrans{\delta}{p}{x} \} & \text{if } \applytrans{\delta}{p}{x} \notin S, \\ 
			S                               & \text{otherwise}.
		\end{cases}
	\end{equation}	
    So, $\mathcal A / {\sim_S}$ has $|Q| - |S| + 1$ states.
    By \cref{prelim:general:order:subset-in-equiv}, 
    as a strongly connected component
    is either contained in~$S$ or disjoint from~$S$, for $p,q \in Q$ and $x \in \Sigma$ with $p \eqreach q$
    we have
    \begin{equation}\label{prelim:eqn:reestransS}
    p = \delta(q, x) \Leftrightarrow 
    \begin{cases}
    p = \delta_{|S}(q, x) & \text{if } p \in S, \\
    \pi_{\sim_S}(p) = \delta_{\sim_S}(\pi_{\sim_S}(q), x) & \text{if } p \notin S.
    \end{cases}
    \end{equation}  	
	For $T \subseteq Q$,
    we have 
    \begin{equation}\label{prelim:eqn:Reesproj}
    \pi_{\sim_S}(T) = 
    \begin{cases}
    \{ S \} \cup \{ \{q\} \mid q\in T \setminus S \} & \text{if } S \cap T \ne \emptyset, \\
    \{ \{q\} \mid q \in T \}                         & \text{if } S \cap T = \emptyset.    	
    \end{cases}
\end{equation}
\end{remark}

As stated next, the projection maps induce bijections between 
certain strongly connected components and upper sets
in $\mathcal A$ and $\mathcal A / {\sim_S}$

\begin{proposition}\label{prelim:automata:Rees:prop:correspondence}
	Let $\semiautomat{A}{Q}{\Sigma}{\delta}$ and let $S\subseteq Q$ be an upper set.
	
	\begin{enumerate}
		\item Let $P$ be a strongly connected component of $\mathcal A/{\sim_S}$,
		distinct from $\{S\}$.
		Then $\pi_{\sim_S}^{-1}(P)$ is a strongly connected component of
		$\mathcal A$ contained in $Q\setminus S$, and
		\[
		\pi_{\sim_S}\bigl(\pi_{\sim_S}^{-1}(P)\bigr)=P.
		\]
		
		Conversely, let $R\subseteq Q\setminus S$ be a strongly connected component
		of $\mathcal A$.
		Then $\pi_{\sim_S}(R)$ is a strongly connected component of
		$\mathcal A/{\sim_S}$ distinct from $\{S\}$, and
		\[
		\pi_{\sim_S}^{-1}\bigl(\pi_{\sim_S}(R)\bigr)=R.
		\]
		
		\item Let $P$ be an upper set of $\mathcal A/{\sim_S}$ containing $S$.
		Then $\pi_{\sim_S}^{-1}(P)$ is an upper set of $\mathcal A$
		containing $S$, and
		\[
		\pi_{\sim_S}\bigl(\pi_{\sim_S}^{-1}(P)\bigr)=P.
		\]
		
		Conversely, let $R$ be an upper set of $\mathcal A$ containing $S$.
		Then $\pi_{\sim_S}(R)$ is an upper set of
		$\mathcal A/{\sim_S}$ containing $S$, and
		\[
		\pi_{\sim_S}^{-1}\bigl(\pi_{\sim_S}(R)\bigr)=R.
		\]
	\end{enumerate}
\end{proposition}
\begin{proof}
	This follows from \cref{prelim:rem:Rees},
	in particular \cref{prelim:eqn:reestrans}, 
	which yield that for $p,q \in Q \setminus S$ and $u \in \Sigma^*$
	we have $\delta(p, u) = q$
	if and only if $\delta_{\sim_S}(\pi_{\sim_S}(p), u) = \pi_{\sim_S}(q)$.
\end{proof}

\begin{corollary}
	The number of strongly connected components in $\mathcal A /{\sim_S}$
	is precisely the number of strongly connected components not in $S$
	plus one accounting for $\{S\}$.
\end{corollary}

\begin{corollary}
	Let $\semiautomat{A}{Q}{\Sigma}{\delta}$ and $S \subseteq Q$ be a proper upper set.
	Then $Q \setminus S$ is a strongly connected component in $\mathcal A$
	if and only if
	$( Q / {\sim_S} ) \setminus \{S\}$ is a strongly connected component in~$\mathcal A / {\sim_S}$.	
\end{corollary}

\begin{corollary}  \label{prelim:fact:factor-strongly-trap}
	For $S \ne \emptyset$ and $S\ne Q$,
	$\mathcal A / {\sim_S}$ is strongly trap-connected 
	if and only if $Q\setminus S$ is strongly connected
	and $\delta(Q \setminus S, x) \cap S \ne \emptyset$ for some $x \in \Sigma$.
\end{corollary}
\begin{proof}
	By \cref{prelim:automata:Rees:prop:correspondence},
	$Q \setminus S$ is strongly connected
	if and only if $\pi_{\sim_S}(Q \setminus S)$
	and $\pi_{\sim_S}(S)$ are the only strongly connected components in $\mathcal A / {\sim_S}$.
	Since~$\pi_{\sim_S}(Q \setminus S)$
	is not a sink component,
	the claim follows from \cref{prelim:automata:prop:char-strongly-trap:1b} of \cref{prelim:automata:prop:char-strongly-trap}.		
\end{proof}

Next, we state isomorphisms between Rees quotient automata.

\begin{proposition} \label{prelim:autmoata:Rees:prop:iso-eqns}
	Let $\semiautomat{A}{Q}{\Sigma}{\delta}$ and $S, T \subseteq Q$ be upper sets.	
	Then we have:
	\begin{align}
		\mathcal A & \cong \mathcal A / {\sim_S} & \text{ for } |S| \le 1  
		\label{prelim:Rees:S-at-most-one} \\
		 ( \mathcal A_{|S} ) / {\sim_{S\cap T}}
		& \cong  ( \mathcal A / {\sim_{T}} )_{|\pi_{\sim_{T}}(S)} 
		\label{prelim:eqn:Rees:ScapT} \\
		(\mathcal A_{|S\cup T}) / {\sim_T} & \cong  ( \mathcal A / {\sim_{T}} )_{|\pi_{\sim_{T}}(S)}  & \text{ for } S \cap T \ne \emptyset \label{prelim:eqn:Rees:ScapTScupT} \\
%		( \mathcal A_{|S} ) / {\sim_T} & \cong ( \mathcal A / {\sim_T} )_{|\pi_{\sim_T}(S)} & \text{ for } T \subseteq S
%		\label{prelim:automata:Rees:eqn:ScupT-T} \\
%		(\mathcal A / {\sim_T})_{|\pi_{\sim_T}(S)} 
%		& \cong 
%		\mathcal A_{|S} & \text{ for } |S \cap T| \le 1 \\
		( \mathcal A / {\sim_S} ) / {\sim_{\pi_{\sim_S}(T)}} & \cong \mathcal A / {\sim_{S \cup T}} 
		\label{prelim:aut:Rees:eqn:quot-ScupT} \\
%		( \mathcal A / {\sim_S} ) / {\sim_{\pi_{\sim_S}(T)}} & \cong \mathcal A / {\sim_T} & \text{ for } S \subseteq T \label{prelim:Rees:doublequotient} \\	
		( \mathcal A / {\sim_S} ) / {\sim_{\pi_{\sim_S}(T)}} & \cong ( \mathcal A / {\sim_T} ) / {\sim_{\pi_{\sim_T}(S)}} \label{prelim:aut:Rees:eqn:quotients-commute}	
	\end{align} 
\end{proposition}
\begin{proof}
\noindent\underline{\cref{prelim:Rees:S-at-most-one}}: For $|S| \le 1$
the Rees congruence $\sim_S$ equals the equality relation
and thus $\mathcal A \cong \mathcal A / {\sim_S}$.

\noindent\underline{\cref{prelim:eqn:Rees:ScapT}}: The map $( \pi_{\sim_T} )_{|S} \colon S \to Q / {\sim_T}$
is a morphism from $\mathcal A_{|S}$ to $\mathcal A / {\sim_T}$
with kernel $\sim_{S\cap T}$.
Thus the stated isomorphism follows from \cref{prelim:thm:homtheorem}.\footnote{Alternatively, 
	by \cref{prelim:thm:homsfactorthrough}, the map
	$
	\pi_{\sim_{S\cap T}}(q)\mapsto \pi_{\sim_T}(q)
	$
	for $q \in S$ is a well-defined surjective morphism from
	$(\mathcal A_{|S})/{\sim_{S\cap T}}$
	onto
	$(\mathcal A/{\sim_T})_{|\pi_{\sim_T}(S)}$.
	Since
	\[
	\pi_{\sim_{S\cap T}}(p)=\pi_{\sim_{S\cap T}}(q)
	\iff
	\pi_{\sim_T}(p)=\pi_{\sim_T}(q)
	\]
	for all $p,q\in S$, this morphism is an isomorphism.
	Thus we have $( \mathcal A_{|S} ) / {\sim_{S\cap T}} \cong  ( \mathcal A / {\sim_{T}} )_{|\pi_{\sim_{T}}(S)}$.}

\noindent\underline{\cref{prelim:eqn:Rees:ScapTScupT}}: Set $S' = S \cup T$.
Then $S'$ is an upper set and, by \cref{prelim:eqn:Rees:ScapT},
we find
\[
(\mathcal A_{|S'}) / {\sim_{S' \cap T}}  \cong  ( \mathcal A / {\sim_{T}} )_{|\pi_{\sim_{T}}(S')}.
\]
Then the claim follows
since $S' \cap T = T$
and for $S \cap T \ne \emptyset$ we have 
$\pi_{\sim_T}(S) = \pi_{\sim_T}(S')$.

\noindent\underline{\cref{prelim:aut:Rees:eqn:quot-ScupT}}: The map 
$\concat{\pi_{\sim_S}}{\pi_{\sim_R}} \colon Q \to ( Q / {\sim_S} ) / {\sim_R}$ with $R = \pi_{\sim_S}(T)$
is a morphism with kernel $\sim_{S\cup T}$.
Thus the stated isomorphism follows from \cref{prelim:thm:homtheorem}.\footnote{Alternatively,
since $\sim_{S}$ refines $\sim_{S\cup T}$,
the map $\pi_{\sim_S}(p) \mapsto \pi_{\sim_{S\cup T}}(p)$ for $q \in Q$
is a well-defined surjective morphism
by \cref{prelim:thm:homsfactorthrough}.
Since it has kernel $\pi_{\sim_S}(T)$
the stated isomorphism follows from \cref{prelim:thm:homtheorem}.}

\noindent\underline{\cref{prelim:aut:Rees:eqn:quotients-commute}}: 
By \cref{prelim:aut:Rees:eqn:quot-ScupT} both automata involving two Rees quotients
are isomorphic to $\mathcal A / {\sim_{S\cup T}}$. Thus they are isomorphic
to each other.
\end{proof}

Note that \cref{prelim:eqn:Rees:ScapT}
also implies 
$( \mathcal A_{|S} ) / {\sim_{S\cap T}} \cong ( \mathcal A / {\sim_{S\cap T}} )_{|\pi_{\sim_{S\cap T}}(S)}$
by applying it to the upper sets $S$ and $S \cap T$ instead of~$S$ and~$T$.


For a morphism and an upper set, we obtain
a morphism between corresponding Rees quotient automata.

\begin{proposition} \label{prelim:prop:Rees-homsfactorthrough}
	Let $\varphi \colon Q \to S$ be a morphism
	from $\semiautomat{A}{Q}{\Sigma}{\delta}$
	to $\semiautomat{B}{S}{\Sigma}{\mu}$ and  $T \subseteq Q$ be an upper set.
	Then $\varphi(T)$ is an upper set %in~$\mathcal B$
	and 
	\[ [q]_{\sim_T} \mapsto [\varphi(q)]_{\sim_{\varphi(T)}}
	\]
	for $q \in Q$ is a well-defined surjective morphism from 
	$\mathcal A / {\sim_T}$ onto $( \mathcal B / {\sim_{\varphi(T)}} )_{|\pi_{\sim_{\varphi(T)}}(\varphi(Q))}$.
\end{proposition}
\begin{proof}
	For $p,q \in Q$, we have 
	\begin{align*}
		p \sim_T q & \Rightarrow p = q \lor \{p,q\} \subseteq T \\ 
		           & \Rightarrow \varphi(p) = \varphi(q) \lor \{\varphi(p), \varphi(q)\} \subseteq \varphi(T) \\ 
                   & \Rightarrow \varphi(p) \sim_{\varphi(T)} \varphi(q).
	\end{align*}
	Hence, \cref{prelim:thm:homsfactorthrough}
	gives the claim. In fact, we have $\varphi([q]_{\sim_T}) = [\varphi(q)]_{\sim_{\varphi(T)}}$
	and so $\varphi$ induces this morphism on the equivalence classes of $\sim_T$.
\end{proof}

Lastly, using the Rees quotient automaton construction
we can derive a necessary condition, originally due to Thierrin~\cite[Prop.~8]{DBLP:journals/kybernetika/Thierrin70},
for primitivity of automata (see also \cite[Cor.~3]{DBLP:journals/actaC/Babcsanyi07}
or, in a more algebraic context, \cite[Prop.~5.1]{Steinberg2010}).

\begin{proposition}[{Thierrin~\cite[Prop.~8]{DBLP:journals/kybernetika/Thierrin70}}]
	\label{prelim:cor:simpleaut}
	Let $\semiautomat{A}{Q}{\Sigma}{\delta}$ be a primivite automaton with $|Q| \ge 3$. 
    Then $\mathcal A$ is strongly connected
    or strongly trap-connected with trap state $t \in Q$
    such that 	$\delta(p, a) = t$ and $\delta(q, a) \ne t$
    for some $p,q \in Q$ and $a \in \Sigma$.
\end{proposition}
\begin{proof}	
	There exists no proper upper set $S \subseteq Q$ of size at least two,
	as otherwise~$\sim_S$ would be a non-trivial $\TS{\mathcal A}$-congruence.
	Hence, by \cref{prelim:fact:no-proper-upper-sets-size-two},
	$\mathcal A$ is either strongly connected or strongly trap-connected
	with a trap state $t \in Q$.
	In the latter case, since $Q \setminus \{t\}$ and $\{t\}$ form a partition
	that does not induce a $\TS{\mathcal A}$-congruence,
	there exist $p,q \in Q$ and $a \in \Sigma$
	such that $\delta(p, a) = t$ and $\delta(q, a) \ne t$.
\end{proof}

The following corollary has been noted in~\cite[Sect.~5.1]{RystsovSzykula2023}
and~\cite[Prop.~3.1]{Volkov24}.

\begin{corollary}
	A primitive automaton $\mathcal A$ with at least three states is either strongly connected or synchronizing.	
\end{corollary}
\begin{proof}
	This follows from \cref{prelim:cor:simpleaut} since a strongly trap-connected automaton
	is synchronizing.
\end{proof}

\begin{remark}
	\Cref{prelim:cor:simpleaut} generalizes 
	\cref{sync_max:prop:prim_implies_transitive} (via the correspondence from \cref{prelim:remarks_trans_semigrp_semiautomata}),
	since a trap state in a permutation automaton
	is not reachable from any other state.
	Thus, by \cref{prelim:cor:simpleaut}, a permutation automata
	whose transformation monoid is a primitive group $G$
	is strongly connected and so~$G$ is transitive.	
\end{remark}


\subsection{The Product Automaton} 
\label{prelim:sec:productaut} 

%!TEX root = ./../../thesis-trier.tex

\todo[inline]{hier kann ich nicht automat-kommando nehmen in Definition von Produktautomaten.}

\IN{product (semi-)automaton}
\label{def:product_aut} 
Given $\semiautomat{A_i}{Q_i}{\Sigma}{\delta_i}$
for $i \in \{1,\ldots,m\}$,  
their \emph{product (semi\mbox{-)}automaton}~\cite{HopUll79,DBLP:journals/ibmrd/RabinS59} 
is $\mathcal A_1 \times \ldots \times \mathcal A_m = (Q_1 \times \ldots \times Q_m,\Sigma, \delta_{\mathcal A_1 \times \ldots \times \mathcal A_m})$ with 
\[ 
\applytrans{\delta}{(q_1, \ldots, q_m)}{x}
= (\applytrans{\delta_1}{q_1}{x}, \ldots, \applytrans{\delta_m}{q_m}{x})
\]
for $q_1 \in Q_1, \ldots, q_m \in Q_m$ and $x \in \Sigma$.
If $\mathcal A_1 = \cdots = \mathcal A_m$, then
we also write $\mathcal A_1^m$
for~$\mathcal A_1 \times \cdots \times \mathcal A_m$.
For the set of synchronizing words we have
\begin{equation}\label{prelim:eqn:product_aut_sync_words}
	\Syn(\mathcal A_1 \times \cdots \times \mathcal A_m)
	= \bigcap_{i=1}^m \Syn(\mathcal A_i).
\end{equation}

For upper sets $S_1 \subseteq Q_1, \ldots, S_m \subseteq Q_m$,
the set $S = S_1 \times \cdots \times S_m$ is an upper set
and since $(q_1, \ldots, q_m) \in S$ if and only if $q_1 \in S_1, \ldots, q_m \in S_m$
we have
\begin{align}
	(\mathcal A_1 \times \cdots \times \mathcal A_m)_{|S} & \cong 
	( \mathcal A_1 )_{|S_1} \times \cdots \times ( \mathcal A_m )_{|S_m} \label{prelim:automata:prod:eqn:sub} \\
	(\mathcal A_1 \times \cdots \times \mathcal A_m) / {\sim_S} 
	& \cong ( \mathcal A_1 / {\sim_{S_1}} ) \times \cdots \times ( \mathcal A_m / {\sim_{S_m}} ) \label{prelim:automata:prod:eqn:quot}
\end{align}

\subsection{The Power Automaton} 
\label{prelim:sec:aut:power}

%!TEX root = ./../../springer-mono-format/book.tex
%!TEX root = ./../../thesis-trier.tex

\label{def:power_aut} \IN{power (semi-)automaton}
Let $\semiautomat{A}{Q}{\Sigma}{\delta}$. The \emph{power (semi-)automaton} 
is $\semiautomat{P({\mathcal A})}{\mathcal P(Q)}{\Sigma}{\delta}$
where $\delta \colon \mathcal P(Q)\times \Sigma \to \Sigma$ is the extension of
$\delta \colon Q \times \Sigma \to Q$ to subsets $\applytrans{\delta}{S}{x} = \{ \applytrans{\delta}{q}{x} \mid q \in Q \}$ 
for $S \subseteq Q$ and $x \in \Sigma$ (see also Section~\ref{intro:sec:basic}).
\Cref{fig:thm_not_usable} shows the power automaton
of a four-state automaton. Also, part of the power 
automaton of the \v{C}ern\'y automaton $\mathcal C_4$
is depicted in \cref{sync_max:fig:C4-poweraut}.

In this section, we present properties of powers of a single word in the power automaton. These can also be interpreted as properties of transformations on subsets of a finite set; however, since our focus is on automata, we express them in automata-theoretic terms.

\todo[inline]{bild kompakter}

\begin{figure}[htb]
\begin{minipage}{0.25\textwidth}
 \scalebox{1}{\begin{tikzpicture}[shorten >=1pt,->,node distance=2.3cm]
  %\tikzstyle{vertex}=[circle,color=red!60,fill=black!10,minimum size=25pt,inner sep=1pt,very thick]
  %\tikzstyle{vertex}=[circle,draw,minimum size=25pt,inner sep=1pt]
  
  \node[state] (0) {$0$};
  \node[state, right of=0] (1) {$1$};
  \node[state, below of=1] (2) {$2$};
  \node[state, left of=2]  (3) {$3$};
  
  \path[->] (0) edge [above] node {$a,b$} (1)
            (1) edge [bend left, right] node {$a,b$} (2)
            (2) edge [below] node {$b$} (3)
            (3) edge [left]  node {$b$} (0)
            (2) edge [bend left, left] node {$a$} (1)
            (3) edge [loop below] node {$a$} (3);
\end{tikzpicture}}%
\end{minipage}%\hspace*{0.25cm}
\begin{minipage}{0.75\textwidth}
\scalebox{1}{\begin{tikzpicture}[shorten >=1pt,->,node distance=1cm and 0.9cm]
  \tikzstyle{vertex}=[]
  
  % \node at (0,-1) (s) {start};  
   
  \node[vertex] (0123)               {$\{0,1,2,3\}$};
  \node[vertex, right = of 0123, xshift=-0.5cm] (123) {$\{1,2,3\}$};
  \node[vertex, below = of 123]  (023) {$\{0,2,3\}$};
  \node[vertex, below = of 023]   (013) {$\{0,1,3\}$};
  \node[vertex, below left = of 013,xshift=0.7cm]  (012) {$\{0,1,2\}$};
  
  \node[vertex, right = of 023] (13) {$\{1,3\}$};
  \node[vertex, above = of 13]  (02) {$\{0,2\}$};
  \node[vertex, below = of 13]  (23) {$\{2,3\}$};
  \node[vertex, right = of 23]  (03) {$\{0,3\}$};
  \node[vertex, below = of 03]  (01) {$\{0,1\}$};
  \node[vertex, left = of 01]  (12) {$\{1,2\}$};
  
  \node[vertex,  right = of 02, yshift=0.2cm] (0) {$\{0\}$};
  \node[vertex,  right = of 0]  (1) {$\{1\}$};
  \node[vertex,  below = of 1]  (2) {$\{2\}$};
  \node[vertex,  left = of 2]   (3) {$\{3\}$};
  
  
  \path (123) edge [loop above] node {$a$} (123);
  \path (012) edge [above] node {$a$} (12);
  \path (013) edge [bend left=45,left] node {$a$} (123);
  
  %\path[->] (s) edge (0123);
  
    \path[->] (0) edge [below] node {$a,b$} (1)
            (1) edge [bend left, right] node {$a,b$} (2)
            (2) edge [below] node {$b$} (3)
            (3) edge [left]  node {$b$} (0)
            (2) edge [bend left, right] node {$a$} (1)
            (3) edge [loop below] node {$a$} (3);
            
  
  \path[->] (0123) edge [above] node {$a$} (123);
  \path[->] (0123) edge [loop above] node {$b$} (0123);
  \path[->] (123)  edge [right] node {$b$} (023); 
  \path[->] (023)  edge [right] node {$b$} (013);
  \path[->] (013)  edge [above,pos=.6] node {$b$} (012);
  \path[->] (012)  edge [bend left,left] node {$b$} (123);
  
  \path[->] (023)  edge [above] node {$a$} (13);
  \path[->] (13)   edge [bend left,left] node {$b$} (02);
  \path[->] (02)   edge [bend left=40,right] node {$b$} (13);
  \path[->] (13)   edge [bend left,right] node {$a$} (23);
  \path[->] (23)   edge [bend left,left] node {$a$} (13);
  \path[->] (23)   edge [above] node {$b$} (03);
  \path[->] (03)   edge [right] node {$b$} (01);
  \path[->] (01)   edge [above] node {$a,b$} (12);
  \path[->] (12)   edge [left] node {$b$} (23);
  \path[->] (03)   edge [above] node {$a$} (13);
  \path[->] (12)   edge [loop below] node {$a$} (12);
  
  \path[->] (02) edge [bend left,above] node {$a$} (1);
  
\end{tikzpicture}}%
\end{minipage}%
 %\caption{Automaton from Example~\ref{ex:I_vs_Gamma1} where $I(\TS{\mathcal A})$ 
% 	is a proper subgraph of $\Gamma_1(\TS{\mathcal A})$.}
\caption[Example power automaton]{Left: An automaton. Right: Its power automaton (empty set omitted).}
 \label{fig:thm_not_usable}
\end{figure}

\begin{proposition} \label{prelim:fact:cycle_states} \label{prelim:automata:power:P}
	Let $\semiautomat{A}{Q}{\Sigma}{\delta}$ with $n = |Q|$ and $u \in \Sigma^*$. 
	Then for $i \ge 0$ we have
	$\delta(Q, u^{i+1}) \subseteq \delta(Q, u^i)$
	and for $P = \delta(Q, u^{n-1})$
	the following is true:
	\begin{enumerate}
	\item $P = \delta(Q, u^{n - |P|})$,
		
	\item $\delta(P, u) = P$ ($u$ permutes~$P$) where~$P$
	is the largest set with this property, and $P = \delta(Q, u^i)$ for all $i \ge n - |P|$,
		
	\item $P \subseteq \delta(Q, u^i)$ for all $i \ge 0$.
	\end{enumerate}
\end{proposition}
\begin{proof}
	For the first claim, since $\delta(Q, u) \subseteq Q$
	and $\delta(Q, u^{i+1}) \subseteq \delta(Q, u^i)$
	implies 
	\[ \delta(\delta(Q, u^{i+1}), u) \subseteq \delta(\delta(Q, u^i), u),
	\]
	by induction we find that $\delta(Q, u^{i+1}) \subseteq \delta(Q, u^i)$
	for all $i \ge 0$.
	
	Next, by finiteness, there exists a minimal $j \ge 0$
	such that $\delta(Q, u^{j+1}) = \delta(Q, u^j)$.
	Then we have $\delta(Q, u^k) = \delta(Q, u^j)$ for all $k \ge j$
	and we get the following chain of proper inclusions
	\begin{equation}\label{prelim:automata:power:eqn:seq}
		\delta(Q, u^j) \subsetneq \delta(Q, u^{j-1}) \subsetneq\cdots \subsetneq \delta(Q, u) \subsetneq Q.
	\end{equation}
	This yields $\delta(Q, u^j) \subseteq \delta(Q, u^i)$ for every $i \ge 0$.
	Since $|\delta(Q, w)| \ge 1$ for all $w \in \Sigma^*$,
	the above chain of proper inclusions has length at most $n - |\delta(Q, u^j)|$, i.e., $j  \le n - |\delta(Q, u^j)|$.
	Hence we find
	\[
	P = \delta(Q, u^{n-1}) = \delta(Q, u^{n - |\delta(Q, u^j)|}) = \delta(Q, u^j)
	= \delta(Q, u^{n}) = \delta(P, u),
	\]
	and so $P = \delta(Q, u^{n - |P|})$,
	$P \subseteq \delta(Q, u^i)$ for all $i \ge 0$,
	and $P$ is permuted by~$u$. This also gives
	$P = \delta(Q, u^i)$ for $i \ge n - |P|$
	by induction.	
	Furthermore, if $\delta(S, u) = S$ for $S \subseteq Q$,
	then $S = \delta(S, u^i) \subseteq \delta(Q, u^{i})$ for all $i \ge 0$.
	In particular, we obtain $S \subseteq \delta(Q, u^{n-1}) = P$
	and so $P$ is the largest set permuted by~$u$. This shows all claims.
\end{proof}

The set~$P$ from \cref{prelim:fact:cycle_states}
is precisely the set of cycle states described in \cref{prelim:rem:intuition-functional-graph}.
By finiteness, after applying a word at most $n - 1$ times
to any state we end up at a cycle state. 

Note that since $\delta(Q, u) \subseteq Q$ for $u \in \Sigma^*$,
we have $\delta(Q, uv) \subseteq \delta(Q, v)$ for each $v \in \Sigma^*$.

\begin{corollary}  \label{prelim:automata:cor:power:P}
	Let $\semiautomat{A}{Q}{\Sigma}{\delta}$ with $n = |Q|$ and $u \in \Sigma^*$. 
	Then we have $\delta(Q, u^i) = \delta(Q, u^{\min\{n-1,i\}})$ for $i \ge 0$.
\end{corollary}
\begin{proof}
	 By \cref{prelim:automata:power:P}, the word~$u$
	 permutes the set $\delta(Q, u^{n-1})$.
	 Thus, for $i \ge n - 1$ we find $\delta(Q, u^i) = \delta(Q, u^{n-1})$, which gives the claim.
\end{proof}

In~\cite[Lem.~1~\&~2]{DBLP:journals/iandc/PribavkinaR11} it was also proved that $P$
is the largest set permuted by~$u$ in \cref{prelim:fact:cycle_states},
together with the fact that we have to apply~$u$ at most $n - |P|$
times to~$Q$ to obtain~$P$.

Note that for the properties collected in \cref{prelim:automata:power:P} to hold
it is essential that we start from the whole state set before applying powers of~$u$.

\begin{proposition} \label{prelim:word_in_PA} \label{prelim:automata:power:prop:S}  
Let $\semiautomat{A}{Q}{\Sigma}{\delta}$, $n = |Q|$, $u \in \Sigma^*$, $S \subseteq Q$, and $i \ge 0$.

\begin{enumerate}
\item \label{prelim:automata:power:prop:S:1} $\delta(S, u^i) = \delta(S, u^j)$ for some $0 \le j < 2^n - 1$.

\item \label{prelim:automata:power:prop:S:2} $|\delta(S, u^i)| = |\delta(S, u^j)|$ for some $0 \le j \le n - |\delta(S, u^i)|$. 

\item \label{prelim:automata:power:prop:S:3} $|\delta(S, u^i)| = |\delta(S, u^{\min\{n-1,i\}})|$. 
\end{enumerate}
\end{proposition}
\begin{proof}
\ref{prelim:automata:power:prop:S:1}: We can assume $S \ne \emptyset$.
 Let $j \ge 0$ be a minimal number
 such that $\delta(S, u^i) = \delta(S, u^j)$.
 Assume $j \ge 2^n - 1$.
 We have $\delta(S, u^k) \ne \emptyset$ for all $k \ge 0$. 
 Since there are at most $2^n - 1$ non-empty subsets, 
 and the sequence $\delta(S, u^0), \delta(S, u), \ldots, \delta(S, u^{2^n-1})$
 has length $2^n$, there exist $0 \le k < k' \le 2^n - 1$
 such that $\delta(S, u^{k}) = \delta(S, u^{k'})$.
 Then $\delta(S, u^{j - (k' - k)}) = \delta(S, u^j)$,
 contradicting the minimal choice of~$j$.
 So, we have $j < 2^n - 1$.

\ref{prelim:automata:power:prop:S:2}: We can assume $i > n - |\delta(S, u^i)|$.
 For the minimal $j \ge 0$ 
 with $\delta(Q, u^{j+1}) = \delta(Q, u^j)$, 
 we have a proper chain of inclusions as in \cref{prelim:automata:power:eqn:seq}.
 Thus, for $k \in \{0,1,\ldots,j\}$,
 we get $|\delta(Q, u^k)| \le n - k$.
 Since $\delta(S, u^j) \subseteq \delta(Q, u^j)$,
 we find $|\delta(S, u^j)| \le n - j$,
 and so $j \le n - |\delta(S, u^j)|$.
 Set $P = \delta(Q, u^{n-1})$.
 With \cref{prelim:automata:power:P} we get
 $\delta(Q, u^j) = P$ and that $u$ permutes~$P$,
 and so $u$ acts injectively on subsets of~$P$.
 Thus, from $\delta(S, u^j) \subseteq P$
 we obtain $|\delta(S, u^k)| = |\delta(S, u^j)|$
 for all $k \ge j$.
 In particular, since $i > j$,
 we find $|\delta(S, u^i)| = |\delta(S, u^j)|$.
  
 \ref{prelim:automata:power:prop:S:3}: We can assume $i \ge n$.
 	By Item~\ref{prelim:automata:power:prop:S:2} 
 	we have $|\delta(S, u^i)| = |\delta(S, u^j)|$
 	for some $0 \le j < n$.
 	With \cref{eqn:subset-mapped-to-smaller-set} 
 	we find
 	\begin{enumerate}
 	\item $
 	|\delta(S, u^{n-1})| = |\delta(\delta(S, u^j), u^{n - 1 - j})| \le |\delta(S, u^j)|
 	$
 	and
 	
 	\item $|\delta(S, u^i)| = |\delta(\delta(S, u^{n - 1}), u^{i - n + 1})| \le |\delta(S, u^{n - 1})|.
 	$
 	\end{enumerate}
 	Thus we obtain $|\delta(S, u^i)| = |\delta(S, u^{n-1})|$.
\end{proof}

\begin{corollary} 	\label{prelim:fact:unary_sync_word}
 Let $\semiautomat{A}{Q}{\Sigma}{\delta}$, $u \in \Sigma^*$, and $S \subseteq Q$.
 Then $|\delta(S, u^i)| = 1$ for some $i \ge 0$
 if and only if $|\delta(S, u^{n - 1})| = 1$ for $n = |Q|$.
\end{corollary}
\begin{proof}
  First, assume $|\delta(S, u^i)| = 1$ for some $i \ge 0$.
  Since $|\delta(S, u^{i})| = |\delta(S, u^{\min\{n-1, i\}})|$
  by \crefitem{prelim:automata:power:prop:S}{prelim:automata:power:prop:S:3} 
  and $|\delta(S, u^{n - 1})| \le |\delta(S, u^{\min\{n-1, i\}})|$, 
  we find $|\delta(S, u^{n - 1})| = 1$.
  The converse implication is obvious.
\end{proof}

\begin{remark}
To obtain \crefitem{prelim:automata:power:prop:S}{prelim:automata:power:prop:S:1},
we have basically used the fact that in any $n$-state automaton 
we enter a cycle when we read a word more than $n-1$ times.
Now, as $u$ permutes $P$ in \cref{prelim:fact:cycle_states},
the smallest $i \ge 0$ such that $\delta(q, u^i) = q$
for all $q \in P$
is asymptotically equal to $\exp((1+o(1))\sqrt{(|P|\ln{|P|})}$
by a result from Landau~\cite{landau03}.
This can be used to improve the bound from \cref{prelim:word_in_PA}
asymptotically. However, as 
the bound given is sufficient for our applications,
we do not bother to state a more precise bound.
\end{remark}

\subsection{The Pair Automaton}
\label{sec:pair_aut}

%!TEX root = ./../../thesis-trier.tex

\IN{pair (semi)automaton}
\label{def:pair_aut}
Let $\semiautomat{A}{Q}{\Sigma}{\delta}$ and $S = \{ \{ p, q\} \mid p,q \in Q\}$.
Then~$S$ is an upper set in $\mathcal P(\mathcal A)$.
We denote $\mathcal P(\mathcal A)_{|S}$
also by $\semiautomat{A^{\le 2}}{S}{\Sigma}{\delta^{\le 2}}$.
We have $\applytrans{\delta^{\le 2}}{\{p,q\}}{x} = \applytrans{\delta}{\{p,q\}}{x}$
for $x \in \Sigma$.
The set $F = \{ \{q\} \mid q \in Q \}$ is an upper set in $\mathcal A^{\le 2}$
and the Rees quotient automaton $\mathcal A^{\le 2} / {\sim_F}$ is called the \emph{pair (semi-)automaton}. 
The state set of the pair automaton is
$\{ \{ \{ p,q\}\} \mid p \ne q \text{ and } p,q \in Q\} \cup \{F\}$.
For ease of notation, for distinct $p,q \in Q$ we identify a state $\{ \{p,q\} \}$ with $\{p,q\}$.
The pair automaton has $\binom{n}{2} + 1$ states.

The map $\varphi \colon Q \times Q \to \{\{p,q\} \mid p,q \in Q \}$
given by $\varphi(p,q) = \{p,q\}$ is a surjective morphism
of $\mathcal A \times \mathcal A$ onto $\mathcal A^{\le 2}$.
Thus, by \cref{prelim:prop:congruence_induces_morphism},
$[(p,q)]_{\sim_{\varphi}} \mapsto \{p,q\}$
is an isomorphism from 
$( \mathcal A \times \mathcal A ) / {\sim_{\varphi}}$ to $\mathcal A^{\le 2}$.
Note that $(p,q) \sim_{\varphi} (s,t)$ if and only if $\{p,q\} = \{s,t\}$
and so $\sim_{\varphi}$ is the $\TM{\mathcal A \times \mathcal A}$-congruence
that precisely identifies $(p,q)$ with $(q,p)$ for $p,q \in Q$.

\section{Sequences of Upper Sets and Maximal States}
\label{prelim:sec:aut:sequppersets}

%!TEX root = ./../../thesis-trier.tex

In this section, we investigate sequences of upper sets
as constructed in \cref{prelim:general:order:sec:seq-upper-sets} for general quasiorders
in the special case of automata and the reachability order. More specifically, we relate these sequences
to subautomata and Rees quotient automata as introduced previously in \cref{prelim:sec:aut:sub,prelim:sec:rees}.

\begin{proposition} \label{prelim:prop:composition-series}
	Let $\semiautomat{A}{Q}{\Sigma}{\delta}$.
	There exists a sequence of upper sets 
	\[
	\emptyset = S_0 \subsetneq S_1 \subsetneq \cdots \subsetneq S_m = Q
	\]
	such that the sets $S_i \setminus S_{i-1}$
	for $i \in \{1,\ldots,m\}$
	are all the strongly connected components in~$\mathcal A$,
	$\mathcal A_{|S_1}$ is strongly connected
	and, for $i \in \{2,\ldots,m\}$,
	$( \mathcal A_{|S_i} ) / {\sim_{S_{i-1}}}$
	is strongly trap-connected
	or $S_i \setminus S_{i-1}$ is a sink component in $( \mathcal A_{|S_i} ) / {\sim_{S_{i-1}}}$ (so that $S_i \setminus S_{i-1}$ 
	and $\{S_{i-1}\}$
	are the only sink components in this case).
\end{proposition}
\begin{proof}
	Consider a sequence $\emptyset = S_0 \subsetneq S_1 \subsetneq \cdots \subsetneq S_m = Q$
	as in \crefitem{prelim:general:order:cor:seq-upper-sets}{prelim:general:order:cor:seq-upper-sets:1},
	where, according to the statement,
	the sets $S_i \setminus S_{i-1}$ for $i \in \{1,\ldots,m\}$
	are all the strongly connected components. Since $S_1$
	is a strongly connected component,~$\mathcal A_{S_1}$
	is strongly connected.
	Next, let $i \in \{2,\ldots,m\}$.
	Now, suppose $S_i \setminus S_{i-1}$ 
	is not a sink component.
	Since the state set of  $( \mathcal A_{|S_i} ) / {\sim_{S_{i-1}}}$
	is $S_i \setminus S_{i-1} \cup \{S_{i-1}\}$ with $S_{i-1} \ne \emptyset$,
	if~$S_i \setminus S_{i-1}$ is not a sink component,
	by \crefitem{prelim:automata:prop:char-strongly-trap}{prelim:automata:prop:char-strongly-trap:1b}
	the automaton is strongly trap-connected;
	otherwise $S_i \setminus S_{i-1}$ is a sink component
	and the claim is proved.	
\end{proof}

In a sequence as in \cref{prelim:prop:composition-series},
for $p \in S_i$ and $q \in S_j$ with $p \lereach q$ we have $i \le j$.

\begin{proposition}\label{prelim:cor:composition-series-maximal-states}
	Let $\semiautomat{A}{Q}{\Sigma}{\delta}$.
	There exists a sequence of upper sets 
	\[
	S_1 \subsetneq \cdots \subsetneq S_m = Q
	\]
	such that~$S_1$ equals the set of maximal states in $\mathcal A$, and for $i \in \{2,\ldots,m\}$,
	$( \mathcal A_{|S_i} )/ {\sim_{S_{i-1}}}$
	is strongly trap-connected 
	and $S_{i} \setminus S_{i-1}$
	are all the strongly connected components in $\mathcal A$ that are not sink components.
\end{proposition} 
\begin{proof}
	Consider a sequence $S_0 \subsetneq S_1 \subsetneq \cdots \subsetneq S_m = Q$
	as in \crefitem{prelim:general:order:cor:seq-upper-sets}{prelim:general:order:cor:seq-upper-sets:2}
	where~$S_0$ equals the set of maximal states
	and $S_{i} \setminus S_{i-1}$
	for $i \in \{1,\ldots,m\}$ are all the non-sink strongly connected components.
	Since the state set of  $( \mathcal A_{|S_i} ) / {\sim_{S_{i-1}}}$ for $i \in \{1,\ldots,m\}$ 
	is $S_i \setminus S_{i-1} \cup \{S_{i-1}\}$ with $S_{i-1} \ne \emptyset$
	and  $S_i \setminus S_{i-1}$ 
	is not a sink component,
	it is strongly trap-connected
	by \crefitem{prelim:automata:prop:char-strongly-trap}{prelim:automata:prop:char-strongly-trap:1b}.	
\end{proof}

\begin{remark}
 The strongly trap-connected quotient automata
 that arise from a series as in \cref{prelim:prop:composition-series} 
 are always the same up to isomorphism.
 However, the remaining quotient automata
 and the subautomaton $\mathcal A_{S_1}$
 might be different for different series, where they
 differ by an isolated, i.e., not reachable, trap state.
 However, the length of such a series is always the same.
 A series as in \cref{prelim:cor:composition-series-maximal-states}
 always has the same length for a given automaton
 and the resulting quotient automata $( \mathcal A_{|S_i} )/ {\sim_{S_{i-1}}}$
 for $i \in \{2,\ldots,m\}$
 are always the same up to isomorphism.
 However, $\mathcal A_{S_1} / {\sim_{\emptyset}}$
 might admit further subautomata.
 For synchronizing automata,
 series obtained according
 to \cref{prelim:prop:composition-series} 
 or \cref{prelim:cor:composition-series-maximal-states}
 have the same length, isomorphic quotient automata,
 and $\mathcal A_{|S_1}$ admits no subautomata. 
 Thus, in this case, we obtain
 an analog of the Jordan-Hölder theorem
 from algebra~\cite{lang02}.	
\end{remark}

\begin{corollary}\label{prelim:automata:cor:composition-series-maximal-states}
	Let $\semiautomat{A}{Q}{\Sigma}{\delta}$.
	Then $Q$ equals the set of maximal states
	or there exists a proper and non-empty upper set $S \subseteq Q$
	such that $\mathcal A / {\sim_S}$
	is strongly trap-connected.
\end{corollary}
\begin{proof}
	Consider a sequence $S_1 \subsetneq \cdots \subsetneq S_m = Q$ 
	as in \cref{prelim:cor:composition-series-maximal-states}.
	For $m = 1$, $Q$ equals the set of maximal states.
	Otherwise, the claim is fulfilled with 
	the upper set $S = S_{m-1}$.	
\end{proof}

\begin{proposition} \label{prelim:automata:prop:max-states}
   	Let $\semiautomat{A}{Q}{\Sigma}{\delta}$ and $S \subseteq Q$ be the set of maximal states.
   	Then we have the following.
   	
   	\begin{enumerate}
   	\item \label{prelim:automata:prop:max-states:1} For all $q \in Q$ there exists $u \in \Sigma^*$ such that $\delta(Q, u) \subseteq S$.
   			
   	\item \label{prelim:automata:prop:max-states:2} $[q]_{\sim_S} \lereach^{\mathcal A / {\sim_S}} S$ in $\mathcal A / {\sim_S}$ for all $q \in Q$.
   	\end{enumerate}
\end{proposition}
\begin{proof}
\ref{prelim:automata:prop:max-states:1} \& \ref{prelim:automata:prop:max-states:2}: 
By finiteness, for each $q \in Q$ there exists a maximal $p \in S$
such that $q \lereach p$. By the definition of the reachability order,
this yields $\delta(q, u) = p$ for some $u \in \Sigma^*$, which proves Item~\ref{prelim:automata:prop:max-states:1}.
Next for the canonical morphism
$\pi_{\sim_S} \colon Q \to Q / {\sim_S}$  from \cref{prelim:prop:congruence_induces_morphism}
we have $\pi_{\sim_S}(q) \lereach^{\mathcal A / {\sim_S}} \pi_{\sim_S}(p)$ by \cref{prelim:eqn:morph-reachability-order}
where $\pi_{\sim_S}(q) = [q]_{\eqreach}$
and $\pi_{\sim_S}(p) = S$ since $p \in S$.
\end{proof}

\begin{proposition} \label{prelim:prop:maximal-states-SCC}
	Let $\semiautomat{A}{Q}{\Sigma}{\delta}$. The set~$S$ of maximal states 
	is strongly connected if and only if 
	there exists a unique sink component in~$\mathcal A$ (which then equals~$S$).	
\end{proposition}
\begin{proof}
	First, suppose~$S$ is strongly connected.
	Since two sink component consist of incomparable 
	elements, it follows that~$S$ equals
	a sink component, which is the only sink component in~$\mathcal A$.
	
	Conversely, suppose there exists a unique sink component.
	Since~$S$ is the union of all sink components by \cref{prelim:general:order:set-of-maximal-elements}
	this unique sink component equals~$S$ and so~$S$ is strongly connected.	
\end{proof}

\section{Automata as Language Acceptors} 
\label{prelim:sec:autacceptors}

\subsection{Regular Languages}
\label{prelim:sec:regularlang}

%!TEX root = ./../../thesis-trier.tex

\todo[inline]{Hier an einigen Stellen nicht die Makros für semiautomaton und automaton benutzt.}
\IN{start state!of an automaton} 
\index{initial state|see{start state}}
\IN{final state set!of an automaton}
\IN{regular language}
\IN{accepted!language!by automaton}
\IN{accepted!word!by automaton}
\IN{automaton!as language acceptor}
\index{language acceptor automaton|see{automaton!as language acceptor}}
\index{recognized language|see{accepted language}}
\label{def:automaton-acceptor}
\Cref{intro:chap:short_intro_sync_aut} introduced semiautomata as mere finite-state transition systems.
Here we associate languages with automata.
A \emph{(deterministic and complete) finite automaton\footnote{In the literature, most of the time, automata are meant to be language
		acceptors without a special term and automata without
		a start state and a final state are called semiautomata.		
	} (as a language acceptor)} 
is a $5$-tuple $\automat{A}{Q}{\Sigma}{\delta}{q_0}{F}$ 
where $(Q, \Sigma, \delta)$ is a semiautomaton as defined before, $q_0$
is the \emph{start} (or \emph{initial) state} and $F \subseteq Q$ is the \emph{set of final states}. 
The language \emph{accepted}, or \emph{recognized}, by $\mathcal A$
is
\[
L(\mathcal A) = \{ u \in \Sigma^* \mid \applytrans{\delta}{q_0}{u} \in F \}, \label{def:accepted_lang_aut}
\]
i.e., the set of labels of paths going from the start state to some final state.
We also say a word $u \in \Sigma^*$ is \emph{accepted} by $\mathcal A$ if $u \in L(\mathcal A)$.\footnote{There should be no problem with using the modifier ``regular'' in the context of languages and in the context of permutation groups since both domains are sufficiently different
so that no confusion should arise.}	
A language $L \subseteq \Sigma^*$ is called \emph{regular} if there exists an automaton $\mathcal A$ 
such that $L = L(\mathcal A)$.\footnote{Eilenberg~\cite{Eilenberg1974} called
	languages acceptable by automata recognizable,
	which makes sense when considering more general monoids
	than the free monoid. Since we do not consider the more general case,
	we do not use this terminology here.	
	See also the related footnote on page~\pageref{def:RegSigma}.}	

Kleene~\cite{Kleene56} has shown that regular expressions denote precisely
the set of regular languages. A proof (in contemporary notation)
can be found in any textbook on formal language theory, e.g.~\cite{HopUll79,Pin2022}.

\begin{theorem}[Kleene's theorem~\cite{Kleene56}]
	% For every regular language $L \subseteq \Sigma^*$ there exists a regular
	% expression $R$ such that $
	The set of regular languages is precisely the set of languages
	describable by regular expressions, i.e., 
	for every $\automat{A}{Q}{\Sigma}{\delta}{q}{F}$
	there exists a regular expression $R$ such that $L(\mathcal A) = L(R)$
	and vice versa.
\end{theorem}


\IN{accessible!state!in automaton} 
\IN{accessible!automaton} 
\IN{coaccessible state!in automaton}
\index{initially connected|see{accessible automaton}}
Let $\automat{A}{Q}{\Sigma}{\delta}{q_0}{F}$
and $q \in Q$. The state $q$ is called \emph{accessible} if it is reachable from the start state, i.e.,  there exists a word $u \in \Sigma^*$
such that $\applytrans{\delta}{q_0}{u} = q$. 
It is called \emph{coaccessible} if there exists a word $u \in \Sigma^*$
such that $\applytrans{\delta}{q}{u} \in F$. We call $\mathcal A$ \emph{accessible} (or \emph{initially connected}) if all states are accessible. 
Obviously, for every regular language there exists an accessible automaton accepting it.


%Observe that a permutation automaton
%as a language acceptor is accessible if and only if it is strongly connected. \todo{im minimalaut teil sagen perm aut mit einem finalzsutand immer minimal}

For a semiautomaton $\semiautomat{A}{Q}{\Sigma}{\delta}$, $q \in Q$
and $S \subseteq Q$, by 
\[\mathcal A_{q,S} = (Q, \Sigma, \delta, q, S) \label{def:add_start_final}
\]
we denote the language acceptor automaton that results from $\mathcal A$ 
by choosing $q$ as a start state and $S$ as a set of final states.

There should be no confusion by calling both types of automata simply automata, as it should be clear
from the context whether we mean automata as language acceptors, or whether we are talking about 
semiautomata. Furthermore, both types are specified by a different sets of data, a $3$-tuple versus a $5$-tuple, which distinguishes them notation-wise. 
In case any ambiguity might arise, and it is necessary to differentiate between both notions,
we revert to using the proper names semiautomaton and automaton as language acceptor.
% or by specifying them explicitly  
%either as a $3$-tuple or a $5$-tuple. \todo{daraus kann man es ableiten.}
%However, if there might arise confusion because both notions appear in a single statement, 
%we use the more specific terminology.

Every automaton as a language acceptor is also a semiautomaton when forgetting about the start state
and the set of final states. So, all semiautomata notions transfer to automata as language acceptors
and we use them also for language acceptor automata in this sense.
%In this sense, all notions that we define and that do not involve a
%start state and a set of final state are to be regarded as definitions for both types of automata.\todo{english.}

Given automata $\automat{A_i}{Q_i}{\Sigma}{\delta_i}{q_i}{F_i}$
where $i \in \{1,\ldots,m\}$ and $m \ge 2$, 
for the product automaton $\mathcal A_1 \times \ldots \times \mathcal A_m$
we have the following:
\begingroup
\allowdisplaybreaks 
\begin{align}
	L(\mathcal A_1) \setminus L(\mathcal A_2) & = L((\mathcal A_1 \times \mathcal A_2)_{(q_1, q_2), ( F_1 \times Q )\setminus (Q \times F_2)}), \\
	\bigcap_{i=1}^m L(\mathcal A_i) & = L((\mathcal A_1 \times \cdots \times \mathcal A_m)_{(q_1,\ldots,q_m), F_1 \times \cdots \times F_m}), \label{prelim:eqn:product-intersection}	\\
	\bigcup_{i=1}^m L(\mathcal A_i) & = L((\mathcal A_1 \times \cdots \times \mathcal A_m)_{(q_1,\ldots,q_m), E}), 
\end{align}
\endgroup%
where $E = \bigcup_{i=1}^n Q \times \cdots \times \underbracket[0.4pt]{F_i}_{\clap{\scriptsize$i$-th position}} \times \cdots \times Q$.

Next, we state two results that relate the accepted language to morphic images and congruences.

\begin{proposition}
\label{prelim:prop:acceptor-morphism}
	Let $\varphi \colon Q \to S$ be a morphism between $\semiautomat{A}{Q}{\Sigma}{\delta}$
	and $\semiautomat{B}{S}{\Sigma}{\mu}$.
	Then $L(\mathcal A_{q_0,F}) \subseteq L(\mathcal B_{\applyfunc{q_0}{\varphi}, \applyfunc{F}{\varphi}})$
	for $q_0 \in Q$ and $F \subseteq Q$.
	Furthermore, if $\applyfunc{\applyfunc{F}{\varphi}}{\varphi^{-1}} = F$,
	then $L(\mathcal B_{\applyfunc{q_0}{\varphi}, \applyfunc{F}{\varphi}}) \subseteq L(\mathcal A_{q_0,F})$.
\end{proposition}
\begin{proof}
	Let $u \in \Sigma^*$.
	For $\applytrans{\delta}{q_0}{u} \in F$,
	we have $\applytrans{\mu}{\applyfunc{q_0}{\varphi}}{u} = \applyfunc{\applytrans{\delta}{q_0}{u}}{\varphi} \in \applyfunc{F}{\varphi}$ and so $u \in L(\mathcal B_{\applyfunc{q_0}{\varphi}, \applyfunc{F}{\varphi}})$.
	Next, if $\applyfunc{\applyfunc{F}{\varphi}}{\varphi^{-1}} = F$
	and $u \in L(\mathcal B_{\applyfunc{q_0}{\varphi}, \applyfunc{F}{\varphi}})$,
	then we get $\applyfunc{\applytrans{\delta}{q_0}{u}}{\varphi} = \applytrans{\mu}{\applyfunc{q_0}{\varphi}}{u} \in  \applyfunc{F}{\varphi}$
	and so $\applytrans{\delta}{q_0}{u} \in F$.
\end{proof}

\Cref{prelim:prop:congruence_induces_morphism} 
gives the following corollary to \cref{prelim:prop:acceptor-morphism}.

\begin{corollary}
\label{prelim:cor:acceptor-congruence}
  Let $\semiautomat{A}{Q}{\Sigma}{\delta}$
  be an automaton and $\sim$ be a $\TS{\mathcal A}$-congruence.
  Then for each $q \in Q$ and $F \subseteq Q$
  we have $L(\mathcal A_{q,F}) \subseteq L(( \mathcal A / {\sim} )_{[q]_{\sim}, [F]_{\sim}})$
  and if $F$ saturates~$\sim$, then
  $L(\mathcal A_{q,F}) = L(( \mathcal A / {\sim} )_{[q]_{\sim}, [F]_{\sim}})$ (note that $[F]_{\sim} = \{ [q]_{\sim} \mid q \in F\}$).  
\end{corollary}

Morphisms as introduced here are a concept defined for semiautomata. \Cref{prelim:prop:acceptor-morphism}
gives a way to extend this notions to morphisms for automata as language acceptors
by also stipulating that the start state is preserved
and the final state set saturates the kernel relation. We do not use such an extended notion of morphisms
for automata as language acceptors and thus can use the term morphism alone without any ambiguity here.

\subsection{The Minimal Automaton}
\label{prelim:subsubsec:min_aut}

%!TEX root = ./../../springer-mono-format/book.tex
%!TEX root = ./../../thesis-trier.tex

Here, we introduce the Myhill-Nerode relation,
the minimal automaton of a given language, the canonical map of
an accessible automaton onto the minimal automaton, and the notion of distinguishability of states.
\Cref{prelim:min_aut:terminal_object} gives a characterization of regularity in terms of the Myhill-Nerode relation
and shows that the canonical map is a surjective morphism.
Then, in \cref{prelim:thm:char_minaut}, we characterize those automata
with the least number of states accepting a given regular language.
\Cref{prelim:cor:uniform-min} characterizes those automata %(which must be strongly connected)
having the least number of states for every distribution of a start state and some non-empty and proper set of final states.

\IN{Myhill-Nerode relation}
\label{def:Nerode}

Let $L \subseteq \Sigma^*$ be a language.
\begin{enumerate}
\item  For $u,v \in \Sigma^*$ define the 
\emph{Myhill-Nerode relation} w.r.t.~$L$~\cite{myhill1957finite,nerode1958linear} by 
\[ 
u \equiv_L v
\Leftrightarrow 
u^{-1}L = v^{-1}L.
\]
The relation $\equiv_L$ is an equivalence relation that saturates $L$ and such that, for all
$u,v,w \in \Sigma^*$, we have\footnote{Every word $w \in \Sigma^*$
	can be associated to a unique transformation ${\sigma}_w \colon \Sigma^* \to \Sigma^*$
	given by $\sigma_w(u) = uw$ for $u \in \Sigma^*$.
	Then $\hat{\Sigma^*} = \{ {\sigma}_w \mid w \in \Sigma^* \}$
	gives a transformation monoid on $\Sigma^*$
	and \cref{prelim:eqn:Nerodecongruence} 
	states that $\equiv_L$ is a $\hat{\Sigma^*}$-congruence.
	%This construction can be carried out for any monoid
	%	and the resulting congruences are called \emph{right-congruences}. \IN{right-congruence relation}
} 
\begin{equation}\label{prelim:eqn:Nerodecongruence}
	u \equiv_L v \Rightarrow uw \equiv_L vw.
\end{equation}

\index{minimal automaton} 
\label{def:minaut}

\item The \emph{minimal automaton (of $L$)} is the
(here possibly infinite) acceptor automaton 
$\minaut{L} = (Q_L, \Sigma, \delta_L, [\varepsilon]_{\equiv_L}, F_L)$
with 
\begin{align*}
	Q_L & = \Sigma^* / {\equiv_L}, \\
	\delta_L([u]_{\equiv_L}, v) & = [uv]_{\equiv_L}, \\
	F_L & = \{ [u]_{\equiv_L} \mid u \in L \}
\end{align*}
where $\delta_L$ is well-defined by \cref{prelim:eqn:Nerodecongruence}.
Clearly, $\minaut{L}$ is accessible.

\index{canonical map of $\mathcal A$ onto $\minaut{L}$}
\label{def:Nerode-morphism}

\item Suppose $L$ is regular
 and $L = L(\mathcal A)$ for some accessible automaton $\automat{A}{Q}{\Sigma}{\delta}{q_0}{F}$.
 Then, for $u,v \in \Sigma^*$,
 \begin{equation}\label{prelim:eqn:can-map}
 	\applytrans{\delta}{q_0}{u} = \applytrans{\delta}{q_0}{v}
 	\Rightarrow
 	u \equiv_L v.
 \end{equation} 
 Let $\minaut{L} = (Q_L, \Sigma, \delta_L, [\varepsilon]_{\equiv_L}, F_L)$.
 By \cref{prelim:eqn:can-map}, the map $\pi_{\mathcal A} \colon Q \to Q_L$, called the \emph{canonical map of~$\mathcal A$ onto~$\minaut{L}$},
 given by \[ 
 \applyfunc{q}{\pi_{\mathcal A}} = [u]_{\equiv_L}
 \]
 for some word $u \in \Sigma^*$ such that $\applytrans{\delta}{q_0}{u} = q$
 is well-defined.
 So, $\pi_{\mathcal A}(\delta(q_0, u)) = [u]_{\equiv_L}$
 and $\pi_{\mathcal A}$ is surjective.
  
\IN{distinguishable states}
\item Let $\automat{A}{Q}{\Sigma}{\delta}{q_0}{F}$ be an automaton. 
Two states $p, q \in Q$ are \emph{distinguishable} (in $\mathcal A$) 
if there exists $u \in \Sigma^*$
such that 
\begin{equation}\label{prelim:eqn:distinguish}
	( \applytrans{\delta}{p}{u} \in F \land  \applytrans{\delta}{q}{u} \notin F )
	\lor 
	( \applytrans{\delta}{p}{u} \notin F \land \applytrans{\delta}{q}{u} \in F ).
\end{equation}
We call $p$ and $q$ \emph{indistinguishable} if they are not distinguishable. The indistinguishable relation
is a quasiorder.
\IN{indistinguishable states}
\end{enumerate}

Suppose $\sim$ is a~$\TS{\mathcal A}$-congruence that saturates $F \subseteq Q$
for  $\automat{A}{Q}{\Sigma}{\delta}{q_0}{F}$.
Then the following implication holds for $p,q \in Q$:
\begin{equation} \label{prelim:eqn:congruence-implies-indist}
	p \sim q \Rightarrow \text{$p$ and $q$ are indistinguishable},
\end{equation}
i.e., the indistinguishability relation is the coarsest $\TS{\mathcal A}$-congruence
saturating~$F$. Note that, by \cref{prelim:eqn:distinguish}, two states $p,q$
are indistinguishable in $\mathcal A$ if and only if
\[
 \forall u \in \Sigma^* \colon \delta(p, u) \in F \Leftrightarrow \delta(q, u) \in F.
\]

% before we characterize ...
We aim to characterize minimality in terms of $\equiv_L$, congruences,
morphisms, and the distinguishability relation in \cref{prelim:thm:char_minaut}.
But before, % below in~\cref{prelim:lem:indist},
we list relations, needed in the proof of~\cref{prelim:thm:char_minaut}, between $\equiv_L$, accessible automata~$\mathcal A$ and $\pi_{\mathcal A}$. The next 
lemma follows readily from the definitions.

\begin{lemma}\label{prelim:lem:indist}
	Let $\automat{A}{Q}{\Sigma}{\delta}{q_0}{F}$ be accessible and $L = L(\mathcal A)$. 
	Then, for $u,v \in \Sigma^*$, the following are equivalent:	
	\begin{enumerate}
		\item $\pi_{\mathcal A}(\applytrans{\delta}{q_0}{u}) = \pi_{\mathcal A}(\applytrans{\delta}{q_0}{v})$,
		
		\item $u \equiv_L v$, 
		
		\item $\applytrans{\delta}{q_0}{u}$ and $\applytrans{\delta}{q_0}{v}$ are indistinguishable.
	\end{enumerate}
\end{lemma}

Next, we summarize the important properties of $\minaut{L}$ and~$\pi_{\mathcal A}$.

\begin{theorem}
	\label{prelim:min_aut:terminal_object}
	Let $L \subseteq \Sigma^*$ and $\minaut{L} = (Q_L, \Sigma, \delta_L, [\varepsilon]_{\equiv_L}, F_L)$.
	Then $L$ is regular if and only 
	if $Q_L$ is finite. Further,
	if $L$ is regular, then for every accessible $\automat{A}{Q}{\Sigma}{\delta}{q_0}{F}$
	with $L = L(\mathcal A)$,
	the map $\pi_{\mathcal A} \colon Q \to Q_L$
	is a surjective morphism 
    such that $\pi_{\mathcal A}(q_0) = [\varepsilon]_{\equiv_L}$, 
    $\pi_{\mathcal A}^{-1}(F_L) = F$ and
    $\minaut{L}$ accepts $L$.
\end{theorem}
\begin{proof}
	First, suppose $Q_L$ is finite.
	As $\equiv_L$ saturates $L$, for $u \in \Sigma^*$, we 
	get
	\[
	u \in L \Leftrightarrow \exists v \in L \colon u \equiv_L v 
	        \Leftrightarrow [u]_{\equiv_L} \in F_L
	        \Leftrightarrow \delta_L([\varepsilon]_{\equiv_L}, u) \in F_L.
	\]
	So, $\minaut{L}$
	is a finite automaton accepting $L$.
	Hence, $L$ is regular.	
	Conversely, suppose $L$ is regular.
	Let $\automat{A}{Q}{\Sigma}{\delta}{q_0}{F}$ be
	a finite automaton with $L = L(\mathcal A)$, which we can assume to be accessible.
	As $\pi_{\mathcal A}$ is surjective, we find
	that $Q_L$ is finite.	
	Next, we show the remaining claims w.r.t.~$\mathcal A$.
	Let $q \in Q$ and $u \in \Sigma^*$ with $\delta(q_0, u) = q$.
	Then, by the definitions of~$\pi_{\mathcal A}$ and~$\delta_L$, 
	we have, for $v \in \Sigma^*$,
	\[
	\pi_{\mathcal A}(\delta(q, v)) = \pi_{\mathcal A}(\delta(q_0, uv)) = [uv]_{\equiv_L} = \delta_L([u]_{\equiv_L}, v) 
	= \delta_L(\pi_{\mathcal A}(q), v).
	\]
    So, $\pi_{\mathcal A}$ is a morphism from $\mathcal A$ onto $\minaut{L}$.
	Lastly, by the definition of $F_L$, $\pi_{\mathcal A}$ and as~$\equiv_L$ saturates $L$, we have
	\begin{align*}
		\pi_{\mathcal A}(q) \in F_L 
		& \Leftrightarrow \exists u \in L \ \exists v \in \Sigma^* \colon \delta(q_0, v) = q \land v \equiv_L u & \tag{Def.~$\pi_{\mathcal A}$ \& $F_L$} \\
		& \Leftrightarrow \exists u \in L \ \exists v \in L \colon \delta(q_0, v) = q \land v \equiv_L u & \tag{$\equiv_L$ saturates $L$} \\
		& \Leftrightarrow \exists v \in L \colon \delta(q_0, v) = q  \\
		& \Leftrightarrow q \in F.
	\end{align*}
	Therefore, $\applyfunc{F_L}{\pi_{\mathcal A}^{-1}} = F$.
	By \cref{prelim:prop:acceptor-morphism}, we find $L(\mathcal A) = L(\minaut{L})$.
\end{proof}

Brzozowski~\cite{DBLP:journals/jacm/Brzozowski64}
has shown that a language is regular if and only if the set of its left quotients is finite, which follows from \cref{prelim:min_aut:terminal_object}.

\begin{corollary}[Brzozowski~\cite{DBLP:journals/jacm/Brzozowski64}]
	\label{prelim:thm:regular_iff_quotients_finite}
	A language $L \subseteq \Sigma^*$
	is regular if and only if $\{ u^{-1}L \mid u \in \Sigma^*\}$
	is finite.
\end{corollary}
\begin{proof}
 By the defining equation of the Myhill-Nerode relation, $[u]_{\equiv} \mapsto u^{-1}L$
 is a bijection between~$Q_L$ and the different left quotients.
 Thus, the claim follows from \cref{prelim:min_aut:terminal_object}.
\end{proof}

Next, we state various characterizations for an automaton
to have the least number of states for its accepted language.

\begin{theorem}\label{prelim:thm:char_minaut}
	Let $L \subseteq \Sigma^*$ be regular, $\minaut{L} = (Q_L, \Sigma, \delta_L, [\varepsilon]_{\equiv_L}, F_L)$
	and $\automat{A}{Q}{\Sigma}{\delta}{q_0}{F}$ with $L = L(\mathcal A)$.
	Then the following are equivalent.
\begin{enumerate}
	
\item \label{thm:minaut:item:leaststates} $\mathcal A$ has the least number of states such that $L = L(\mathcal A)$.

\item \label{thm:minaut:item:morpism} $\mathcal A$ is accessible and there exists no automaton $\automat{B}{S}{\Sigma}{\mu}{s_0}{E}$
and a morphism $\varphi \colon Q \to S$
with $\varphi(q_0) = s_0$ and $\applyfunc{E}{\varphi^{-1}} = F$
such that $|\applyfunc{Q}{\varphi}| < |Q|$. %$1 < |\applyfunc{Q}{\varphi}| < |Q|$.

% den satz oben löschen	und eingangs sagen item 2 rechtfertig name minimal
\item \label{thm:minaut:item:piA} $\mathcal A$ is accessible and $\pi_{\mathcal A} \colon Q \to Q_L$ is an isomorphism of $\mathcal A$ onto $\minaut{L}$.
% with $\pi_{\mathcal A}(q_0) = [\varepsilon]_{\equiv_L}$ and $\pi_{\mathcal A}^{-1}(F_L) = F$.

\item \label{thm:minaut:item:statesdist} $\mathcal A$ is accessible and all pairs of distinct states are distinguishable.

\item \label{thm:minaut:item:congruence} $\mathcal A$ is accessible and the equality relation is the only $\TS{\mathcal A}$-congruence 
 saturating~$F$.
 
\end{enumerate}
\end{theorem}
\begin{proof} $\ref{thm:minaut:item:leaststates} \Rightarrow \ref{thm:minaut:item:morpism}$:
	Set $T = \,\upperset{\{q_0\}}$.
	The automaton $( \mathcal A_{|T} )_{q_0, T \cap F}$ is accessible 
	and we have $L(( \mathcal A_{|T} )_{q_0, T \cap F}) = L(\mathcal A)$.
	By assumption, we find $T = Q$. So, $\mathcal A$ is accessible.
	Next, let $\mathcal B$ and $\varphi$ be as in the statement.
	We can assume $\varphi$ is surjective, i.e., $S = \varphi(Q)$.
	So, $\varphi(F) = \varphi(\varphi^{-1}(E)) = E$.
	Furthermore, we have
	\[
	\varphi^{-1}(\varphi(F)) = \varphi^{-1}(\varphi(\varphi^{-1}(E))) = \varphi^{-1}(E) = F.
	\]
	Then, by \cref{prelim:prop:acceptor-morphism},
	we find $L(\mathcal A) = L(\mathcal B)$.
	Since $\mathcal A$ has the least number of states
	among all automata accepting $L(\mathcal A)$,
	we get $|S| = |Q|$.

$\ref{thm:minaut:item:morpism} \Rightarrow \ref{thm:minaut:item:piA}$: 
By \cref{prelim:min_aut:terminal_object}
we have  $\pi_{\mathcal A}(q_0) = [\varepsilon]_{\equiv_L}$, $\pi_{\mathcal A}^{-1}(F_L) = F$
and $\pi_{\mathcal A}$ is surjective so that $|Q_L| \le |Q|$.
By assumption, we get $|Q| = |Q_L|$
and so $\pi_{\mathcal A}$ is bijective.
	
$\ref{thm:minaut:item:piA} \Rightarrow~\ref{thm:minaut:item:statesdist}$: The automaton
$\minaut{L}$ is accessible and as $\minaut{L} \cong \mathcal A$,
also $\mathcal A$ is accessible.
Let $p,q \in Q$ be two distinct states.
As $\pi_{\mathcal A}$ is injective, we have $\pi_{\mathcal A}(p) \ne \pi_{\mathcal A}(q)$.
By \cref{prelim:lem:indist},
it follows that $p$ and $q$ are distinguishable.

$\ref{thm:minaut:item:statesdist} \Rightarrow~\ref{thm:minaut:item:congruence}$: Suppose $\sim$ 
is some $\TS{\mathcal A}$-congruence that saturates $F$.
As any pair of distinct states $p,q \in Q$ is distinguishable,
we find, by Formula~\eqref{prelim:eqn:congruence-implies-indist},
that $p \not\sim q$. So $\sim$ equals the equality relation.

$\ref{thm:minaut:item:congruence} \Rightarrow \ref{thm:minaut:item:leaststates}:$ Suppose there exists an accessible 
automaton $\automat{B}{S}{\Sigma}{\mu}{s_0}{E}$ 
with $|S| < |Q|$ and $L = L(\mathcal B)$.
By \cref{prelim:min_aut:terminal_object}, as $\pi_{\mathcal B}$ is surjective, 
we have $|Q_L| \le |S|$. 
Thus $|Q_L| < |Q|$
and so~$\pi_{\mathcal A}$ maps to distinct states to one state in $\minaut{L}$,
giving that $\sim_{\pi_{\mathcal A}}$ does not equal the equality relation.
By \cref{prelim:thm:homtheorem,prelim:min_aut:terminal_object}, $\sim_{\pi_{\mathcal A}}$
is a $\TM{\mathcal A}$-congruence, which, as $\pi_{\mathcal A}^{-1}(F_L) = F$,
saturates $F$.
\end{proof}

\IN{state complexity}
For regular $L \subseteq \Sigma^*$, we call $\min\{ |Q| : \automat{A}{Q}{\Sigma}{\delta}{q_0}{F} \text{ with } L = L(\mathcal A)\}$
the \emph{state complexity} of $L$,
which equals, by \crefitem{prelim:thm:char_minaut}{thm:minaut:item:piA},
the number of states of $\minaut{L}$.
Furthermore, all automata with the least number of states
accepting~$L$ are isomorphic by	\crefitem{prelim:thm:char_minaut}{thm:minaut:item:piA}, which justifies calling $\minaut{L}$
the minimal automaton. 
In view of this, we 
also call $\automat{A}{Q}{\Sigma}{\delta}{q_0}{F}$ \emph{minimal}
if~$\pi_{\mathcal A}$ is an isomorphism.
Note that two non-coaccessible states are indistinguishable. So, a minimal automaton has at most one
state that is not coaccessible and this state must be a trap state.

\begin{remark}
	Since the indistinguishability relation
	is the coarsest $\TS{\mathcal A}$-congruence saturating~$F$,
	by \crefitem{prelim:thm:char_minaut}{thm:minaut:item:congruence}
	the automaton 	$\mathcal A = (Q,\Sigma,\delta,q_0,F)$ is minimal if and only if it is accessible
	and the indistinguishability relation equals the equality relation.
\end{remark}

\begin{remark}
 Accessibility in Items~\ref{thm:minaut:item:morpism},\ref{thm:minaut:item:statesdist} and~\ref{thm:minaut:item:congruence} 
 of \cref{prelim:thm:char_minaut} 
 is crucial. For example, in an automaton $\automat{A}{Q}{\Sigma}{\delta}{q_0}{F}$
 whose state set forms two cycles of size at least two induced by different letters $a$ and~$b$ that induce self-loops on the other cycle
 and with one final state on each cycle, all pairs of distinct states are distinguishable,
 equality is the only $\TS{\mathcal A}$-congruence saturating~$F$,
 and there exists no morphism as stated in Item~\ref{thm:minaut:item:morpism}.
 This follows by considering $\mathcal A_{|\{a\}}$
 it either permutes the states on the $a$-cycle or maps all those states
 to a single state outside the $a$-cycle. But since this is also a morphism of $\mathcal A_{|\{b\}}$,
 it must permute the full $b$-cycle if it maps outside the $a$-cycle.
 The same argument for the $b$-cycle gives the it has to permute
 the $a$-cycle and the $b$-cycle.
 Thus it is an automorphism.
\end{remark}

Davies~\cite[Lem.~3.7]{DBLP:journals/mst/Davies18} (see also~\cite[Lem.~4.2.8]{DBLP:phd/basesearch/Davies19}) proved 
that an automaton with a non-empty and proper final state set is minimal
if and only if it admits no non-trivial $\TM{\mathcal A}$-congruence saturating
the final state set. This readily follows 
from \crefitem{prelim:thm:char_minaut}{thm:minaut:item:congruence},
as the final state set is empty or equals the whole state
set if and only if the accepted language is empty or equals $\Sigma^*$, respectively.
Further, she also characterized
strongly connected automata that are minimal for every start state and proper non-empty set of final states
by the absence of non-trivial $\TS{\mathcal A}$-congruences~\cite[Prop.~4.2.3]{DBLP:journals/mst/Davies18} (or~\cite[Prop.~3.2]{DBLP:phd/basesearch/Davies19}).
Next, we state this together with other characterizations as a corollary to \cref{prelim:thm:char_minaut}.

Note that for strongly connected $\semiautomat{A}{Q}{\Sigma}{\delta}$ 
with at least two states, $L(\mathcal A_{q,F})$
is a non-empty proper subset of~$\Sigma^*$ for $q \in Q$
and non-empty and proper $F \subsetneq Q$
since $\delta(q, u) \in F$
and $\delta(q, v) \notin F$
for some $u,v \in \Sigma^*$ by accessibility and strong connectedness.

\begin{corollary}\label{prelim:cor:uniform-min}
Let $\semiautomat{A}{Q}{\Sigma}{\delta}$.
Then the following are equivalent.
\begin{enumerate}
\item \label{prelim:cor:uniform-min:1} $\mathcal A_{q,F}$ is minimal for each~$q \in Q$ and non-empty and proper $F \subsetneq Q$.

\item \label{prelim:cor:uniform-min:2} $\mathcal A$ is strongly connected and
 all pairs of distinct states are distinguishable in~$\mathcal A_{q,F}$
 for all $q \in Q$ and $\emptyset \ne F \subsetneq Q$.

\item \label{prelim:cor:uniform-min:3} $\mathcal A$ is strongly connected and there exists no non-trivial $\TS{\mathcal A}$-congruence.

\item \label{prelim:cor:uniform-min:4} $\mathcal A$ is strongly connected and there exists no $\semiautomat{B}{S}{\Sigma}{\mu}$
with a morphism $\varphi \colon Q \to S$
such that $1 < |\applyfunc{Q}{\varphi}| < |Q|$.
\end{enumerate}
\end{corollary} 
\begin{proof}
$\ref{prelim:cor:uniform-min:1} \Rightarrow \ref{prelim:cor:uniform-min:2}$:
Suppose $\mathcal A_{q, F}$ is minimal for each $q \in Q$
and $\emptyset \ne F \subsetneq Q$.
By \crefitem{prelim:thm:char_minaut}{thm:minaut:item:statesdist},
all pairs of distinct states are distinguishable in $\mathcal A_{q,F}$.
Fix some $\emptyset \ne F \subsetneq Q$. By \crefitem{prelim:thm:char_minaut}{thm:minaut:item:statesdist},
$\mathcal A_{q,F}$ is accessible w.r.t.~each state~$q$. Now,
by \crefitem{prelim:prop:sink-comp}{prelim:eqn:sink-comp},
 $\mathcal A_{q,F}$ being accessible for any $q \in Q$
is equivalent to $\mathcal A$ being strongly connected. 
Alternatively, let $S \subseteq Q$ be a sink component
and $q \in S$. As the states in $Q \setminus S$ are not accessible
in $\mathcal A_{q, S}$, we must have $Q = S$ since minimality implies that all states
are accessible by \cref{prelim:thm:char_minaut}.
So $\mathcal A$ is strongly connected. 

$\ref{prelim:cor:uniform-min:2} \Rightarrow \ref{prelim:cor:uniform-min:3}$:
This follows from \cref{prelim:eqn:congruence-implies-indist}. 

$\ref{prelim:cor:uniform-min:3} \Rightarrow \ref{prelim:cor:uniform-min:4}$:
This follows from \crefitem{prelim:thm:homtheorem}{prelim:thm:homtheorem:2}.

$\ref{prelim:cor:uniform-min:4} \Rightarrow \ref{prelim:cor:uniform-min:1}$:
Let $q \in Q$ and $\emptyset \ne F \subsetneq Q$. Since $\mathcal A$
is strongly connected, $\mathcal A_{q,F}$ is accessible.
As there exists no morphism onto an automaton with less states,
by \crefitem{prelim:thm:char_minaut}{thm:minaut:item:morpism}, $\mathcal A_{q, F}$ is minimal.
\end{proof}

Note that the equivalences from \cref{prelim:cor:uniform-min} 
also hold for $|Q| = 1$ vacuously since then there exists no
proper and non-empty final state set and no morphism $\varphi$
with $1 < |\varphi(Q)| < |Q|$. A one-state automaton
only accepts $\emptyset$ or $\Sigma^*$.

\Cref{prelim:cor:simpleaut} stated that an automaton admitting no non-trivial $\TM{\mathcal A}$-congruence
is either strongly connected or the state set is partitioned into a trap state and a strongly connected component.
\Cref{prelim:cor:uniform-min} gives more details in case the automaton is strongly connected
and admits no non-trivial $\TM{\mathcal A}$-congruence.

\begin{remark}
	Strong connectivity is necessary in Items~\ref{prelim:cor:uniform-min:2},\ref{prelim:cor:uniform-min:3},
	and~\ref{prelim:cor:uniform-min:4} of \cref{prelim:cor:uniform-min}.
	For example, consider $\mathcal A = (\{0,1,\ldots,n\}, \{a,b\}, \delta)$
	with $\fword{\delta}{a} = (0)(1 ~ 2 ~ \cdots n)$
	and $\fword{\delta}{b} = \SI{1}{0}$ for some number $n > 1$. 
	Then $\mathcal A_{0, \{0\}}$ is not minimal.
	However, we have the following:
	
	\begin{enumerate}
	\item All pairs of distinct states are distinguishable
	in $\mathcal A_{q,F}$ for each $q \in Q$
	and non-empty and proper $F \subseteq Q$,
	for we can move one state to~$0$ while the other remains on the cycle, and
	then can be moved to a state that is final if $0 \notin F$
	or not final if $0 \in F$.
		
	\item The automaton admits no non-trivial $\TS{\mathcal A}$-congruence.
	For suppose $\sim$ is a $\TS{\mathcal A}$-congruence
	such that $p \sim q$ for $p,q \in Q$ with $p < q$.
	Then $\delta(p, b^{n - p + 1}a) = 0$, $\delta(q, b^{n - p + 1}a) \ne 0$
	and thus $\delta(q, b^{n - p + 1}a) \sim 0$.
	This yields $\delta(q, b^{n - p + 1}a b^i ) \sim 0$ for each $i \ge 0$,
	and since $p = \delta(q, b^{n - p + 1}a b^i )$ for some $i \ge 0$
	for each $p \in Q \setminus \{0\}$,
	we find that all pairs of states are congruent w.r.t.~$\sim$.
	
	\item There exists no morphism to a strictly smaller automaton
	with at least two states.	
	\end{enumerate}
	So, each condition except strong connectivity is fulfilled.
\end{remark}

Note that any automaton $\mathcal A$ with at most two states, even those that are not strongly connected,
admits no non-trivial $\TS{\mathcal A}$-congruence, no non-trivial morphic images,
and has the property that distinct states are indistinguishable in~$\mathcal A_{q,F}$ for $q \in Q$ and $\emptyset \ne F \subsetneq Q$.

Restivo \& Vaglica~\cite{DBLP:journals/tcs/RestivoV12,DBLP:journals/tcs/RestivoV12a}
also investigated automata that are minimal for all $k$-sets of final states for some $k \ge 1$
or for all non-empty and proper subsets of final states.

\section{Remarks on Transformation Semigroups and Automata}
\label{prelim:remarks_trans_semigrp_semiautomata}

%!TEX root = ./../../springer-mono-format/book.tex
%!TEX root = ./../../thesis-trier.tex

To every (complete and deterministic) semiautomaton $\semiautomat{A}{Q}{\Sigma}{\delta}$
corresponds its transformation semigroup $\TS{\mathcal A}$
with generating set $\{ \fword{\delta}{a} \colon Q \to Q \mid a \in \Sigma \}$.
Conversely, let $S \le_S \mathcal T_n$ 
be a transformation semigroup
generated by $g_1, \ldots, g_m$. Then for  $\Sigma = \{g_1, \ldots, g_n\}$ (viewing the generating sequence as a set of abstract symbols)
we obtain
a semiautomaton $\semiautomat{A}{[n]}{\Sigma}{\delta}$
with $\TS{\mathcal A} = S$
by setting $\applytrans{\delta}{q}{g_i} = \applyfunc{q}{g_i}$
for $q \in [n]$ and $i \in \{1,\ldots,m\}$.\footnote{A transformation $g_i$ is allowed to appear more than once among  $g_1,\ldots,g_m$.} 
Thus, a semiautomaton may be viewed as a transformation semigroup together
with a distinguished generating sequence.

\begin{example}	\label{prelim:aut:ex:Z6genset}
	The choice of a generating set for a given transformation semigroup $S \le_S \mathcal T_n$
	can lead to quite distinct automata.
	For example, consider the permutation group $G = \langle g \rangle_G$ with $g = (0 ~ 1 ~ \cdots ~ 5)$.
	When constructing the automaton with generating set $\{ g \}$ 
	we obtain an automaton that cyclically permutes its states.
	Another generating set is given by $g^2 = (0 ~ 2 ~ 4) (1 ~ 3 ~ 5)$ and $g^3 = (0 ~ 3) (1 ~ 4) (2 ~ 5)$
	since $g^2 g^2 g^3 = g$.
	The automaton we obtain from this generating set is depicted in \cref{prelim:aut:fig:Z6genset}.	
	
	\begin{figure}[ht]
		\centering	
		\begin{tikzpicture}[node distance = 6em and 10em, on grid, auto, ->]
			\node (q0) [state] {$0$};
			\node (q2) [state, right = of q0] {$2$};
			\node (q4) [state, right = of q2] {$4$};
			\node (q3) [state, below = of q0] {$3$};
			\node (q5) [state, below = of q2] {$5$};
			\node (q1) [state, below = of q4] {$1$};
			
			\path (q0) edge node {$g^2$} (q2)
			(q2) edge node {$g^2$} (q4)
			(q4) edge [bend right=23,above] node {$g^2$} (q0);
			
			\path (q3) edge node {$g^2$} (q5)
			(q5) edge node {$g^2$} (q1)
			(q1) edge [bend left=23] node {$g^2$} (q3);			
			
			\path (q0) edge [bend left] node {$g^3$} (q3)
			(q3) edge [bend left] node {$g^3$} (q0);
			
			\path (q2) edge [bend left] node {$g^3$} (q5)
			(q5) edge [bend left] node {$g^3$} (q2);
						
			\path (q4) edge [bend left] node {$g^3$} (q1)
			(q1) edge [bend left] node {$g^3$} (q4);
			
		\end{tikzpicture}
		\caption[Automaton constructed from transformation semigroup.]{Automaton constructed from permutation group, see \cref{prelim:aut:ex:Z6genset}.}
		\label{prelim:aut:fig:Z6genset}
	\end{figure}		
\end{example}

Many of our notions and results
can be formulated either for automata or for transformation semigroups.
For example, $\mathcal A$ is synchronizing if and only if $\TS{\mathcal A}$
contains a constant transformation.

Our results are motivated by automata-theoretical problems, but properties only
related to subset mapping properties do not depend on the choice of some generating set 
and are actually properties of transformation
semigroups. So, these properties are better expressed in an algebraic setting
with transformation semigroups, as done in \cref{part:aut}. The corresponding automata-theoretic 
reformulations are usually straightforward and are therefore omitted.

A semiautomaton is a machine model that processes, or reacts to, inputs.
In this setting, the alphabet is crucial, as it is the means by which
the machine interacts with its environment.
Since inputs correspond to physical actions and an operator may want
to minimize the effort to achieve a given goal,
it is natural to ask for the minimum number of such 
actions required to do this.
Thus, subset mapping properties with length bounds on shortest words
realizing some property are naturally expressed in an automata-theoretic formulation.

Later in \cref{part:constrsync}, we use automata-theoretic notation exclusively.

We note that there is a generalization of the \Cerny conjecture for transformation monoids~\cite{AnanichevVolkov04,Rystsov2024}.
Shortest words in some generating sets have also been investigated
in the context of the diameter of a Cayley graph depends~\cite{DBLP:journals/ijfcs/GonzeGJGV19}.
%Coxeter groups are intimately related to their generating sets and the element represented
%by the longest possible word in the generators~\cite{}, 
%or the set of words representing geodesics in geometric group theory~\cite{jamesworringten}.

\section{Discussion} 
\label{prelim:automata:discussion}

%!TEX root = ./../../springer-mono-format/book.tex
%!TEX root = ./../../thesis-trier.tex

Some material on automata is presented differently here or is not covered in standard texts, such as~\cite{HopUll79,Ito04,DBLP:books/daglib/0088160}. Specifically, the various facts on semiautomata morphisms from \cref{prelim:subsubsec:endo_aut} appear to be scattered in the literature or, as far as we know, not stated explicitly in the form presented here.
Moreover, the isomorphisms in \cref{prelim:autmoata:Rees:prop:iso-eqns} concerning Rees quotient automata 
and \cref{prelim:prop:Rees-homsfactorthrough} are, to the best of our knowledge, new.
Automata can also be viewed as unary algebras with finitely many operations on a finite domain. Results from automata theory and formal language theory have been extended to unary algebras in~\cite{DBLP:journals/tcs/PetkovicCB04}; see also~\cite[p.~835]{Volkov22} and~\cite{Bogdanovic99} for connections between unary algebras, their identities, and synchronizing automata.


\paragraph{Overview on Automata Theory.} The theory of finite automata as an independent field of study is generally considered to have started with the seminal book 
\emph{Automata Studies}~\cite{ShannonM56}. Medvedev translated this book into Russian and added his own article~\cite{Medvedev56}, in which he introduced semiautomata (therefore sometimes referred to as \emph{Medvedev automata}, see~\cite{DBLP:conf/birthday/Salomaa12,Volkov22}). Early work in automata theory was primarily concerned with the processing of events, signals, and languages by automata rather than with automata viewed merely as finite transition systems.
While models of computation were initially developed in the 1930s to investigate the limits of computability, the focus later shifted to modeling the operation of nervous systems using electronic components. See, for example,~\cite{vonNeumann1951GeneralLogicalTheoryAutomata,Shannon1953ComputersAutomata} for surveys of the work and motivations of the 1940s and 1950s (where automata is also used in a broader sense than here).
Synchronization is a concept tied to the transition structure of an automaton,
or more generally a machine, and thus is a semiautomaton concept.\footnote{In~\cite{HolzerJ18} the notion of synchronization
	was transferred to languages via the minimal automaton
	by calling a language synchronizing if there exists a word
	always ending in the same state in the (possibly infinite) 
	minimal automaton of the language.}

Fleck~\cite{Fleck61,DBLP:journals/jacm/Fleck62} introduced morphisms on automata
and initiated a structure theory for automata centered on their automorphism groups.
For example, among the results obtained in this line of research are the following.
\begin{enumerate}
	\item The automorphism group of a strongly connected automaton
	is a semiregular permutation group~\cite{DBLP:journals/jacm/Fleck62,DBLP:journals/jacm/Weeg62}
	(here stated in \cref{prelim:automata:morphism:prop:endo-SCC-is-iso-b}).\footnote{%
		In~\cite{DBLP:journals/jacm/Fleck62,DBLP:journals/jacm/Weeg62}, semiregularity is referred to simply as ``regular''.
		We follow the terminology that is now standard~\cite{cameron_1999,dixon1996permutation,wielandt1964finite}.}
	\item Let $\semiautomat{A}{\{q_1, \ldots, q_n\}}{\Sigma}{\delta}$ be strongly connected.
	In~\cite{DBLP:journals/jacm/Weeg62}, for $i \in \{1,\ldots,n\}$,
	an equivalence relation $\sim_i$ on~$\Sigma^*$ was introduced
	where, for $u,v \in \Sigma^*$,
	\[
	u \sim_i v \quad\text{if and only if}\quad \applytrans{\delta}{q_i}{u} = \applytrans{\delta}{q_i}{v}.
	\]
	Furthermore, in~\cite[Def.~5]{DBLP:journals/jacm/Weeg62},
	the quotient $\Sigma^* / {\sim_i}$ was equipped with a semigroup structure $(\Sigma^* / {\sim_i}, \ast_i)$
	by setting $[u]_{\sim_i} \ast_i [v]_{\sim_i} = [w]_{\sim_i}$
	precisely if $\applytrans{\delta}{q_i}{u'v'} = \applytrans{\delta}{q_i}{w}$
	for all $u', v' \in \Sigma^*$ with $u' \sim_i u$ and $v' \sim_i v$.
	Barnes~\cite{DBLP:journals/jacm/Barnes65,DBLP:journals/mst/Barnes70}
	used these semigroups to characterize the automorphism group of~$\mathcal A$.
	For instance, this semigroup structure is isomorphic to the automorphism group
	whenever the latter is regular, i.e.\ transitive~\cite[Cor.~1 and Thms.~3 \&~4]{DBLP:journals/jacm/Barnes65}.
\end{enumerate}
The investigation of automata via their automorphism groups was continued in the following works:
\cite{
	DBLP:journals/jacm/Barnes65,
	DBLP:journals/mst/Barnes70,
	DBLP:journals/jacm/Bavel68,
	DBLP:journals/cacm/Bavel70,
	DBLP:journals/jacm/Fleck62,
	DBLP:journals/jacm/Fleck65,
	DBLP:journals/mst/Grzymala-Busse70,
	DBLP:journals/jcss/Grzymala-Busse71,
	DBLP:conf/mfcs/Grzymala-Busse75,
	DBLP:journals/jcss/MasunagaNO73,
	DBLP:journals/jacm/Oehmke63,
	DBLP:journals/jacm/Weeg65,
	DBLP:journals/jacm/Weeg67,
}.

\IN{monoid action}
\IN{$M$-act}
A \emph{monoid action} or \emph{(right) $M$-act}~\cite{KilpKnauerMikhalev2000_MonoidsActs}
is a tuple $(X, M, \mu)$ where~$X$ is some set, $M$ a monoid with identity $1\in M$,
and $\mu \colon X \times M \to X$
a function such that 
\begin{enumerate}
\item $\mu(x, 1) = x$ and

\item $\mu(x, st) = \mu(\mu(x, s), t)$ for $x \in X$ and $s,t \in M$.
\end{enumerate}
A morphism of two $M$-acts $(X,M,\mu)$ and $(Y, M, \rho)$
is a map $\varphi \colon X \to Y$ such that $\varphi(\mu(x, s)) = \rho(\varphi(x),s)$
for all $x \in X$ and $s \in M$.
The morphism is called an endomorphism of the $M$-act
if $X = Y$. Via their extensions to~$\Sigma^*$, semiautomata are precisely the $\Sigma^*$-acts, and morphisms of semiautomata
coincide with the morphisms of these $\Sigma^*$-acts (see also \cite[Prop.~4.5]{KilpKnauerMikhalev2000_MonoidsActs}).
Hence, the study of semiautomata may be viewed as the study of finite $\Sigma^*$-acts.
For a group~$G$, a $G$-act
is also called a group action, $G$-action, or a $G$-set~\cite{cameron_1999,dixon1996permutation}.
Isomorphisms of $G$-acts are called representation equivalences
in~\cite[p.~21]{dixon1996permutation}.
The notion of $G$-acts also arises in algebra, topology,
and geometry~\cite{NeumannStoyThompson1994GroupsAndGeometry}.

\begin{remark}
	An $M$-act can also be identified with a monoid morphism from~$M$ to~$\mathcal T_X$
	equipped with the left-most-first composition operation $;$ from \cref{prelim:sec:remarks_func_comp}
	as its operation, in the following way.
	Given an $M$-act $(X, M, \mu)$, the map $s \mapsto \mu_s$ for $s \in M$
	with $\mu_s(x) = \mu(x, s)$ for $x \in X$ gives such a morphism.
	Conversely, given a morphism $\varphi \colon M \to \mathcal T_X$ of $(M,\cdot)$ to $(T_X, ;)$,
	defining $\mu \colon X \times M$ by $\mu(x, s) = \varphi(s)(x)$
	yields an $M$-act $(X, M, \mu)$.
	For example, every transformation monoid~$M \le_M \mathcal T_X$
	can be considered as an $M$-act by applying
	this construction to the identity transformation on~$M$.	
	
	In this setting, semiautomata can be seen as morphisms $\Sigma^* \to \mathcal T_Q$.
	The image of such a morphism is the transformation monoid.
	This viewpoint was used in~\cite{Rystsov97}.
	Since morphisms from $\Sigma^*$ are uniquely determined
	by their images on~$\Sigma$, it suffices to specify
	a semiautomaton by a map $\Sigma \to \mathcal T_Q$ and
	\cref{prelim:prop:UMP} then yields the extension to all of~$\Sigma^*$.
\end{remark}

\IN{simple automata}
The notion of an $\TS{\mathcal A}$-congruence was first
considered by Hartmanis~\cite{DBLP:journals/iandc/Hartmanis60} (therein called
partitions with the substitution property) in connection
with series-parallel decompositions of automata.
Slightly earlier, Rabin and Scott~\cite{DBLP:journals/ibmrd/RabinS59} had considered the closely related concept
of right-invariant congruences on~$\Sigma^*$.
The notion of a quotient automaton
is implicit in Hartmanis's work,
but was formally introduced for the first time
by Bayer~\cite{DBLP:conf/focs/Bayer66}.
However, he defined it~\cite[Def.~9]{DBLP:conf/focs/Bayer66}
only when $\TS{\mathcal A}$ is a group.
One year later, Bayer~\cite{DBLP:conf/focs/Bayer67} gave the general definition and related
the lattice of $\TS{\mathcal A}$-congruences
to endomorphisms of general automata.
Thierrin~\cite{DBLP:journals/kybernetika/Thierrin70} and Babcs\'anyi~\cite{DBLP:journals/actaC/Babcsanyi07}
investigated primitive automata $\mathcal A$ (therein called \emph{simple automata}).
For example, \cref{prelim:cor:simpleaut}
is due to Thierrin~\cite[Prop.~8]{DBLP:journals/kybernetika/Thierrin70}.
Primitivity has been linked to minimality conditions
for automata in~\cite{DBLP:journals/mst/Davies18,
	DBLP:phd/basesearch/Davies19,
	DBLP:journals/tcs/RestivoV12,
	DBLP:journals/tcs/RestivoV12a}.


In the context of coding theory and code parsing,
Schützenberger~\cite[Sect.~5 \& 6]{DBLP:journals/tit/Schutzenberger56}
introduced the syntactic semigroup of a code (which he called the
\emph{fundamental semigroup}) and represented its elements using matrices.
His construction in fact yields the transition matrices of the minimal
automaton of the code. Although automata are not mentioned explicitly,
this relation to the minimal automaton shows that his work can be viewed
as the first investigation of the transformation semigroup of an automaton.
An early work that explicitly introduced transition matrices of automata 
is the work by Burks and Wang~\cite{DBLP:journals/jacm/BurksW57a}.
Closely related stochastic transition matrices had, however, already been
used decades earlier in the theory of Markov chains.

The congruences $\sim_i$ from
Weeg~\cite[Def.~5]{DBLP:journals/jacm/Weeg62}, discussed above,
are related to the transformation semigroup as follows.
Let $\sim$ be the equivalence relation on $\Sigma^*$ defined by
\[
u \sim v
\quad\Longleftrightarrow\quad
u \sim_i v \text{ for all } i \in \{1,\ldots,n\}.
\]
Then $u \sim v$ if and only if
$\delta(q_i,u)=\delta(q_i,v)$ for all $i\in\{1,\ldots,n\}$.
Moreover, if $u\sim v$, then
$wu\sim wv$ and $uw\sim vw$ for every $w\in\Sigma^*$.
Thus, for the product on subsets we have
\[
[u]_{\sim}\,[v]_{\sim} = [uv]_{\sim},
\]
and so we obtain a well-defined operation on $\Sigma^* / {\sim}$,
making it into a semigroup.
Now the map
$
[u]_{\sim} \mapsto \fword{\delta}{u}
$
is an isomorphism onto the transformation semigroup.

\IN{homing sequence} \IN{automata with output}
\IN{sequential transducer}
In the early days, automata were considered as \emph{(sequential) transducers} (or \emph{automata with output}), see~\cite{HopUll79,Sandberg04} for details.
A notion related to synchronization in this setting is that of a \emph{homing sequence} for a transducer automaton model. Such a sequence is an input word whose produced outputs uniquely determine the state reached by the automaton. For more information, we refer to the survey by Sandberg~\cite{Sandberg04}, which discusses homing sequences and their relation to synchronizing sequences.

The seminal paper by Rabin and Scott~\cite{DBLP:journals/ibmrd/RabinS59} is widely regarded as the first to systematically consider finite automata as language acceptors.\footnote{Behind the Iron Curtain, however, automata had already been defined as language acceptors by Medvedev~\cite{Medvedev56}.} They also introduced nondeterministic automata and the power-set construction and established much of the standard terminology and notation.
Nevertheless, it was not the first paper to relate finite automata to language acceptance. Shannon~\cite{Shannon:1948:MTCa} investigated them as language-generating devices in the context of coding theory and Chomsky~\cite{DBLP:journals/tit/Chomsky56,DBLP:journals/iandc/ChomskyM58} studied them in the context of the processing, understanding, and production of natural languages by humans.\footnote{Chomsky concluded that regular languages are not sufficiently expressive to model natural languages adequately and therefore introduced more powerful grammar formalisms in~\cite{DBLP:journals/tit/Chomsky56}.} Furthermore, building on the ideas of McCulloch and Pitts~\cite{McCullochPitts43}, several authors, for example Burks and Wang~\cite{DBLP:journals/jacm/BurksW57,DBLP:journals/jacm/BurksW57a} and Kleene~\cite{Kleene56},
investigated finite automata as models of nerve nets. Kleene~\cite{Kleene56}
introduced regular expressions in this context.

The characterization of minimal automata and the Myhill--Nerode relation
go back to Nerode~\cite{nerode1958linear} (input-output version)
and Myhill~\cite{myhill1957finite} (automata/semigroup version).
Regular languages have been classified by algebraic properties of 
the transformation semigroups of minimal accepting automata
via the theory of (pseudo-)varieties of languages
and semigroups~\cite{Eilenberg1976,EILENBERG1976413,Pin86}. Seminal
results in this direction include the characterization of star-free
languages by counter-free automata~\cite{McNaughton71}
and aperiodic monoids~\cite{DBLP:journals/iandc/Schutzenberger65a},
see~\cite{Eilenberg1974,Eilenberg1976,Pin86,DBLP:reference/hfl/Pin97,Pin2022} 
for more information.


\IN{Krohn-Rhodes decomposition theory}
Krohn and Rhodes~\cite{KrohnRhodes65} developed the \emph{Krohn--Rhodes decomposition theory}, a decomposition theory for automata.
In this theory, automata are represented as homomorphic images of certain
basic permutation automata and aperiodic automata related to the given automaton.
The theory identifies basic building blocks from which every automaton can be
simulated by an appropriate product of these basic automata.
There are several formulations of this theory:

\begin{enumerate}
	\item The original semigroup-theoretic formulation from~\cite{KrohnRhodes65}
	and~\cite[Chap.~5]{Arbib68} uses the wreath product.
	Other formulations based on the same notions were given
	in~\cite{DBLP:journals/mst/Lallement71}
	and~\cite{DBLP:journals/tcs/Esik00}.
	Textbook expositions close to the original formulation can be found
	in~\cite{rhodes2009applications,RhodesSteinberg2009qTheory}.
	The monograph~\cite{rhodes2009applications} also discusses broader
	applications to biology, physics, and psychology, together with some
	philosophical reflections.
	
	\item A formulation in terms of the cascade product was developed
	in~\cite{Ginzburg68,DBLP:journals/iandc/Zeiger67a,DBLP:journals/mst/Zeiger67}
	and~\cite[Chap.~4]{Arbib68}.
	In this formulation, every automaton is
	covered\footnote{Cover maps are related to morphisms.} by a cascade
	product of certain permutation-reset automata.
	Ginzburg~\cite{Ginzburg68} gives a correction to the proof
	in~\cite{DBLP:journals/iandc/Zeiger67a,DBLP:journals/mst/Zeiger67}.
	We also refer to the expository article~\cite{Zimmermann2020KrohnRhodesSemiautomata}.
	
	\item Meyer and Thompson's~\cite{DBLP:journals/mst/MeyerT69}
	formulation in terms of cascade products connects the theory to an
	earlier decomposition result of Hartmanis and
	Stearns~\cite{HartmanisStearns1966}.

	\IN{holonomy decomposition}	
	\item Based on Zeiger's
	approach~\cite{DBLP:journals/iandc/Zeiger67a,DBLP:journals/mst/Zeiger67},
	Eilenberg~\cite{Eilenberg1976} gave an elegant formulation in the
	language of transformation semigroups, termed the
	\emph{holonomy decomposition}, see also~\cite{Holcombe82}.
	Egri-Nagy and
	Nehaniv~\cite{doemosinehaniv05,DBLP:phd/ethos/EgriNagy05,DBLP:conf/afl/Egri-NagyMN11,DBLP:conf/wia/Egri-NagyN04,DBLP:journals/actaC/Egri-NagyN06}
	gave an automata-theoretic reformulation of the holonomy
	decomposition and studied its complexity.
	
	\item Diekert, Kufleitner, and
	Steinberg~\cite{DBLP:journals/fuin/DiekertKS12} gave an algebraic
	formulation utilizing the notion of local divisors.
\end{enumerate}

\IN{Krasner–Kaloujnine universal embedding theorem}
In the special case of groups, the Krohn–Rhodes theorem reduces to the \emph{(Krasner–Kaloujnine) universal embedding theorem}~\cite{KK513}, see~\cite[Sect.~2.6]{dixon1996permutation} for a textbook exposition of this special case.

Maler and Pnueli~\cite{DBLP:conf/focs/MalerP90} carried out an in-depth study of the size of the Krohn--Rhodes decomposition and used their results to obtain an improved algorithm for translating counter-free automata into past temporal logic formulae. A corresponding bound for future linear temporal logic (LTL) was later established in~\cite{DBLP:conf/fossacs/BokerLS22}. We refer to~\cite{DBLP:conf/birthday/Maler10,MalerPnueli94} for further discussion and pointers to the literature. Results on the relationship between temporal logics and automata can be found in~\cite{DBLP:journals/jcss/CohenPP93}.

The survey by Glushkov~\cite{Glu61} also gives a good overview of the early developments in automata theory, and includes also work carried out behind the Iron Curtain.

\paragraph{More on Automata Morphisms.} 
In the transformation semigroup viewpoint, the endomorphism semigroup corresponds to the centralizer of a transformation semigroup, that is, the set of transformations commuting with every transformation in the semigroup. The centralizer of a single transformation has been characterized in~\cite{AraujoKonieczny2013_CentralizersFullTransformation}; we also refer to this work for further pointers to the literature. For permutation groups, the centralizer of a single permutation in the full symmetric group is described in~\cite[pp.~295]{suzuki1982group}, \cite[Thm.~15.2]{Wells1976_ApplicationsWreath}, and~\cite{KezlanRhee99}, see also the generalization in~\cite{Lipscomb1988_CentralizerPermutation}. 
It is a folklore result in permutation group theory that the centralizer in the symmetric group of a transitive (resp.~semiregular) permutation group is semiregular (resp.~transitive). Consequently, a regular permutation group on a finite set has the same order as its centralizer and coincides with its centralizer when it is abelian. That the centralizer of a transitive permutation group is semiregular also follows from \cref{prelim:automata:morphism:prop:endo-SCC-is-iso-b}.

\Cref{prelim:automata:morphism:prop:image-SCC,prelim:automata:morphism:cor:image-SCC}
state that morphisms between strongly trap-connected automata are either full isomorphisms or map all states to a single trap state.
This is somehow reminiscent of the well-known Schur's lemma from ring theory~\cite{lang02},
see also a related remark from~\cite[p.~275]{Rystsov97}.

In our definition of morphisms, the alphabet was fixed. We can generalize this
by including a mapping $\alpha \colon \Sigma \to \Gamma$ of alphabets
in the definition and let $\varphi \colon Q \to S$ be a morphism between $\mathcal A = (Q, \Sigma, \delta)$
and $\mathcal B = (S, \Gamma, \mu)$
if 
\[
\varphi(\delta(q, x)) = \mu(\varphi(q), \alpha(x))
\]
for all $q \in Q$ and $x \in \Sigma$ (a definition
in this direction was given in~\cite{Wells1976_ApplicationsWreath}).

Given such a morphism that is surjective, let $F \colon \TM{\mathcal A} \to \TM{\mathcal B}$
be given by $F(\delta_{u}) = \mu_{\alpha(u)}$
for $u \in \Sigma^*$ with $\alpha$ extended to a mapping (also denoted by~$\alpha$) 
from $\Sigma^*$ into $\Gamma^*$ via \cref{prelim:prop:UMP}.
This is well-defined since for $u,v \in \Sigma^*$
with $\delta_u = \delta_v$,
by choosing for $s \in S$ some $q \in Q$ with $s = \varphi(q)$
we have 
\[
\mu_{\alpha(u)}(s) = \mu(\varphi(q), \alpha(u)) = \varphi(\delta(q, u)) = \varphi(\delta(q, v)) = \mu(\varphi(q), \alpha(v)) = \mu_{\alpha(v)}(s)
\]
Then we have the following:
\begin{enumerate}
\item $F$ is a monoid morphism since
\[
F(\concat{\delta_u}{\delta_v}) = F(\delta_{uv}) = \mu_{\alpha{uv}} = \concat{\mu_{\alpha(u)}}{\mu_{\alpha(v)}} = \concat{F(\delta_u)}{F(\delta_v)}
\]

\item For $t \in \TM{\mathcal A}$
with $t = \delta_u$ for $u \in \Sigma^*$ we have
\[
F(t)(\varphi(q))
= F(\delta_u)(\varphi(q))
= \mu(\varphi(q), u) = \varphi(\delta(q, u)) = \varphi(\delta_u(q))
= \varphi(t(q))
\]
for $q \in Q$.
\end{enumerate}
This is related to the notion of permutation isomorphism
from~\cite[p.~17]{dixon1996permutation}
which is defined between two permutation groups $G \le_G \mathcal S_X$
and $H \le_G \mathcal S_Y$ as a bijection $\lambda \colon X \to Y$
with a group isomorphism $\psi \colon G \to H$
such that $\lambda(g(p)) = \psi(g)(\lambda(p))$
for all $p \in X$ and $g \in G$.



\label{def:R-quasiorder}
\paragraph{Reachability Quasiorder.}
The reachability relation appears to have been introduced independently
by Meyer and Thompson~\cite{DBLP:journals/mst/MeyerT69},
Brzozowski and Fich~\cite{DBLP:journals/jcss/BrzozowskiF80},
and Letichevsky~\cite{Letichevsky65}; it was later used
in~\cite{DBLP:journals/tcs/Ito12}.
Brzozowski and Fich~\cite{DBLP:journals/jcss/BrzozowskiF80}
proved that a language is accepted by a partially ordered automaton
if and only if Green's $\mathcal R$-relation~\cite{Green51}
on its transition semigroup is the equality relation
(this was also proved independently in~\cite[Thm.~5.1]{DBLP:journals/jcss/CohenPP93}).
We obtain the $\mathcal R$-relation as the equivalence relation induced by
the quasiorder~$\le_{\mathcal R}$ on $\TS{\mathcal A}$ given by
\[
\fword{\delta}{u} \le_{\mathcal R} \fword{\delta}{v}
\Leftrightarrow
\exists w \in \Sigma^* \ \forall q \in Q :
\applytrans{\delta}{q}{u} = \applytrans{\delta}{q}{vw}
\]
for $\semiautomat{A}{Q}{\Sigma}{\delta}$ and $u,v \in \Sigma^*$.
Clearly, if $v$ is a prefix of $u$, then
$\fword{\delta}{u} \le_{\mathcal R} \fword{\delta}{v}$.
Furthermore,
$\fword{\delta}{u} \le_{\mathcal R} \fword{\delta}{v}$
implies
$\applytrans{\delta}{q}{v} \lereach \applytrans{\delta}{q}{u}$
for every $q \in Q$. Consequently,
\begin{equation}\label{Rrelation}
	\fword{\delta}{u} \mathcal R \fword{\delta}{v}
	\Rightarrow
	\forall q \in Q \colon
	\applytrans{\delta}{q}{u}
	\eqreach
	\applytrans{\delta}{q}{v}.
\end{equation}
It follows from~\cref{Rrelation} that, for a partially ordered automaton,
the $\mathcal R$-relation is the equality relation.
Conversely, let
$\mathcal A = (Q, \Sigma, \delta, q_0, F)$
be such that the $\mathcal R$-relation on $\TS{\mathcal A}$
is the equality relation.
Then the automaton
$\mathcal B = (\TS{\mathcal A}, \Sigma, \mu,
\fword{\delta}{\varepsilon}, E)$,
where
$\mu(\fword{\delta}{u}, a) = \fword{\delta}{ua}$
and
$E = \{\fword{\delta}{u} \mid \fword{\delta}{u}(q_0) \in F\}$,
is partially ordered and satisfies
$L(\mathcal A) = L(\mathcal B)$.\footnote{Alternatively, let
	$p,q \in Q$ satisfy
	$\delta(p,u)=q$ and $\delta(q,v)=p$
	for some $u,v\in\Sigma^*$.
	Set $w=uv$.
	Since $\mathcal A$ is finite,
	$\fword{\delta}{w^i}=\fword{\delta}{w^{i+p}}$
	for some $i\ge0$ and $p>0$.
	This implies
	$\fword{\delta}{w^iu}\mathcal R\fword{\delta}{w^i}$.
	By assumption,
	$\fword{\delta}{w^iu}=\fword{\delta}{w^i}$.
	Therefore,
	$q=\delta(p,(uv)^iu)=\delta(p,(uv)^i)=p$,
	showing that $\mathcal A$ is partially ordered.
	In fact, this proves that
	$\le_{\mathcal R}$ is a partial order
	if and only if $\lereach$ is a partial order.}
This gives a short proof of the characterization due to
Brzozowski and Fich~\cite{DBLP:journals/jcss/BrzozowskiF80}.

The notion of an upper set (for the reachability quasiorder)
appears under the name \emph{subautomaton} in~\cite{Rystsov97}.
It also appears in other contexts, for example in the notion of
coverage of two automata introduced in~\cite{DBLP:journals/dm/BrzozowskiS73}.
As mentioned in the main text, the Rees congruence goes back to
a related congruence introduced by Rees~\cite{Rees40} for semigroups.

The notion of a sink component (and a corresponding dual notion) can actually be defined
for any quasiordered set, and the results from \cref{prelim:sec:reach}
transfer to this more general setting. Furthermore, reachability is a property of the underlying automaton digraph, so the reachability order can also be defined in the context of digraphs.
Similarly,
the notions of strong and unilateral connectivity in digraphs
correspond to these notions in automata,
and a strongly connected component in an automaton is precisely
a strongly connected component of its underlying automaton digraph.
However, as we need these concepts only in the context of automata,
we introduce them here for automata and rather than for general quasiorders and digraphs. 